% fgib_theory.tex —— FGIB 完整理论推导（四个问题的证明 + FGIB 机制的证明）
% 本文件合并自 fgib_ceb_collapse.tex / fgib_degeneration.tex / fgib_bound.tex / fgib_linear.tex，
% 原四文件可删除；全部标签名保持不变，fgib_intro.tex / fgib_method.tex 的 \ref 不受影响。
% 用法：\usepackage{amsmath,amssymb,amsthm}，正文处 % fgib_theory.tex —— FGIB 完整理论推导（四个问题的证明 + FGIB 机制的证明）
% 本文件合并自 fgib_ceb_collapse.tex / fgib_degeneration.tex / fgib_bound.tex / fgib_linear.tex，
% 原四文件可删除；全部标签名保持不变，fgib_intro.tex / fgib_method.tex 的 \ref 不受影响。
% 用法：\usepackage{amsmath,amssymb,amsthm}，正文处 % fgib_theory.tex —— FGIB 完整理论推导（四个问题的证明 + FGIB 机制的证明）
% 本文件合并自 fgib_ceb_collapse.tex / fgib_degeneration.tex / fgib_bound.tex / fgib_linear.tex，
% 原四文件可删除；全部标签名保持不变，fgib_intro.tex / fgib_method.tex 的 \ref 不受影响。
% 用法：\usepackage{amsmath,amssymb,amsthm}，正文处 % fgib_theory.tex —— FGIB 完整理论推导（四个问题的证明 + FGIB 机制的证明）
% 本文件合并自 fgib_ceb_collapse.tex / fgib_degeneration.tex / fgib_bound.tex / fgib_linear.tex，
% 原四文件可删除；全部标签名保持不变，fgib_intro.tex / fgib_method.tex 的 \ref 不受影响。
% 用法：\usepackage{amsmath,amssymb,amsthm}，正文处 \input{fgib_theory.tex}。
% 记号沿用 Fischer (2020)、Kolchinsky & Tracey (2019, Caveats)、Kolchinsky et al. (2017, NIB)、
% Alemi et al. (2017, VIB)。
%
% 核心结论（FGIB 论文问题陈述与解法）：
% 第一部分（四条失效通道）：
%   P1 退化继承（死旋钮）：CEB ≡ IB(β=1+γ>1)，确定性场景下所有 γ>0 的唯一最优解钉在角点，
%      γ 扫不出任何中间点（SIB Caveat 1 的退化半区，经 β'=1/(1+γ)∈(0,1) 映射）。
%      该角点恰是 CEB 自己的 MNI 目标，rem:mni 正面处理这一反驳；
%   P2 证书盲视 + 1:β 交换率：T_α 流形上 ReX 恒为 0（证书对标签保留量完全盲视），
%      且目标以 1 nat 标签信息 ↔ β nat 认证残差的比率交换；
%   P3 方差竞赛：KL 两侧皆可训练，分类项压 σ→0、钳位处退化为权重 e^{10} 的均值匹配。
%      (ii)(iii) 是受力方向分析而非收敛证明，且可实测（rem:p3-status）；
%   P4 对象错位：两种方法的界都只界住各自的含噪码，都不界住确定性连续推理表示的残差
%      （后者为 +∞）。差别是结构性的（噪声在不在推理路径上），不是紧度排序。
% 第二部分（FGIB 机制，以及它没能关闭的通道）：
%   M1 等价性：旁路 = 固定 A 下作用在 h+已知噪声上的 CEB 瓶颈（逐步成立，非全程固定参考）；
%   M2 扫参：固定分类器尺度 c 时 L_β(s)=CE(s)+(β/2)(s−a)² 严格凸、每 β 唯一内点、曲线严格凹；
%      c 自由时下确界对所有 β 均为 0（rem:c-scale，给出权重归一化分类器的架构修法）；
%   M3 证书：gap 恒等式 + 矩匹配 + 固定编码器下的值排序；固定参考的力持续性成立，
%      但侧头 A 是推理无关的第二个自由方（rem:A-channel），构成未关闭的通道。
%
% 修订记录（第二轮）：理论/实现分层后统一引入锚点方差 τ²（先验 N(a v_y, τ²I)，默认 τ²=1，
% 后验方差最优 σ²*=τ²，σ_p 为公共方差记号；point-anchor 极限统一写为 τ→0）；理论旁路固定为
% 可逆线性头 T(h)=Ah（实现用带 bias 的 nn.Linear 近似，见理论-代码分层 caveat）；
% 明确 sweep 证明的是对齐族上 KL–CE surrogate 曲线（rem:surrogate），不是 IB 互信息曲线；
% 无穷互信息结论统一挂靠分布性 standing assumption；rem:c-scale 的架构修法收紧为
% 固定温度 cosine 分类器（weight decay 只缓解不发散）；a=0 退化点、条件数表述、
% 梯度消失条件、row-space 表述等小项修正。
%
% 修订记录（第一轮）：删除三处错误论断——(1) P3 原称 ReX≈0 时 I(X;Z|Y) 发散，与 gap 恒等式矛盾；
% (2) prop:object(c) 原称 FGIB 的界"更紧/唯一有效"，由 DPI 不成立；(3) prop:conditioning(iii)
% 原从 Gram 恒等式推出条件数为 1，有反例。修正 P2 的换算、Caveat 2 的主张、rem:frame 的
% A^T/V^T、mu_head 的仿射与非推理路径身份。未证结论统一收入 sec:limits。

\newtheorem{lemma}{Lemma}
\newtheorem{proposition}{Proposition}
\newtheorem{remark}{Remark}

% =====================================================================
\section*{Part I. The four failure channels of the conditional bottleneck}
% =====================================================================

\textbf{Notation.}
$X$ is the input and $Y$ the label with $K$ classes; the encoder induces the Markov chain
$Y - X - Z$. The \emph{deterministic scenario} is $Y = f(X)$ ($H(Y|X) = 0$). Following
Fischer (2020), we write the MNI point $I(X;Y) = I(X;Z) = I(Y;Z)$ (Eq.~1), the
chain-rule identity $I(X;Z|Y) = I(X;Z) - I(Y;Z)$ (Eq.~4), the CEB objective
$\mathrm{CEB} \equiv \min_Z I(X;Z|Y) - \gamma I(Y;Z)$ (Eq.~5), the IB objective
$\mathrm{IB} \equiv \min_Z I(X;Z) - \beta_{\mathrm{IB}} I(Y;Z)$ (Eq.~16), the
variational objective
$\mathrm{VCEB} \equiv \min_{e,b,c} \langle \log e(z|x)\rangle - \langle \log b(z|y)\rangle
- \gamma \langle \log c(y|z)\rangle$ (Eq.~15), and the residual certificate
$\mathrm{ReX} \equiv \langle \log e(z|x)\rangle - \langle \log b(z|y)\rangle
\ge I(X;Z|Y)$ (Eq.~20), where $b(z|y)$ is a \emph{trainable} backward encoder (we write
$q(z|x)$ for Fischer's encoder $e(z|x)$; the two notations are interchangeable
throughout). In the loss convention $\mathrm{CE} + \beta\,\mathrm{KL}$ used throughout, the weights are
$\beta = 1/\gamma$ and $\beta_{\mathrm{IB}} = 1 + \gamma = 1 + 1/\beta$. Following
Kolchinsky \& Tracey (2019), we write the IB curve $F(r)$ (their Eq.~1), the IB
Lagrangian $L^{\mathrm{IB}}_\beta = I(Y;T) - \beta I(X;T)$ (their Eq.~2), the
DPI-saturating manifold $T_\alpha := B_\alpha \cdot Y$ (their Eq.~6), the bound
$L^{\mathrm{IB}}_\beta \le (1-\beta)H(Y)$ (their Eq.~7), and their footnote~4 (the
$\beta > 1$ regime of trivial maximizers). The variational objects are
\begin{align*}
R &= \mathbb{E}\big[\operatorname{KL}(q(z|x)\,\|\,m(z))\big] \ge I(X;Z)
&&\text{(VIB rate, marginal prior $m$; Fischer Eq.~19, Alemi et al.\ Eq.~14)} \\
\mathrm{ReX} &= \mathbb{E}\big[\operatorname{KL}(q(z|x)\,\|\,b_\theta(z|y))\big]
\ge I(X;Z|Y)
&&\text{(CEB residual; Fischer Eq.~20)} \\
\widehat I &= -\frac{1}{N}\sum_i \ln\Big[\frac{1}{N}\sum_j
e^{-\|f_i-f_j\|^2/(2\sigma^2)}\Big] \ge I(X;M)
&&\text{(NIB, non-parametric; Eq.~16, the homoscedastic case of Eq.~10)} \\
\mathrm{ReX}_F &= \mathbb{E}\big[\operatorname{KL}(q(z|x)\,\|\,r(z|y))\big]
\ge I(X;Z|Y)
&&\text{(FGIB, fixed anchor)}
\end{align*}
(Fischer 2020, Eqs.~18--20; Alemi et al.\ 2017, Eq.~14; Kolchinsky et al.\ 2017,
Eqs.~10 and 16). Fischer's comparison of the first two: at the respective prior optima
($m = q(z)$, $b = q(z|y)$), $\mathrm{ReX} = R - I(Y;Z)$ --- the conditional prior removes
exactly the retained label information, so CEB's bound is tighter than VIB's by
$I(Y;Z)$. We ask the same question one level down: where does FGIB's fixed-anchor bound
sit relative to CEB's learnable-backward bound?

\begin{proposition}[Channel P1, dead knob: every $\gamma > 0$ selects the same corner]\label{prop:dead}
Assume $Y = f(X)$. For every $\gamma > 0$,
\begin{equation}
\min_Z L_\gamma(Z) = -\gamma H(Y),
\qquad
L_\gamma(Z) := I(X;Z|Y) - \gamma I(Y;Z),
\label{eq:min}
\end{equation}
attained exactly on the corner family
$\{Z : I(Y;Z) = H(Y),\; I(X;Z|Y) = 0\}$ (e.g.\ $Z = Y$, or $Z = (Y, N)$ with
$N \perp X \mid Y$): the corner is unique in the information plane, while at the
representation level it is an entire family. Consequently the map
$\gamma \mapsto \big(I(X;Z^*|Y),\, I(Y;Z^*)\big)$ is constant: no interior point of the
IB curve is ever a Lagrangian minimizer, and CEB inherits the deterministic degeneration
of the Caveats paper (their Issue~1) over its \emph{entire} parameter range, since
$\beta_{\mathrm{IB}} = 1 + \gamma > 1$ always.
\end{proposition}

\begin{proof}
By Fischer's Eq.~4, $L_\gamma(Z) = I(X;Z) - (1+\gamma)I(Y;Z)$. The DPI
($I(Y;Z) \le I(X;Z)$, by the Markov chain $Y - X - Z$) gives
\begin{equation}
L_\gamma(Z) \ge I(Y;Z) - (1+\gamma)\,I(Y;Z)
= -\gamma\,I(Y;Z) \ge -\gamma H(Y),
\label{eq:chain}
\end{equation}
with equality iff both inequalities bind: $I(X;Z) = I(Y;Z)$ and $I(Y;Z) = H(Y)$.
$Z = Y$ attains the bound. Fischer (2020, p.~7) states the underlying equivalence
explicitly: CEB and IB are the same objective for $\gamma = \beta_{\mathrm{IB}} - 1$.
In the Caveats paper's maximizing convention, minimizing $L_\gamma$ is maximizing
$L^{\mathrm{IB}}$ with $\beta' = 1/\beta_{\mathrm{IB}} = 1/(1+\gamma) \in (0,1)$ ---
precisely the range in which their Eq.~7 pins every maximizer at the copy corner (their
Issue~1); their footnote~4 covers the complementary regime $\beta' > 1$ (i.e.\
$\gamma < 0$, outside CEB's range), whose maximizers are the trivial collapse. \hfill\qed
\end{proof}

\begin{remark}[The corner is the MNI point: why that does not rescue the knob]\label{rem:mni}
In the deterministic scenario the corner of Proposition~\ref{prop:dead} \emph{is} the MNI
point: $I(Y;Z) = H(Y) = I(X;Y)$ together with $I(X;Z|Y) = 0$ gives
$I(X;Y) = I(X;Z) = I(Y;Z)$, which is Fischer's Eq.~1. CEB's declared target and the
degenerate attractor therefore coincide, and Kolchinsky \& Tracey concede the point
themselves (their Section~4): ``if $Y = f(X)$ and one's goal is precisely to find the corner
point $\langle H(Y), H(Y)\rangle$ (maximum compression at no prediction loss), then our
results suggest that the IB Lagrangian provides a robust way to do so, being invariant to
the particular choice of $\beta$.'' Read that way, Proposition~\ref{prop:dead} states
hyperparameter robustness rather than failure. Three things keep it from closing the matter.
\begin{enumerate}
\item[(i)] \emph{The corner is one point in the information plane and an entire equivalence
class of representations.} The objective is exactly constant on
$\{Z : I(Y;Z) = H(Y),\ I(X;Z|Y) = 0\}$, which contains $Z = Y$, $Z = (Y,N)$ with
$N \perp X \mid Y$, and every bijective recoding of these. Adversarial robustness, OoD
detection and calibration --- the properties CEB is argued for --- are properties of the
\emph{geometry} of $Z$, on which the objective has no opinion; which member of the class is
reached is settled by the optimizer. That is channels P3 and P4 below.
\item[(ii)] \emph{CEB's own reading of the knob is unavailable at the exact optimum.}
Fischer (2020, p.~7) writes ``$\rho < 0$ may capture less than the MNI. $\rho > 0$ may
capture more'', whereas Proposition~\ref{prop:dead} places every $\gamma > 0$ at the MNI
point exactly. His argument for preferring $\gamma = 1$ is correspondingly not about
minimizers but about the optimizer: weighting the two appearances of $I(Y;Z)$ differently
``forces the optimization procedure to count bits of $I(Y;Z)$ in two different ways'' (his
Eq.~8 and surrounding discussion).
\item[(iii)] \emph{The knob does move trained models, which locates what it acts on.} On
Fashion MNIST --- a deterministically labelled dataset --- Fischer's own Figure~4 reports
the rate lower bound $R_X$ rising from $\approx 2.0$ to $\approx 5.0$ nats as $\rho$ goes
from $-1$ to $5$, with $\rho = 0$ landing on $I(X;Y) = \log 10 \approx 2.3$ nats. Since
every $\gamma > 0$ shares one exact optimum, that axis is not a path along the IB curve: it
is the distance between the optimum of the variational objective over a restricted family
and the point at which training was stopped. A knob acting through the approximation error
is precisely the knob that Propositions~\ref{prop:flat} and~\ref{prop:variance} show cannot
be read off the certificate.
\end{enumerate}
Kolchinsky \& Tracey state that their results ``are based on analytically-provable
properties of the IB curve, i.e., global optima'' and ``do not concern practical issues of
optimization''; Proposition~\ref{prop:flat} supplies the finite-$\beta$ variational
counterpart. The account is owed by any method reporting a $\beta$-sweep on
deterministically labelled data --- FGIB and the experiments of this paper included
(Remark~\ref{rem:c-scale}).
\end{remark}

\begin{proposition}[Channel P2, a blind certificate and a $1{:}\beta$ exchange rate]\label{prop:flat}
On the manifold $T_\alpha = B_\alpha \cdot Y$ of the Caveats paper (their Eq.~6),
$I(X;T_\alpha|Y) = 0$ for every $\alpha$, and with the optimal backward encoder and
classifier,
\begin{equation}
\mathrm{ReX}(T_\alpha) = 0,
\qquad
\mathrm{CE}(T_\alpha) = H(Y) - I(Y;T_\alpha) =: \delta_\alpha,
\qquad
\mathrm{VCEB}(T_\alpha) = \gamma\,\delta_\alpha .
\label{eq:vceb}
\end{equation}
Two consequences.
\emph{(a) The certificate is blind to label retention.} $\mathrm{ReX} = 0$ holds from the
total collapse ($\alpha = 0$, no label information) to the full copy ($\alpha = 1$), so its
tightness is compatible with retaining all, part, or none of $Y$. Fischer's reading of the
certificate as a meter --- ``observing ReX tells us exactly how many bits we are from the
MNI point'' (p.~7) --- therefore holds only along the $I(X;Z|Y)$ axis; on the necessity axis
the meter reads $0$ at chance level.
\emph{(b) Exchange rate.} With the optimal classifier the objective is
$L = \mathrm{CE} + \beta\,\mathrm{ReX} = \delta + \beta\,\mathrm{ReX}$, so it is indifferent
between giving up $\delta$ nats of label information and saving $\delta/\beta$ nats of
certified residual. In the heavy-regularization regime the bottleneck is supposed to act in,
the whole label axis --- up to $H(Y)$ nats --- is purchasable for $H(Y)/\beta$ nats of
certified compression.
\end{proposition}

\begin{proof}
$T_\alpha$ is a function of $Y$ alone, so $T_\alpha \perp X \mid Y$ and
$I(X;T_\alpha|Y) = 0$; the Caveats paper shows $I(Y;T_\alpha)$ sweeps $[0, H(Y)]$.
With the true backward encoder ($b = p(z|y)$) the certificate is tight:
$\mathrm{ReX} = I(X;T_\alpha|Y) = 0$ (Fischer Eqs.~9--11, 20). The optimal classifier
achieves $\langle \log c(y|z)\rangle = -H(Y|T_\alpha)$ (Fischer Eq.~12), i.e.\
$\mathrm{CE} = H(Y|T_\alpha) = \delta_\alpha$, and the VCEB value follows from Eq.~15.
Part~(b) is the definition of the objective evaluated at $\mathrm{CE} = \delta$.
\hfill\qed
\end{proof}

\begin{remark}[What the flatness claim is, and what it is not]\label{rem:convention}
An earlier form of this argument read the near-optimality band off the \emph{value}
$\mathrm{VCEB} = \gamma\delta \to 0$. That reading is a scaling artifact:
$\mathrm{VCEB} = \gamma\,(\mathrm{CE} + \beta\,\mathrm{ReX}) = \gamma L$, so the two
conventions differ by a positive factor and share their minimizers, and on $T_\alpha$ itself
the objective is $L = \delta$ in either convention and strictly prefers $\alpha = 1$. What
is convention-free is the \emph{ratio} of Proposition~\ref{prop:flat}(b): the prediction
term is weaker than the compression term by exactly the factor $\beta$, so in units of the
compression term the entire manifold lies within $H(Y)/\beta$ of optimal. Under a
scale-invariant optimizer (Adam, used throughout the experiments) only the ratio matters;
under plain SGD at fixed learning rate the convention additionally rescales the effective
step size. Combined with Proposition~\ref{prop:flat}(a) --- no certified compression is
bought anywhere on the manifold --- the operative statement is that the label axis is cheap
and unmetered, not that the objective is literally flat along it.
Pushing $\beta$ up also drives $\beta_{\mathrm{IB}} = 1 + 1/\beta \to 1^{+}$: the objective
approaches the Caveats paper's $\beta = 1$ degeneracy (their Eq.~7 discussion), at which
\emph{every} point of the diagonal is simultaneously optimal, from above. Fischer (2020,
p.~7) anticipates the direction --- ``$\rho < 0$ may capture less than the MNI'' --- but by
Proposition~\ref{prop:dead} that cannot happen at the exact optimum;
Proposition~\ref{prop:flat} gives the rate at which the variational objective is willing to
let it happen.
\end{remark}

\begin{remark}[Approximately deterministic labels]\label{rem:approx}
The Caveats paper (Appendix~C, Issue~1) shows that when $Y$ is only $\epsilon$-close to
deterministic, all optimizers fall within $O(-\epsilon\log\epsilon)$ of a single corner:
the dead-knob statement of Proposition~\ref{prop:dead} degrades gracefully, and the
near-optimality band of Proposition~\ref{prop:flat} is perturbed by the same order.
\end{remark}

\begin{proposition}[Channel P3, collusion collapse: the variance race and the overfit channel]\label{prop:variance}
For the diagonal Gaussian family $e(z|x) = \mathcal{N}(m(x), \sigma^2 I)$,
$b(z|y) = \mathcal{N}(g(y), \tau^2 I)$ (the implementation's setup, with the
\texttt{logvar} clamp $[-10, 10]$ of \texttt{kl\_divergence}), and a \emph{linear}
classifier $c(z) = Wz + b$ (the implementation's classifier is \texttt{nn.Linear}):
\begin{enumerate}
\item[(i)] for fixed means, $\operatorname{KL}$ is minimized over $\sigma^2$ at
$\sigma^2 = \tau^2$, with value $\tfrac{1}{2\tau^2}\|m - g\|^2$ --- a pure mean-matching
term;
\item[(ii)] the classifier term favors $\sigma^2 \to 0$:
$\mathrm{CE}(c(z), y) = -\log \mathrm{softmax}_y(Wz+b)$ is convex in $z$ (log-sum-exp
of an affine function of $z$, minus a linear term --- the convexity claim requires $c$
affine), so Jensen gives
$\mathbb{E}_z[\mathrm{CE}(c(z), y)] \ge \mathrm{CE}(c(m), y)$. The KL's derivative in
$\sigma^2$ vanishes at $\sigma^2 = \tau^2$, but along the matched line it \emph{resists}
shrinking $\tau^2$ while the means are unmatched
($\partial \operatorname{KL}/\partial \tau^2|_{\sigma^2=\tau^2}
= -\|m-g\|^2/(2\tau^4) < 0$). The joint dynamics therefore runs in two stages: the mean
matching $m \to g$ completes first (with weight $1/\tau^2$, which grows as $\tau$
shrinks), and once $\|m-g\| \to 0$ the variance race is unopposed --- the classifier
drives $\sigma^2 \to 0$ with $\tau^2$ following the moment-matched value --- until the
clamp at $s = e^{-10}$, where the encoder collapses toward the deterministic code
$z \approx g(y)$ and
$\mathrm{ReX} = \tfrac{1}{2s}\mathbb{E}\|m(x) - g(y)\|^2$: a mean-matching force of
weight $1/s = e^{10} \approx 2.2\times10^{4}$ between two \emph{trainable} maps;
\item[(iii)] the certified object and the deployed object separate. By the gap lemma
(Lemma~\ref{lem:gap}) $\mathrm{ReX} = I(X;Z|Y) + \Delta \ge I(X;Z|Y)$, so a small
$\mathrm{ReX}$ \emph{entails} a small residual for the stochastic training code --- the
certificate cannot under-report its own object. What it does not control is the relation
between the training code and the deployed representation. Exact mean matching
($m(x) = g(y)$) makes both objects class-constant, so both residuals vanish; the
scenario of interest is the \emph{near}-match regime, where
$\mathbb{E}\,\|m(X) - g(Y)\|^2 = o(s)$ keeps $\mathrm{ReX}$ small while the deployed
continuous $m$ retains within-class variation, and any non-class-constant continuous
$m$ has $I(X;m(X)|Y) = +\infty$ under the standing non-atomic assumption (the
pathology of deterministic networks, Amjad \&
Geiger 2018; Kolchinsky et al.\ 2017). The variance race therefore moves the
\emph{certified} object toward a noise-dominated class-constant code while the
\emph{deployed} object can keep diverging residual --- and the encoder is fitted to the
backward encoder's per-class targets, memorizing the label code $g(y)$ rather than
being told what residual to erase, while the certificate reports $\mathrm{ReX} \approx
0$ on a code that inference does not use.
\end{enumerate}
\end{proposition}

\begin{proof}
The closed form
$\operatorname{KL} = \tfrac{d}{2}\big[\sigma^2/\tau^2 + \|m-g\|^2/(d\tau^2)
- 1 + \ln(\tau^2/\sigma^2)\big]$ has derivative $\tfrac{d}{2}(1/\tau^2 - 1/\sigma^2)$
in $\sigma^2$, so the variance optimum is $\sigma^2 = \tau^2$, giving (i); at
$\sigma^2 = \tau^2 = s$ the value is $\|m-g\|^2/(2s)$. For (ii): $\mathrm{CE}(c(z), y)$ has Hessian
$\operatorname{diag}(p) - pp^{\top} \succeq 0$ (the categorical Fisher information), so
it is convex in $z$ and Jensen gives the stated inequality; the same Hessian identity makes
$\sigma^2 \mapsto \mathbb{E}_z[\mathrm{CE}(c(z),y)]$ non-decreasing (Gaussians of larger
variance dominate in convex order), strictly increasing in $\sigma^2$ wherever the
classifier is unsaturated, so that term pushes $\sigma^2$ down. The KL's partial derivatives at the
matched point are $\partial_{\sigma^2}\operatorname{KL} = 0$ and
$\partial_{\tau^2}\operatorname{KL} = -\|m-g\|^2/(2\tau^4) < 0$, so the KL pulls
$\tau^2$ up until the means match, while the mean-matching gradient on the encoder is
$(m - g)/\tau^2$; once $\|m-g\| \to 0$ the race to the clamp is unopposed. For (iii): the
gap identity (Lemma~\ref{lem:gap}) forces the certified residual and the certificate to
vanish together; the divergence attaches to the deployed mean $m(x)$ whenever it is not
class-constant (Amjad \& Geiger; Kolchinsky et al.). \hfill\qed
\end{proof}

\begin{remark}[Status of the two-stage account]\label{rem:p3-status}
Parts~(ii) and~(iii) are a stationarity-and-direction-of-force argument, not a convergence
proof: the ``two stages'' presume a time-scale separation (mean matching completes before
the variance race) that is not established here. The account is directly falsifiable in
this codebase: log the mean posterior log-variance and the fraction of units at the clamp
$s = e^{-10}$ during CEB training. \emph{We ran the test.} On MNIST, across all six
$\beta$ settings and 5 seeds (Section~\ref{sec:ablation} of the main paper), the trained
posterior and prior log-variances of CEB remain near their initialization ($\sigma^2 =
\tau^2 = 1$, mean $\log\sigma^2$ within $0.74$ of zero, clamp fraction $0.0$ at every
epoch of every run). At the early-stopping time scale of the experimental protocol, the
variance race therefore does not materialize, and
Proposition~\ref{prop:variance}(iii) is empirically void for those runs: Part~(iii)
stands as a sharp hypothesis about what longer training or different schedules would
produce, the force analysis of (i)--(ii) as the only proved content, and the
measurement itself as the finding.
\end{remark}

\begin{remark}[Train--evaluation mismatch]\label{rem:mismatch}
The classifier is trained on the sampled code $z$ and evaluated on the mean $m(x)$
(Fischer 2020, Appendix~A.1). Proposition~\ref{prop:variance}(ii) is the force that
closes this gap by destroying the noise altogether ($\sigma \to 0$), at which point the
bottleneck's stochasticity --- the object the bound is defined on --- disappears.
\end{remark}

\begin{proposition}[Channel P4, certificate-object mismatch: neither bound is an upper bound on the deployed representation]\label{prop:object}
(a)~\emph{CEB}. The classifier consumes the sampled code $Z$ at training and the mean
$\mu(X)$ at evaluation (Fischer 2020, Appendix~A.1). At the variance-optimal posterior,
$\sigma(x) \equiv \sqrt{S_\theta(y)}$ (the backward encoder's per-class standard
deviation) is class-constant, so $Z = \mu(X) + \sqrt{S_\theta(Y)}\,\varepsilon$ is a
noisy function of $\mu(X)$ given $Y$; by the DPI,
\begin{equation}
I(X; Z | Y) \;\le\; I(X;\, \mu(X) | Y) ,
\label{eq:dpi}
\end{equation}
strictly so whenever the noise destroys within-class detail. Hence CEB's certificate
$\mathrm{ReX} \ge I(X;Z|Y)$ is a bound on the \emph{training code}, not on the residual
of the representation deployed at inference. The gap is not a small correction: whenever
$\mu$ is a continuous map that is not class-constant, $I(X;\mu(X)|Y) = +\infty$ under
the standing non-atomic assumption (Amjad \& Geiger 2020; Kolchinsky et al.\ 2017), so
no finite quantity --- CEB's certificate included --- bounds the deployed residual from
above. The CE term presses the noise
toward $0$ (by Jensen, $\mathbb{E}_z[\mathrm{CE}(c(z),y)] \ge \mathrm{CE}(c(\mu),y)$),
but no finite $\sigma$ closes the gap.
(b)~\emph{FGIB}. The inference path is $h$ at both training and evaluation; the side
noise of Proposition~\ref{prop:inverse} is not part of any inference path. By the DPI
applied to the same perturbed-variable construction,
\begin{equation}
I(X;\, h + \tau (A^{\top}A)^{-1/2}\varepsilon \,|\, Y)
\;\le\; I(X; H | Y) ,
\label{eq:dpi-f}
\end{equation}
so the direction of the bound is the same as CEB's: $\mathrm{ReX}_F$ upper-bounds the
noisy side-path object, and \emph{no} finite value upper-bounds the deployed residual.
The point-anchor limit $\tau \to 0$ does not repair this: it is a change of model whose
bound value $\|\mu - a v_y\|^2/(2\tau^2)$ typically diverges
(Remark~\ref{rem:limit}), and even the limiting object needs $A$ well-conditioned
(Proposition~\ref{prop:conditioning}(iii)).
(c)~\emph{Conclusion}. In the pure variational value sense (same encoder, same Gaussian
family), CEB's bound is tighter (Proposition~\ref{prop:ordering}). In the object sense
the two are not ordered: both are valid upper bounds on their respective noisy codes, and
neither bounds the deterministic deployed representation. What remains structural --- not
a tightness claim --- is \emph{where the gap lives}: CEB's perturbation is inside the
inference path and is driven by the prediction term (so removing it changes the model),
while FGIB's perturbation lives on a side path that inference never reads and is
controlled by side-path parameters alone. The latter is a cleaner separation of
regularizer from model, not a tighter bound.
\end{proposition}

\begin{remark}[The bound is a training instrument, not an inference one]\label{rem:training-role}
No method computes its bound at inference: the bound's roles are \emph{training-time}
--- a regularizer (a differentiable surrogate for the intractable information term) and
a meter (Fischer 2020, p.~7: ``observing ReX tells us exactly how many bits we are from
the MNI point''). In the $Z \perp X \mid Y$ limit CEB's mean is also class-constant
($m(x) = \mathbb{E}[Z|x] = \mathbb{E}[Z|y]$), so the mismatch of
Proposition~\ref{prop:object}(a) is not a limit failure: it is a finite-$\beta$
\emph{under-reporting} --- the regularizer controls, and the meter reports, the training
code, which is strictly more compressed than the deployed representation $\mu(X)$, by a
positive, trained deficit that cannot be measured from the certificate. FGIB's
side-path noise is a known side-path parameter and is removable in principle
($\tau \to 0$ changes the side path, never the deployed model), but for finite
$\tau$ its regularizer and meter act on a \emph{perturbed copy} of the deployed
object, with the perturbation size governed by the conditioning of $A$
(Proposition~\ref{prop:conditioning}(iii)).
The one genuine inference-time use of a certificate value is as a detection score
(Fischer 2020, RG3, uses the marginal rate $R$): CEB's label-conditioned
$\mathrm{ReX}$ is not even computable at inference (the label is unknown --- Fischer
trains a separate marginal for this purpose), while FGIB's geometry gives a label-free
score, the distance to the nearest anchor
$\min_y \operatorname{KL}(q(z|x)\|\mathcal{N}(a v_y, \tau^2 I))$, with no extra network
(the minimum is an internal enumeration over the $K$ fixed anchor buffers --- the single
QR orthonormal frame --- and consults no ground-truth label, no more than softmax
consults labels when it enumerates the $K$ logits). Whether this score is useful for OoD
detection is a hypothesis; nothing here proves it.
\end{remark}

\textbf{The problem addressed by FGIB.}
The four channels are addressed by the FGIB design as follows; two are closed by
construction, one holds conditionally, and one is only partially addressed. The precise
scope of each claim is documented in Part~II and the open items are collected in
Section~\ref{sec:limits}.
\begin{itemize}
\item \emph{Dead knob and cheap label axis} (Props.~\ref{prop:dead}, \ref{prop:flat}):
FGIB's compression term is the anchor gap $\tfrac{1}{2\tau^2}(s-a)^2$, quadratic in the
alignment scalar, and \emph{at fixed classifier scale} the Lagrangian
$\mathrm{CE}(s) + \tfrac{\beta}{2\tau^2}(s-a)^2$ is strictly convex with a unique interior
minimizer per $\beta$ --- the $\beta$-sweep is one-to-one onto the whole trade-off curve
(Props.~\ref{prop:sweep}, \ref{prop:curve}). With the scale free the infimum is $0$ for
every $\beta$; Remark~\ref{rem:c-scale} states the condition and the architectural fix.
\item \emph{Blind certificate} (Prop.~\ref{prop:flat}): with the fixed anchor
$r(z|y) = \mathcal{N}(a v_y, \tau^2 I)$ the gap
$\mathbb{E}_y[\operatorname{KL}(q(z|y)\,\|\,r(z|y))]$ can shrink only if the
\emph{posterior} moves toward the anchors --- the reference itself cannot chase the
encoder, so the force is persistent (Prop.~\ref{prop:ordering} and
Remark~\ref{rem:force}). The side head $A$, however, is a second trainable party in that
matching, and it is not part of the inference path
(Remark~\ref{rem:A-channel}).
\item \emph{Variance collapse} (Prop.~\ref{prop:variance}): the classifier consumes the
deterministic $h$ at both training and evaluation, so $\sigma^{2*} = \tau^2$ for every $\beta$
(the variance decouples from the prediction term; Lemma~\ref{lem:lagrangian}), the
side path keeps a finite covariance and the bound stays well-defined
(Remark~\ref{rem:limit}), and the trivial collapse $\mu \equiv 0$ pays
the fixed positive price $\operatorname{KL} = \tfrac12 a^2 > 0$ for every class --- the
degenerate corner is excluded by geometry rather than by the optimizer's good behavior.
This channel is closed cleanly, by construction.
\item \emph{Certificate-object mismatch} (Prop.~\ref{prop:object}): the classifier
consumes the deterministic $h$ at both training and evaluation --- there is no
training-code/evaluation-mean split --- but the bound is an upper bound on a
\emph{noisy side-path copy} of $h$, not on $I(X;H|Y)$, which no finite certificate
bounds. This channel is \emph{not} closed; what FGIB buys is that the perturbation is
side-path-only, and its size is governed by the conditioning of $A$
(Prop.~\ref{prop:conditioning}(iii)).
\end{itemize}

% =====================================================================
\section*{Part II. The FGIB theory}
% =====================================================================

\textbf{FGIB notation.}
The deterministic main path computes $h = f_{\mathrm{enc}}(X)$ and the classifier
consumes $h$ directly; the side path parameterizes $q(z|x) = \mathcal{N}(\mu(x),
\Sigma(x))$ with \emph{linear} mean head $\mu(x) = A\,h(x)$ for
$A \in \mathbb{R}^{z \times z}$ (the theory assumes $z = d$ and $A$ invertible; the
inverse transforms of Propositions~\ref{prop:fgib-ceb} and~\ref{prop:inverse} are
defined under this assumption), and a separate logvar head for $\Sigma(x)$. The
side path is \emph{not part of inference}: at evaluation the model returns the classifier
output and never evaluates the mean head (so $A$ is a trainable module that exists only
inside the regularizer --- structurally the same position CEB's backward encoder
occupies on the other side of the KL). The label-conditional prior is the
\emph{fixed} anchor
$r(z|y) = \mathcal{N}(a\,v_y, \tau^2 I)$ with \emph{anchor variance} $\tau^2 > 0$
(default $1.0$): the $v_y$ are orthonormal directions obtained by
QR-decomposing a random matrix ($K \le z$; the orthonormal factor of the random matrix),
scaled by the anchor scale $a$ (default $4.0$), stored as
non-trainable buffers. At the variance optimum of the KL the posterior variance matches
the anchor variance, $\sigma^{2*} = \tau^2$; we write $\sigma_p^2$ for this common
variance where the distinction does not matter. The orthogonality of the method lies in
this fixed frame and in the equilibrium it induces --- the KL drives the
class-conditional means $\mu_y = A\,\mathbb{E}[h|y]$ toward the orthonormal anchors
$a v_y$, so the transformed class representations become mutually orthogonal (the
aligned family below) --- while $A$ itself remains an unconstrained linear map. For
regression, fixed random-Fourier anchors
$\mu_p(y) = a\sqrt{2/z}\,\cos(2\pi\omega y + b)$ with frozen
$\omega \sim \mathcal{N}(0, 2^2)$, $b \sim U[0, 2\pi)$ replace the per-class table
($\|\mu_p(y)\| \approx a$). The loss is
$L_{\mathrm{FGIB}} = \mathbb{E}[\mathrm{CE}(c(h), y)]
+ \beta\,\mathbb{E}[\operatorname{KL}(q(z|x)\,\|\,r(z|y))]$, with gradients reaching the
main path only through the shared encoder; the zero-initialized logvar head fixes
$\sigma^2 = 1$ at initialization, which equals $\tau^2$ in the default configuration
$\tau^2 = 1$ (for general $\tau$ the theory's variance optimum $\sigma^2 = \tau^2$ is
reached by the logvar head, whose initialization is a scheduling choice), so training
starts at $\operatorname{KL} \approx \tfrac{1}{2\tau^2} a^2$.

\emph{Implementation caveat (theory--code layering).} The theory works with a fixed
orthonormal factor, an invertible linear side head $T(h) = Ah$, and per-class Gaussian
priors with anchor variance $\tau^2$. The code implements an \emph{affine} head
(\texttt{nn.Linear} with bias) --- which the theory absorbs as an idealization, the bias
being immaterial when $A$ is learned jointly --- and does not enforce orthogonality,
invertibility, or conditioning constraints on $A$. The guarantees below (orthonormal
frames, invariance, conditioning) are properties of the idealized model; the
implementation approximates them, and the guarantees do not transfer automatically
(Remarks~\ref{rem:A-channel}, \ref{rem:rank}).
\textbf{Distributional standing assumption.} Several results state that a deterministic
continuous representation carries infinite conditional information. Those statements
assume the standard continuous setting: $X \mid Y$ is non-atomic, and the induced
representation $H(X) \mid Y$ is a non-atomic continuous variable. For the finite
empirical distribution of the training data no such conclusion is claimed; there the
quantities are finite and only their growth with sample size is governed by these
results.

% ---------------------------------------------------------------
\subsection{The equivalence: the side path is a CEB bottleneck on the main path}
% ---------------------------------------------------------------

\textbf{Motivation.}
Stochastic bottlenecks such as VIB and CEB are \emph{stochastic at training time but
deterministic at evaluation time}: the classifier is trained on sampled codes
$z \sim q(z|x)$ yet evaluated on the mean $\mu_q(x)$ (Fischer 2020, Appendix A.1), so
its training and evaluation distributions differ. In the same models the cross-entropy
gradient and the KL gradient act on the same heads: the classification term can exploit
the variance (by Jensen, $\mathbb{E}_z[\mathrm{CE}(c(z), y)] \ge \mathrm{CE}(c(\mu_q),
y)$, so loosening the noise relaxes the upper bound), while the KL term pushes the
posterior back toward the prior. Finally, the regularizer constrains the noisy code $z$,
whereas inference consumes the noise-free mean --- the object being regularized is not
the object being used. FGIB resolves all three mismatches by design: the classifier
consumes the deterministic representation $h$ at both training and evaluation time, the
two objectives act on disjoint parameter sets (the classification gradient does not
enter the side heads, and the KL gradient reaches the main path only through the shared
encoder), and, by Proposition~\ref{prop:fgib-ceb} below, the side-path regularizer is,
\emph{at each fixed value of the side head}, a CEB bottleneck read in the coordinates in
which the posterior is centered at the representation that inference actually uses
(Remark~\ref{rem:A-channel} states why this is weaker than a fixed constraint on $h$).
Confining the stochasticity to a side path also keeps the variational objective
well-defined: for continuous deterministic representations the residual information
diverges (Amjad \& Geiger; Kolchinsky et al.), so a purely deterministic bottleneck
admits no finite variational bound, while the finite side-path covariance of FGIB does
(Remark~\ref{rem:limit}). Whether the robustness benefits of the stochastic bottleneck
survive the deterministic inference path is the empirical question studied in the
experiments.

\begin{lemma}[Linear invariance of the Gaussian KL]\label{lem:kl-invariance}
Let $q = \mathcal{N}(\mu_q, \Sigma_q)$ and $p = \mathcal{N}(\mu_p, \Sigma_p)$ on
$\mathbb{R}^{z}$, and let $A \in \mathbb{R}^{z \times z}$ be invertible. Define the
common pushforwards $q' = A_\# q = \mathcal{N}(A\mu_q, A\Sigma_q A^\top)$ and
$p' = A_\# p$. Then $\operatorname{KL}(q'\,\|\,p') = \operatorname{KL}(q\,\|\,p)$.
\end{lemma}

\begin{proof}
With $d = z$, the closed form reads
$\operatorname{KL}(q\,\|\,p)
= \tfrac{1}{2}\big[\operatorname{tr}(\Sigma_p^{-1}\Sigma_q)
+ (\mu_q - \mu_p)^\top \Sigma_p^{-1}(\mu_q - \mu_p)
- d + \ln|\Sigma_p| - \ln|\Sigma_q|\big]$.
Under the common map $A$ the three non-constant terms are unchanged:
\begin{align*}
\operatorname{tr}\big((A\Sigma_p A^\top)^{-1}(A\Sigma_q A^\top)\big)
&= \operatorname{tr}(A^{-\top}\Sigma_p^{-1}A^{-1} A\,\Sigma_q A^\top)
 = \operatorname{tr}(A^{-\top}\Sigma_p^{-1}\Sigma_q A^\top) \\
&= \operatorname{tr}(\Sigma_p^{-1}\Sigma_q A^\top A^{-\top})
 = \operatorname{tr}(\Sigma_p^{-1}\Sigma_q), \\[2pt]
(A\Delta\mu)^\top (A\Sigma_p A^\top)^{-1} (A\Delta\mu)
&= \Delta\mu^\top A^\top A^{-\top}\Sigma_p^{-1}A^{-1}A\,\Delta\mu
 = \Delta\mu^\top \Sigma_p^{-1}\Delta\mu, \\[2pt]
\ln|A\Sigma_p A^\top| - \ln|A\Sigma_q A^\top|
&= \big(\ln|\Sigma_p| + 2\ln|A|\big) - \big(\ln|\Sigma_q| + 2\ln|A|\big)
 = \ln|\Sigma_p| - \ln|\Sigma_q|,
\end{align*}
where $\Delta\mu = \mu_q - \mu_p$ and the trace step uses the cyclic property. \hfill\qed
\end{proof}

\begin{lemma}[Invariance of the CEB objective and the MNI point]\label{lem:mi-invariance}
For any bijective measurable map $\phi$,
$I(X;\phi(Z)|Y) = I(X;Z|Y)$ and $I(Y;\phi(Z)) = I(Y;Z)$.
Consequently the CEB objective (Eq.~5) is invariant under bijections of $Z$, and
$Z$ attains the MNI point (Eq.~1) if and only if $\phi(Z)$ does.
\end{lemma}

\begin{proof}
Mutual information is invariant under bijective transformations of either argument
(data processing inequality applied in both directions). \hfill\qed
\end{proof}

\begin{proposition}[The FGIB side path is a CEB bottleneck on the main path, at each fixed side head]\label{prop:fgib-ceb}
Let $z = d$ and let the learned linear head $T(h) = A h$ be invertible (unconstrained,
so a.s.\ so at initialization). The \emph{fixed}
prior $r(z|y) = \mathcal{N}(\mu_p(y), \Sigma_p(y))$ --- for FGIB, the orthonormal anchor
table $\mu_p(y) = a\,v_y$, $\Sigma_p = \tau^2 I$ --- is, by construction, the image
under $T$ of the canonical prior
$r'(\tilde{z}|y) = \mathcal{N}(A^{-1}\mu_p(y),\, A^{-1}\Sigma_p(y) A^{-\top})$
(since $A A^{-1} = I$), and the posterior $q(z|x) = \mathcal{N}(A\,h(x), \Sigma_q(x))$
is the image of
$q'(\tilde{z}|x) = \mathcal{N}(h(x), A^{-1}\Sigma_q(x) A^{-\top})$.
\emph{At a fixed value of $A$}, the prior is never transformed by training:
$A^{-1}$ only reads the fixed anchor geometry in the coordinates in which the posterior
is centered at $h$; during training $A$ moves, so the pulled-back prior moves with
it (Remark~\ref{rem:A-channel}).
By KL invariance under the common invertible linear map $T^{-1} = A^{-1}$
(Lemma~\ref{lem:kl-invariance}), pointwise for every $(x, y)$,
\begin{equation}
\operatorname{KL}\big(q(z|x)\,\|\,r(z|y)\big)
= \operatorname{KL}\big(q'(\tilde{z}|x)\,\|\,r'(\tilde{z}|y)\big),
\label{eq:fgib-ceb}
\end{equation}
That is, at each fixed side head the side-path regularizer is pointwise identical to a
CEB bottleneck with a backward encoder \emph{frozen at the current head}, applied to
a noisy version of the main-path representation,
$\tilde{z} = h + \tilde{\epsilon}$ with
$\tilde{\epsilon} \sim \mathcal{N}(0, A^{-1}\Sigma_q A^{-\top})$,
using the canonical prior $r'$ as that backward encoder. For FGIB the canonical backward
encoder is the frame $\mathcal{N}(a\,A^{-1} v_y, \tau^2 A^{-1}A^{-\top})$, whose
mean anchors $A^{-1}(a v_y)$ are orthonormal at scale $a$ in the induced metric
$A^{\top}A$
($\langle A^{-1}v_i, A^{-1}v_j\rangle_{A^{\top}A} = \delta_{ij}$): the orthogonality
enters through the fixed frame $\{v_y\}$ and the equilibrium it induces, not through a
constraint on $A$.
Hence the regularizer's gradient with respect to $h$ (through the shared encoder)
coincides, up to the fixed linear map $A^{\top}$, with the gradient that a CEB model
applies to its representation mean.
\end{proposition}

\begin{proof}
The map $z \mapsto \tilde{z} = A^{-1}z$ is a bijection, and
$q' = (A^{-1})_\# q$, $r' = (A^{-1})_\# r$:
$A^{-1}A\,h = h$ and $A^{-1}\Sigma_q A^{-\top}$ give the canonical posterior, and
$A^{-1}\mu_p$, $A^{-1}\Sigma_p A^{-\top}$ give the canonical prior, with
$A\,(A^{-1}\mu_p) = \mu_p$ and $A\,(A^{-1}\Sigma_p A^{-\top})\,A^{\top}
= \Sigma_p$, so $r$ is exactly the image of $r'$ under $T$.
Equation~\eqref{eq:fgib-ceb} is Lemma~\ref{lem:kl-invariance} applied to $A^{-1}$.
Sampling $\tilde{z} \sim q'$ is equivalent to
$\tilde{z} = h + \tilde{\epsilon}$ with the stated noise.
Combined with the variational bound of Fischer (2020, Eqs.~9--11), the regularizer
upper-bounds the residual information of the noisy main-path representation:
$I(X;\tilde{Z}|Y) \le \mathbb{E}[\operatorname{KL}(q'(\tilde{z}|x)\,\|\,r'(\tilde{z}|y))]$.
\hfill\qed
\end{proof}

\begin{remark}[Trainable-implementation caveat: the side head is a second trainable party in the matching]\label{rem:A-channel}
The identity of Proposition~\ref{prop:fgib-ceb} holds pointwise at each value of the
side head, so the asymmetry advertised against CEB --- one trainable side versus two ---
is weaker than it looks. In the idealized theory $A$ may be treated as fixed
(orthonormal or any fixed invertible map); in the \emph{trainable implementation} $A$ is
free and \emph{unused at inference}:
like CEB's backward encoder $b(z|y)$, it lives inside the KL and can lower it without
moving the deployed representation $h$. Since $h$ is linear-separable by the jointly
trained classifier, the linear fit of the anchor target $a v_y$ on $h$ fits well, and
much of that fit can be absorbed by $A$ alone. The variance term
$\tfrac{z}{2}(\sigma^2/\tau^2 - 1 - \ln(\sigma^2/\tau^2))$ does not depend on $A$ at
all, so nothing in
the regularizer resists $A$ drifting toward a projector onto the $K$-dimensional
between-class subspace with a vanishing component on within-class directions. Full
collapse is blocked --- with invertible $A$, $\operatorname{KL} = 0$ forces $h$ itself
to be class-constant --- but the approach to the rank-$K$ regime is not blocked, and
that regime is exactly where the conditioning failure and the perturbation blow-up of
Propositions~\ref{prop:conditioning} and~\ref{prop:inverse} live. This is the one
channel of Part~I that FGIB has not closed; Section~\ref{sec:limits} collects it.
\end{remark}

\begin{remark}[Family closure]\label{rem:closure}
Lemma~\ref{lem:kl-invariance} holds for the underlying Gaussian distributions. Within a
restricted variational family the equality of Proposition~\ref{prop:fgib-ceb} is exact
only if the family is closed under the inverse map $A^{-1}$: the full-covariance family
is closed (and is the parameterization used by Fischer 2020); the diagonal family is not
--- for diagonal $q, r$ the transformed pair is diagonal only if $A^{-1}$ preserves
diagonality (e.g., permutations and sign flips) and approximate otherwise. The diagonal
implementation of the codebase is therefore covered at the distribution level, with the
family approximation confined to the noise covariance $A^{-1}\Sigma_q A^{-\top}$
(Remark~\ref{rem:orth}).
\end{remark}

\begin{remark}[Rank deficiency]\label{rem:rank}
If $z < d$ or $A$ is rank-deficient, only the projection of $h$ onto the row space of
$A$ (the image of $A^{\top}$ in $h$-space) receives gradient from the regularizer; the
null-space component is unconstrained.
In the default configuration $z = d$ and a randomly initialized $A$ is almost surely
invertible, but invertibility carries no conditioning margin --- near-singularity is
enough to produce the effects of Proposition~\ref{prop:conditioning}(iii), and it is the
direction the gap gradient points in (Remark~\ref{rem:A-channel}). If $A \in
\mathbb{R}^{z \times d}$ has full column rank ($z \ge d$), a left inverse $B$ with
$BA = I$ exists and the data-processing inequality for KL gives, per $(x,y)$,
$\operatorname{KL}(q\|r) \ge \operatorname{KL}(B_\# q\,\|\,B_\# r)
\ge I(X;\,BZ\,|\,Y)$ --- a valid lower-level bound, but one whose equality
(invariance) is lost, and whose pulled-back prior again depends on $B$. The exact
invariance argument of Proposition~\ref{prop:fgib-ceb} genuinely requires the invertible
linear map $T(h) = Ah$ with $z = d$.
\end{remark}

\begin{remark}[Noise-free limit]\label{rem:limit}
The equivalence of Proposition~\ref{prop:fgib-ceb} should be understood at the level of
the variational objective, not as a mutual-information limit: as the side-path noise
vanishes, $I(X;Z|Y)$ diverges for continuous deterministic $H$ under the standing
non-atomic assumption (the pathology of
deterministic networks discussed by Amjad \& Geiger and Kolchinsky et al.), and the
variational KL bound becomes loose. Keeping the side-path covariance finite --- without
injecting noise into the inference path --- is precisely what keeps the regularizer
well-defined; this is the design rationale of FGIB.
\end{remark}

% ---------------------------------------------------------------
\subsection{The sweep}
% ---------------------------------------------------------------

\textbf{The degeneracy to be addressed (restated in our terms).}
When $Y = f(X)$, the IB curve saturates the data-processing bound and is piecewise
linear, $F(r) = \min\{r, H(Y)\}$ (Caveats paper, Section~3). Two consequences (their
Caveat~1): (i)~by the DPI, $L^{\mathrm{IB}}_\beta(T) = I(Y;T) - \beta I(X;T)
\le (1-\beta)I(Y;T) \le (1-\beta)H(Y)$, and $T_{\mathrm{copy}} := f(X) = Y$ attains the
bound for \emph{every} $\beta \in [0,1]$: all $\beta \in (0,1)$ are pinned at the single
corner point $\langle H(Y), H(Y)\rangle$ of the curve; (ii)~conversely, at $\beta = 1$
every point of the increasing part simultaneously maximizes $L^{\mathrm{IB}}_1$. The
mechanism is exposed by writing $\max_T L^{\mathrm{IB}}_\beta = \max_r [F(r) - \beta r]$
(the Legendre transform of $-F$): on the linear increasing part every point has the
\emph{same} derivative $F'(r) = 1$, hence the same supporting slope --- a Lagrangian of
the form ``prediction $-\beta\cdot{}$compression'' sweeps a Pareto curve only if each
point of the curve has a \emph{unique} supporting slope, i.e.\ the curve is strictly
concave. The dIB variant fails identically: on the hard-clustering family $T = g(Y)$ one
has $H(T|Y) = 0$, hence $I(Y;T) = H(T)$, so $L^{\mathrm{dIB}}_\beta = (1-\beta)H(T)$
--- a \emph{monotone} function of the single compression scalar $H(T)$, maximized by the
finest partition for every $\beta < 1$. The squared fixes repair this by making the
objective strictly concave in the compression scalar: $F'(r)/(2r) = \beta$ is strictly
decreasing, so each $\beta$ selects a unique $r$ (Eq.~8 argument), and squared-dIB has
the unique maximizer $H(T) = 1/(2\beta)$ (their Eq.~A20). By Proposition~\ref{prop:dead}
this degeneracy is inherited verbatim by CEB ($\beta_{\mathrm{IB}} = 1+\gamma > 1$),
which is the first channel of Part~I.

\textbf{Reduced model of FGIB.}
The KL depends on the side mean $\mu$ alone, and the classifier term is invariant under
the coordinate change $\tilde w_j = A^{-\top} w_j$ (since $w_j \cdot h =
\tilde w_j \cdot \mu$): with $A$ invertible ($z = d$ in the default configuration, so a
random $A$ is a.s.\ invertible), every equilibrium in $(h, w)$ coordinates maps
bijectively to one in $(\mu, \tilde w)$ coordinates. We therefore analyze the reduced
model $\mu(x) = h(x)$ without loss of generality.
Consider the \emph{class-symmetric aligned family}: the encoder maps every input of
class $y$ to $h_y = s\,v_y$ ($s \ge 0$), the classifier is held at the symmetric
configuration $w_j = c\,v_j$, $b_j = 0$ with fixed scale $c > 0$, and the posterior is
$q_y = \mathcal{N}(s v_y, \sigma^2 I)$. The family is the natural equilibrium manifold of
the two forces: the KL gradient
$\nabla_h \operatorname{KL} = (\mu - a v_y)/\tau^2$ acts exactly
along $v_y$, and the classification force on the class center acts along $v_y$ at the
symmetric configuration; the ansatz is a \emph{working approximation}: it captures the
direction of both forces, and is motivated as exact in the $\beta$-dominant regime, but
that exactness is not proved here --- the cross terms are of order $O(c/\beta)$ and are
idealized away (see Remark~\ref{rem:caveats}).

\begin{lemma}[FGIB Lagrangian on the aligned family]\label{lem:lagrangian}
On the aligned family,
\begin{equation}
\operatorname{KL}(s, \sigma^2)
= \tfrac{z}{2}\big(\sigma^2/\tau^2 - 1 - \ln(\sigma^2/\tau^2)\big)
+ \tfrac{1}{2\tau^2}(s-a)^2,
\qquad
\mathrm{CE}(s) = \ln\!\big(1 + (K-1)e^{-cs}\big),
\label{eq:kl-ce}
\end{equation}
and the optimizer over $\sigma^2$ is $\sigma^2 = \tau^2$ for \emph{every} $\beta$, since
the CE term does not depend on $\sigma^2$ (the classifier consumes the deterministic
$h$, not $z$).
\end{lemma}

\begin{proof}
With $h_y = s v_y$ and $w_j = c v_j$, the logits are $c\,\langle v_j, s v_y\rangle
= cs\,\delta_{jy}$, so $p_y = e^{cs}/(e^{cs} + K - 1)$ and
$\mathrm{CE} = -\ln p_y = \ln(1 + (K-1)e^{-cs})$.
For the KL, the diagonal closed form with $\Sigma_q = \sigma^2 I$,
$\Sigma_p = \tau^2 I$ reads
$\tfrac{1}{2}\sum_d\big[\sigma^2/\tau^2 + (\mu_d - a v_{yd})^2/\tau^2 - 1
- \ln(\sigma^2/\tau^2)\big]$, giving
Eq.~\eqref{eq:kl-ce} since $\|v_y\| = 1$. The variance term is non-negative with a
unique minimum at $\sigma^2 = \tau^2$, and $\mathrm{CE}$ is independent of $\sigma^2$.
\hfill\qed
\end{proof}

\begin{proposition}[Unique interior optimum per $\beta$ at fixed classifier scale --- the sweep]\label{prop:sweep}
\emph{At a fixed classifier scale $c$}, for every $\beta > 0$, the FGIB Lagrangian on
the aligned family,
\begin{equation}
L_\beta(s) := \ln\!\big(1 + (K-1)e^{-cs}\big) + \tfrac{\beta}{2\tau^2}(s-a)^2,
\label{eq:fgib-lag}
\end{equation}
is strictly convex in $s$ and has a unique minimizer $s^*(\beta) > a$, given implicitly
by
\begin{equation}
\frac{\beta}{\tau^2}\,(s-a) = c\,\varphi(cs),
\qquad
\varphi(u) := \frac{(K-1)e^{-u}}{1 + (K-1)e^{-u}} \in (0,1) .
\label{eq:stationarity}
\end{equation}
The map $\beta \mapsto s^*(\beta)$ is continuous and strictly decreasing, with
$s^*(\beta) \to a$ as $\beta \to \infty$ and $s^*(\beta) \to \infty$ as $\beta \to 0$.
Consequently the map $\beta \mapsto \big(\operatorname{KL}(s^*), \mathrm{CE}(s^*)\big)$
is one-to-one onto the entire curve
$\{\big(\tfrac{1}{2\tau^2}(s-a)^2,\; \ln(1+(K-1)e^{-cs})\big) : s \in (a, \infty)\}$:
no two
$\beta$ share the same optimum, and every interior point of the trade-off curve is
attained --- the exact failure mode of $L^{\mathrm{IB}}_\beta$ (all $\beta$ pinned at
$T_{\mathrm{copy}}$) is absent. The fixed-$c$ condition is load-bearing and is analyzed
in Remark~\ref{rem:c-scale}.
\end{proposition}

\begin{proof}
Compute $\mathrm{CE}'(s) = -c\varphi(cs)$ and
$\mathrm{CE}''(s) = c^2 (K-1) e^{-cs}/(1 + (K-1)e^{-cs})^2 > 0$, so CE is strictly
convex and strictly decreasing in $s$; the compression term $\tfrac{1}{2\tau^2}(s-a)^2$
is strictly convex. Hence $L_\beta''(s) = \mathrm{CE}''(s) + \beta/\tau^2 > 0$:
$L_\beta$ is
strictly convex with at most one minimizer, characterized by
$0 = L_\beta'(s) = -c\varphi(cs) + \beta(s-a)/\tau^2$, i.e.\
Eq.~\eqref{eq:stationarity}.
The left side $\beta(s-a)/\tau^2$ is strictly increasing from $-\beta a/\tau^2$
(at $s=0$) through $0$
(at $s=a$) to $+\infty$; the right side $c\varphi(cs)$ is strictly decreasing in $s$
with $c\varphi(cs) \in (0, c)$. The unique crossing therefore satisfies $s^* > a$.
Differentiating Eq.~\eqref{eq:stationarity} with respect to $\beta$ gives
$(s-a)/\tau^2 + (\beta/\tau^2) s' = c^2\varphi'(cs)\,s'$ with $\varphi' < 0$, so
$s' = -(s-a)/(\beta - c^2\tau^2\varphi'(cs)) < 0$: $s^*$ is strictly decreasing.
The limits follow by monotonicity: as $\beta \to \infty$ the left side can only match
the bounded right side when $s \to a$; as $\beta \to 0$ the crossing requires
$\varphi(cs)\to 0$, i.e.\ $s \to \infty$. \hfill\qed
\end{proof}

\begin{proposition}[The trade-off curve is strictly concave in the compression--prediction plane]\label{prop:curve}
On the curve of Proposition~\ref{prop:sweep} (fixed classifier scale $c$),
\begin{equation}
\frac{d\,\mathrm{CE}}{d\,\operatorname{KL}}
= \frac{\mathrm{CE}'(s)}{(s-a)/\tau^2}
= -\frac{\tau^2\,c\varphi(cs)}{s-a},
\label{eq:slope}
\end{equation}
which is strictly increasing from $-\infty$ (as $s \to a^+$) to $0^-$ (as $s \to
\infty$). Thus the curve $\operatorname{KL} \mapsto -\mathrm{CE}$ (compression versus
achievable prediction, on the KL--CE surrogate plane of
Remark~\ref{rem:surrogate}) is strictly concave, and each of its points has a \emph{unique}
supporting slope $-\beta$. By the criterion of the Caveats paper (their Fig.~2 right
panel and the $F'(r)/(2r)$ argument of Eq.~8), this is precisely the property that makes
the Lagrangian $-\mathrm{CE} - \beta\,\operatorname{KL}$ sweep the full \emph{surrogate}
curve on the aligned family at fixed scale: the convexity of the compression term
supplies the strict-concavity condition \emph{without squaring it}. This is a statement
about the KL--CE surrogate of Proposition~\ref{prop:sweep}, not about the Caveats
information plane (Remark~\ref{rem:surrogate}).
\end{proposition}

\begin{proof}
$\varphi$ is positive and strictly decreasing and $s-a$ is positive and strictly
increasing, so $\tau^2 c\varphi(cs)/(s-a)$ is strictly decreasing in $s$; the signs give
the stated limits. The stationarity condition Eq.~\eqref{eq:stationarity} is exactly
$-d\mathrm{CE}/d\operatorname{KL} = \beta$, i.e.\ $\beta$ is the supporting slope at the
optimum. \hfill\qed
\end{proof}

\begin{remark}[What the sweep sweeps: a KL--CE surrogate curve, not the IB curve]\label{rem:surrogate}
Propositions~\ref{prop:sweep} and~\ref{prop:curve} prove strict convexity of the
FGIB Lagrangian and one-to-one parameterization by $\beta$ \emph{on the class-symmetric
aligned family, at fixed classifier scale}, in the plane
$\big(\operatorname{KL}_{\mathrm{anchor}},\, \mathrm{CE}\big)$. That plane is not the
Caveats paper's IB plane $\big(I(X;T),\, I(Y;T)\big)$, and the claim does not extend to
it. On the aligned family the posterior is class-constant by construction, so
$I(X;Z|Y) = 0$ for every $s$, while the deterministic class prototypes $h_y = s v_y$
already retain the label completely for any $s > 0$ ($I(Y;H) = H(Y)$); what varies
along the sweep is the anchor mismatch and the classifier margin, not a proved
compression of mutual information. Whether the deployed representation's residual
$I(X;H|Y)$ moves monotonically with $\beta$ is a separate, empirical question; the
theory guarantees a strictly concave surrogate, not a recovered IB curve. This
qualification is carried into Section~\ref{sec:limits}.
\end{remark}

\begin{remark}[The classifier scale is load-bearing, and the current implementation leaves it free]\label{rem:c-scale}
If $c$ is free, the pair $(s, c) = (a, c)$ with $c \to \infty$ sends
$\operatorname{KL} = \tfrac{1}{2\tau^2}(s-a)^2 = 0$ and
$\mathrm{CE} = \ln(1 + (K-1)e^{-ca}) \to 0$ simultaneously (for $a > 0$): the infimum
of $L_\beta$ over the aligned family is $0$ for \emph{every} $\beta$, no minimizer is
attained, and the one-to-one sweep of Proposition~\ref{prop:sweep} has no global
counterpart. (At $a = 0$ the margin growth alone cannot drive CE to zero.) This is
the FGIB analogue of the dead knob --- the same one-corner attractor, now arrived at by
growing the margin rather than by the linearity of the deterministic IB curve.
Three consequences.
\emph{(a) The present implementation is exposed.} The classifier is
\texttt{nn.Linear} trained by Adam with no weight decay, so $c$ is unconstrained. On
separable data the CE gradient drives $\|w_y\|$ upward; it does not diverge in finite
time (the gradient vanishes as $p_y \to 1$, giving logistic growth), but at any finite
scale $c$ the equilibrium $s^*(\beta; c)$ sits on the curve of
Proposition~\ref{prop:sweep} \emph{for that $c$}, and $s^*$ drifts toward $a$ as $c$
grows.
\emph{(b) Testable prediction.} The $\beta$-dependence of the equilibrium should decay
with training length; since the experiments use early stopping with $\texttt{patience} =
10$, part of any observed $\beta$-effect is plausibly an early-stopping effect. This can
be tested by training well past saturation and comparing the converged KL across
$\beta$.
\emph{(c) Architectural fix.} The device that satisfies the fixed-$c$ assumption
\emph{by construction} is a cosine classifier: normalized features, normalized rows,
and a fixed temperature (scale). Weight normalization with a learnable gain does not
guarantee a bounded scale, and weight decay only replaces the free-scale problem with a
joint optimization problem over $(s, c)$ that must be re-derived --- it mitigates the
divergence but does not make Proposition~\ref{prop:sweep} exact. Under a fixed-norm
classifier with fixed temperature the proposition holds as stated, at the price of one
architecture change.
\end{remark}

\textbf{Where the convexity comes from (mechanism).}
Compare the three objectives on their natural solution families, as functions of the
single compression scalar:
\begin{itemize}
\item \emph{dIB on hard clusterings} (Caveats Appendix~B): $L^{\mathrm{dIB}}_\beta
= (1-\beta)H(T)$ --- compression $H(T)$ and prediction $I(Y;T) = H(T)$ are the
\emph{same} scalar, so the objective is monotone in it and the optimum sits at the
endpoint for every $\beta < 1$. Squared-dIB restores interior optima by hand,
$H(T) - \beta H(T)^2$ with unique maximizer $H(T) = 1/(2\beta)$.
\item \emph{Standard IB} (deterministic): the curve is linear ($F'(r) \equiv 1$), every
point shares one supporting slope, and the corner $\langle H(Y), H(Y)\rangle$ is
selected for all $\beta \in (0,1)$. Squared-IB convexifies via $F'(r)/(2r)$ strictly
decreasing.
\item \emph{FGIB}: the compression term is $\operatorname{KL}(q_y \| r_y) =
\tfrac{1}{2\tau^2}\|\mu_y - a v_y\|^2 + {}$(variance terms), equal to
$\tfrac{1}{2\tau^2}(s-a)^2$
at the variance optimum --- \emph{quadratic} in the alignment scalar $s$, hence strictly
convex, because the KL is measured against a \emph{fixed} reference $a v_y$:
$\operatorname{KL}(\mathcal{N}(\mu,\cdot)\|\mathcal{N}(\mu_p,\cdot))
= \tfrac{1}{2\tau^2}\|\mu - \mu_p\|^2 + {}$(variance terms) is quadratic in the
displacement
$\mu - \mu_p$ whenever $\mu_p$ does not move. Entropy-based compression ($H(T)$,
$I(X;T)$) is self-referential --- it depends only on the representation's own
distribution and enters linearly in the clustering probabilities; the fixed anchor
supplies an absolute reference frame and turns compression into a squared distance. The
convexification that squared-IB installs by squaring the compression term is thus built
into FGIB's fixed-anchor KL. Together with the strictly convex prediction term
(Lemma~\ref{lem:lagrangian}), the Lagrangian is strictly convex with a unique interior
minimizer per $\beta$ \emph{at fixed classifier scale} (Proposition~\ref{prop:sweep}),
and the resulting curve is strictly concave (Proposition~\ref{prop:curve}). With the
scale free, the same quadratic gap is driven to zero by margin growth alone
(Remark~\ref{rem:c-scale}) --- the convexification is necessary but not sufficient
without the scale condition.
\end{itemize}

\begin{remark}[The compressed endpoint is non-collapsed but still a label code --- Caveat 2 remains]\label{rem:endpoint}
VIB's fixed marginal $m(z) = \mathcal{N}(0, I)$ places the $\operatorname{KL}=0$ point
at $\mu = 0$ for \emph{all} classes: the fully compressed solution is the trivial
collapse with $\mathrm{CE} = \ln K$ (chance level), and squaring (SVIB) leaves this
corner in place since $\operatorname{KL}^2$ is stationary at $\operatorname{KL}=0$.
FGIB's per-class orthonormal anchors make the $\operatorname{KL}=0$ point a
\emph{class-separated} configuration: $\mu_y = a v_y$ with
$\langle v_y, v_{y'}\rangle = \delta_{yy'}$, so at the corner
$\mathrm{CE}(a) = \ln(1 + (K-1)e^{-ca})$, which is bounded and tends to $0$ as the
anchor scale $a$ grows, and $a = 0$ recovers the VIB corner. But class-separation does
not buy non-triviality in the sense of the Caveats paper: their Caveat~2 identifies as
uninteresting the variables on the manifold $T_\alpha$ that mix a constant map with an
exact copy of $Y$, compressing by ``forgetting the input with some probability'' rather
than by merging perceptually similar classes --- and FGIB's compressed corner
$\mu_y = a v_y$ is a deterministic function of $Y$ alone, i.e.\ a label code of exactly
that kind. The orthonormal frame is in fact \emph{opposed} to the kind of compression
Caveat~2 holds up as interesting: it places every pair of classes at the same distance
and can never express ``merge collie with retriever, keep both far from teapot''. What
the anchors do is choose \emph{one specific} label-code geometry --- class means on a
scaled simplex, a neural-collapse-like configuration --- instead of leaving it to the
optimizer, and refuse the chance-level corner. The compressed endpoint is thus designed,
not certified, and any claim that the method compresses \emph{input} semantics at that
endpoint is exactly the claim Caveat~2 rules out.
\end{remark}

\begin{remark}[Residual-information interpretation (CEB)]\label{rem:residual}
For any backward encoder, $\mathbb{E}[\operatorname{KL}(q(z|x)\|r(z|y))]
= I(X;Z|Y) + \mathbb{E}_y[\operatorname{KL}(q(z|y)\|r(z|y))]
\ge I(X;Z|Y)$ (Fischer 2020, Eqs.~9--11, 20; Lemma~\ref{lem:gap}). FGIB is therefore a
CEB bottleneck with a \emph{fixed} backward encoder: the regularizer upper-bounds the
residual information, and on the aligned family the residual vanishes identically
($q(z|x)$ depends on $y$ alone), so the regularizer equals the anchor-matching gap
$\tfrac{1}{2\tau^2}(s-a)^2$. The quadratic mechanism of Proposition~\ref{prop:sweep} is
inherited from this gap term, which exists only because the anchor is not allowed to
follow the encoder.
\end{remark}

\begin{remark}[Variance decoupling and training--evaluation consistency]
Since the classifier consumes the deterministic $h$ (and the KL does not see the
classifier), the posterior variance decouples: $\sigma^{2*} = \tau^2$ for all $\beta$, and
the trade-off is purely geometric (mean alignment). This contrasts with VIB, where the
classifier consumes the sampled $z$: the noise couples into the prediction term (by
Jensen, $\mathbb{E}_z[\mathrm{CE}(c(z),y)] \ge \mathrm{CE}(c(\mu),y)$), the rate
$R \ge I(X;Z)$ over-counts compression, and training-time sampling versus
evaluation-time mean creates a distribution mismatch (Fischer 2020, Appendix~A.1).
\end{remark}

\begin{remark}[Idealizations]\label{rem:caveats}
The class-symmetric ansatz idealizes the equilibrium geometry; it is exact in the
$\beta$-dominant regime and captures the direction of the KL force in general. For
regression, the continuous RFF anchors satisfy $\|\mu_p(y)\| \approx a$, so the KL again
contains the quadratic displacement term
$\tfrac{1}{2}\|\mu - \mu_p(y)\|^2 + {}$(variance terms) --- the same convexity mechanism
carries over, again subject to the fixed-scale condition of Remark~\ref{rem:c-scale}
(the regression classifier has no margin scale; its norm is likewise unconstrained in
the implementation). Training starts at $\operatorname{KL} \approx a^2/2$
(zero-initialized logvar head, $\sigma = 1$, $\mu \approx 0$), so the anchor scale sets
the initial regularization pressure; larger $a$ delays the approach to $s^*$. The
finite-scale caveat formerly in this remark is stated, with its fix, in
Remark~\ref{rem:c-scale}.
\end{remark}

% ---------------------------------------------------------------
\subsection{The certificate}
% ---------------------------------------------------------------

\begin{lemma}[Gap identity]\label{lem:gap}
For any conditional $r(z|y)$,
\begin{equation}
\mathbb{E}\big[\operatorname{KL}(q(z|x)\,\|\,r(z|y))\big]
= I(X;Z|Y) + \mathbb{E}_y\big[\operatorname{KL}(q(z|y)\,\|\,r(z|y))\big] .
\label{eq:gap}
\end{equation}
\end{lemma}

\begin{proof}
$\mathbb{E}[\operatorname{KL}(q(z|x)\|r(z|y))] = \mathbb{E}[\ln q(z|x) - \ln r(z|y)]$
and $I(X;Z|Y) = \mathbb{E}[\ln q(z|x) - \ln q(z|y)]$ (Fischer 2020, Eqs.~9--10), where
both expectations are over $p(x,y)q(z|x)$; subtracting,
$\mathbb{E}[\ln q(z|y) - \ln r(z|y)] = \mathbb{E}_y[\operatorname{KL}(q(z|y)\|r(z|y))]$
since the argument depends on $z, y$ only. \hfill\qed
\end{proof}

\begin{lemma}[Moment matching]\label{lem:mm}
Within the per-class Gaussian family with \emph{diagonal covariance}, the gap of
Lemma~\ref{lem:gap} is
minimized by the moment-matched Gaussian $b^*_y = \mathcal{N}(m_y, S_y)$ with
$m_y = \mathbb{E}\,q(z|y)$, $S_y = \operatorname{diag}\operatorname{Var} q(z|y)$.
\end{lemma}

\begin{proof}
$\operatorname{KL}(p\|q)$ as a function of $q$ is minimized by matching the first two
moments of $p$ (standard; for Gaussians the minimizer is unique). \hfill\qed
\end{proof}

\begin{proposition}[Variational ordering: CEB's bound is tighter in value]\label{prop:ordering}
For every fixed encoder $q(z|x)$,
$\mathrm{ReX}(\theta^*) \le \mathrm{ReX}_F$ where $\theta^*$ is CEB's optimal backward
encoder, and the excess is the \emph{anchor gap}
\begin{equation}
\Delta := \mathrm{ReX}_F - \mathrm{ReX}(\theta^*)
= \mathbb{E}_y\big[\operatorname{KL}(q(z|y)\,\|\,N(a v_y, \tau^2 I))\big]
- \mathbb{E}_y\big[\operatorname{KL}(q(z|y)\,\|\,N(m_y, S_y))\big] \ge 0 ,
\label{eq:anchor-gap}
\end{equation}
with equality iff the anchors equal the moment-matched marginals. Here $\theta^*$ ranges
over the \emph{ideal} CEB backward encoder family --- per-class Gaussians with free mean
and \emph{free diagonal covariance} (the family of Lemma~\ref{lem:mm}, matching the
diagonal implementation), i.e.\ any conditional $r(z|y)$ of that family the theory can
express; the code's
single affine \texttt{prior\_net} is only an approximation of this family, and the
ordering is a statement about the idealized CEB, not about the implemented
parameterization.
On the aligned family of the sweep analysis ($q(z|y) = \mathcal{N}(s v_y, \sigma^2 I)$,
so $I(X;Z|Y) = 0$ and $b^* = q(z|y)$ exactly),
\begin{equation}
\Delta = \tfrac{1}{2\tau^2}(s-a)^2 + \tfrac{z}{2}\big(\sigma^2/\tau^2 - 1 - \ln(\sigma^2/\tau^2)\big) ,
\label{eq:aligned-gap}
\end{equation}
i.e.\ $\mathrm{ReX}_F = \Delta$ and $\mathrm{ReX}(\theta^*) = 0$: FGIB's bound value
\emph{equals} the sweep-driving compression term of Proposition~\ref{prop:sweep}, while
CEB's bound collapses to the (zero) residual.
\end{proposition}

\begin{proof}
The anchor $N(a v_y, \tau^2 I)$ is a member of CEB's per-class Gaussian family, so the
minimum over the family is no larger than the value at the anchor: the ordering and
Eq.~\eqref{eq:anchor-gap} follow from Lemmas~\ref{lem:gap} and~\ref{lem:mm}.
On the aligned family the marginal is itself Gaussian, so $b^* = q(z|y)$ gives gap $0$,
$I(X;Z|Y) = 0$ (the posterior is class-constant), and
$\operatorname{KL}(\mathcal{N}(s v_y, \sigma^2 I)\|\mathcal{N}(a v_y, \tau^2 I))
= \tfrac{1}{2\tau^2}(s-a)^2 + \tfrac{z}{2}(\sigma^2/\tau^2-1-\ln(\sigma^2/\tau^2))$.
\hfill\qed
\end{proof}

\begin{remark}[Tightness--force trade-off]\label{rem:force}
Proposition~\ref{prop:ordering} is the price of the fixed anchor, and it is paid for a
reason: in CEB the gap is a function of the \emph{backward encoder alone} ---
$b_\theta \to q(z|y)$ eliminates it with the encoder held fixed, so CEB's bound can
become tight while the representation stays uncompressed: its tightness certifies
nothing about compression, and the gap is an unobservable optimization residual that
tracks the forward encoder. FGIB's gap can shrink only if the \emph{posterior} moves
toward the anchors (the anchor itself cannot move), so the force is persistent and, on
the aligned family, the bound value \emph{is} the anchor-matching gap
$\tfrac{1}{2\tau^2}(s-a)^2$, the sweep-driving compression term
(Eq.~\eqref{eq:aligned-gap}). The qualifier is necessary: ``the posterior'' includes
the side head $A$, which is trainable and unused at inference, so part of the movement
can be absorbed without the deployed representation $h$ moving
(Remark~\ref{rem:A-channel}). On the aligned family, where $A$ is fixed by the ansatz,
the equivalence is exact; off it, bound tightness is monotonically equivalent to
alignment of the posterior, not to compression of the deployed $h$.
\end{remark}

\begin{proposition}[The inverse linear transform: the object of the bound]\label{prop:inverse}
Let the side head $T(h) = Ah$ be \emph{invertible} ($z = d$, unconstrained), and let
the trained posterior be isotropic, $\Sigma(x) = \sigma_p^2 I$ with $\sigma_p^2 = \tau^2$
the anchor variance (the KL variance term is minimized at $\sigma^2 = \tau^2$, so this
holds at the variance optimum).
Then $\tilde z := A^{-1} z$ gives
$q'(\tilde z|x) = \mathcal{N}(h(x), \tau^2 (A^{\top}A)^{-1})$ --- the posterior is
centered at the \emph{inference representation} $h$ with side noise
$\tau^2 (A^{\top}A)^{-1}$ --- and the transformed anchors
$r'(\tilde z|y) = \mathcal{N}(a\, A^{-1} v_y, \tau^2 (A^{\top}A)^{-1})$ form an
orthonormal frame at scale $a$ in the induced metric $A^{\top}A$
($\langle A^{-1}v_i, A^{-1}v_j\rangle_{A^{\top}A} = \delta_{ij}$: the fixed frame
$\{v_y\}$ read in canonical coordinates). By KL invariance under invertible linear maps
(Lemma~\ref{lem:kl-invariance}),
\begin{equation}
\mathrm{ReX}_F
= I(X;\, h + \tau\,(A^{\top}A)^{-1/2}\varepsilon\,|Y)
+ \mathbb{E}_y\big[\operatorname{KL}\big(q'(\tilde z|y)\,\|\,r'(\tilde z|y)\big)\big],
\qquad \varepsilon \sim \mathcal{N}(0, I) .
\label{eq:object}
\end{equation}
The bound's object is the inference representation perturbed by \emph{known,
side-path-only} noise (its covariance is a model parameter), and the anchor geometry is
preserved exactly under the transform, in the induced metric. The perturbation,
however, is small only if the side head is well-conditioned: its magnitude is governed
by $A^{\top}A$, which is unconstrained
(Proposition~\ref{prop:conditioning}(iii)).
\end{proposition}

\begin{proof}
$A^{-1}A = I$ and $A^{-1} z = h + A^{-1}\varepsilon$ for $\varepsilon \sim
\mathcal{N}(0, \Sigma)$; $A^{-1}\Sigma A^{-\top} = \tau^2 A^{-1}A^{-\top}
= \tau^2 (A^{\top}A)^{-1}$ for $\Sigma = \tau^2 I$; the prior covariance maps to
$A^{-1} (\tau^2 I) A^{-\top} = \tau^2 (A^{\top}A)^{-1}$; and
$\langle A^{-1}v_i, A^{-1}v_j\rangle_{A^{\top}A}
= v_i^{\top} A^{-\top} A^{\top} A\, A^{-1} v_j = \delta_{ij}$.
Equation~\eqref{eq:object} is Lemma~\ref{lem:gap} applied to the transformed pair.
\hfill\qed
\end{proof}

\begin{proposition}[Orthogonal-factor conditioning: the reference is fixed and well-conditioned]\label{prop:conditioning}
FGIB's anchor frame $V = [v_1, \dots, v_K]$ is the orthonormal factor of a random matrix
(QR: $G = QR$, $V := Q$, $V^{\top}V = I_K$), and the anchor covariance is
$\tau^2 I$ with $\tau^2 > 0$.
Then, for every anchor scale $a \ge 0$ and throughout training:
\begin{enumerate}
\item[(i)] the \emph{reference} of the matching --- the \emph{unscaled} mean frame $V$
and the covariance $\tau^2 I$ --- is fixed (non-trainable buffers) and has condition
number $1$: the certificate is a distance to a compact, non-degenerate,
class-separated (for $a > 0$) target. (At $a = 0$ the mean matrix $a V$ is zero and
has no condition number; the condition-$1$ statement attaches to $V$ and $\tau^2 I$,
not to $a V$.)
\item[(ii)] in side-mean coordinates the certificate's displacement metric is a
scaled identity: the KL's displacement term is $\tfrac{1}{2\tau^2}\|\mu - a v_y\|^2$
with gradient $(\mu - a v_y)/\tau^2$;
\item[(iii)] on the shared-encoder side the metric is $A^{\top}A$, where $A$ is a
learnable unconstrained map, so its \emph{global} condition number is not controlled by
construction --- and the aligned family does not control it either. An earlier version
inferred condition number $1$ on the class-mean subspace from the identity
$h_i^{\top} A^{\top} A\, h_j = s^2\delta_{ij}$; that inference is invalid, since
$\{h_i\}$ is not a Euclidean orthonormal basis. Counterexample: with
$A = \operatorname{diag}(\varepsilon, 1)$, $h_1 = (1/\varepsilon, 0)$,
$h_2 = (0, 1)$ and $s = 1$, one has $A h_i = v_i$ yet
$\kappa(A^{\top}A) = \varepsilon^{-2}$ --- and the inverse-transform perturbation of
Proposition~\ref{prop:inverse} is correspondingly $\varepsilon^{-2}$-large on the first
direction. The counterexample is not adversarial: it is the direction the gap gradient
points in, since the KL is minimized by fitting the anchor target on $h$ in least
squares, which suppresses within-class directions
(Remark~\ref{rem:A-channel}). Any guarantee on this side therefore requires an explicit
constraint on $A$ (orthogonal parameterization, weight norm, or the like) or an
empirical measurement of $\kappa(A^{\top}A)$ over training.
\end{enumerate}
CEB's backward encoder $b_\theta(z|y) = \mathcal{N}(m_\theta(y), S_\theta(y))$ is an
unconstrained linear map of the label: \emph{both} sides of its matching are trainable,
its covariance diagonal can lie anywhere in the clamp $[e^{-10}, e^{10}]$ (condition
number up to $e^{20}$), its means can be rank-deficient (class-merged codes, or the
collapse solution $m_\theta \equiv 0$), and the variance race drives $S_\theta$ to the
clamp, degenerating the certificate into a mean-matching force of weight $1/s = e^{10}$
(Prop.~\ref{prop:variance}, subject to Remark~\ref{rem:p3-status}). What the fixed QR
reference guarantees for FGIB is therefore the \emph{reference} side only: a compact,
non-degenerate, class-separated target whose conditioning cannot degrade. On the
encoder side FGIB's bound is exposed to the same degeneration mechanism CEB's is, with
$A$ playing the role of the second trainable party.
\end{proposition}

\begin{proof}
$V^{\top}V = I_K$: all singular values of $V$ are $1$; the covariance $\tau^2 I$ has
all eigenvalues $\tau^2$; both are frozen buffers, so (i) holds. The diagonal closed
form of
$\operatorname{KL}(\mathcal{N}(\mu, \sigma^2 I)\,\|\,\mathcal{N}(a v_y, \tau^2 I))$
has displacement term $\tfrac{1}{2\tau^2}\|\mu - a v_y\|^2$, giving (ii). For (iii):
the identity
$h_i^{\top} A^{\top} A\, h_j = \mu_i^{\top}\mu_j = s^2 v_i^{\top} v_j = s^2\delta_{ij}$
holds on the aligned family and only says that $\{h_i\}$ is $A^{\top}A$-orthogonal by
construction; it implies nothing about the spectrum of $A^{\top}A$, and the displayed
example verifies the failure. The KL gradient on $A$ is
$(\mu - a v_y)\,h^{\top}/\tau^2$, an alignment force toward the orthonormal frame that
vanishes only at perfect alignment \emph{on non-degenerate samples} (it also vanishes
pointwise at $h = 0$, or by cancellation across a batch) --- but nothing in that force
resists a vanishing within-class component of $A$. For CEB, the displacement term is
$\tfrac12(m - g)^{\top} S_\theta^{-1}(m - g)$: the metric $S_\theta^{-1}$ has condition
number up to $e^{20}$ under the logvar clamp, and the gradient weight $1/s$ at the
collapse equilibrium is Prop.~\ref{prop:variance}. \hfill\qed
\end{proof}

\begin{remark}[Two senses of tightness]\label{rem:senses}
In the pure variational \emph{value} sense (same encoder, same Gaussian family), CEB's
bound is tighter: $\mathrm{ReX}(\theta^*) \le \mathrm{ReX}_F$
(Prop.~\ref{prop:ordering}), the excess being the anchor gap --- the price of the fixed
anchor. The \emph{object} sense does not reverse this ordering: both bounds are valid
upper bounds on their own noisy codes, and neither bounds the residual of a deployed
deterministic continuous representation, which is infinite under the standing
non-atomic assumption (Prop.~\ref{prop:object}). The structural difference is where the certified-versus-
deployed gap lives --- in CEB it is driven by the prediction term inside the inference
path; in FGIB it is a side-path quantity whose size is governed by the conditioning of
$A$ (Prop.~\ref{prop:conditioning}(iii)) --- not which bound is ``tighter'' or
``valid''. The learnability of the side head $A$ bears on that difference: it is a
second trainable party in the matching, unused at inference
(Remark~\ref{rem:A-channel}), so it can absorb part of the anchor gap without moving
the deployed representation. On the aligned family, where $A$ is fixed by the ansatz,
its optimization drives the gap to the equilibrium value $\tfrac12(s^*-a)^2$ at fixed
classifier scale (Prop.~\ref{prop:sweep}); off the ansatz this is an open question
(Section~\ref{sec:limits}).
\end{remark}

\begin{remark}[Unconstrained side head (the implementation)]\label{rem:orth}
Proposition~\ref{prop:inverse} holds for any invertible linear head: the KL invariance
is exact, and the object claim needs only that the perturbation covariance
$\tau^2(A^{\top}A)^{-1}$ be \emph{known} ($A$ is a model parameter) --- but the
object is informative about $h$ only insofar as that covariance is moderate, which the
construction does not guarantee (Prop.~\ref{prop:conditioning}(iii)). The diagonal
variational family is not closed under $A^{-1}$ (Remark~\ref{rem:closure}), so within
the diagonal parameterization of the codebase the equivalence is exact at the
distribution level and approximate as a variational identity; a full-covariance side
head would make it exact within the family. No orthogonality constraint on $A$ is
required --- the orthogonality of the method lives in the fixed anchor frame and in the
equilibrium it induces (the aligned family).
\end{remark}

\begin{remark}[Anchor-frame reading]\label{rem:frame}
With the orthonormal anchor matrix $V$ (columns $v_y$, $K \le z$), left multiplication
by $V^{\top}$ expresses the anchor-space components of $z$ in class coordinates: the
anchors become $a\,e_y$ and the gap of Lemma~\ref{lem:gap} decomposes into $K$
per-class terms plus, when $K < z$, an orthogonal-complement term that the reading
ignores ($V^{\top}$ is a projection, not a coordinate change). The orthogonal frame is
the geometric content of the prior, and it is preserved by the invertible linear
pull-back $A^{-1}$ of Proposition~\ref{prop:inverse} --- not by $A^{\top}$, which is
the Jacobian factor that appears in gradients and equals $A^{-1}$ only when $A$ is
orthogonal. An earlier version conflated the two.
\end{remark}

\begin{remark}[Positioning against VIB and NIB]\label{rem:position}
VIB's rate bounds the \emph{full} mutual information $I(X;Z)$ with a marginal prior;
CEB removes the retained label information ($\mathrm{ReX} = R - I(Y;Z)$ at the optima),
and FGIB inherits the residual-level object while replacing the learnable backward
encoder with a fixed geometric prior (Prop.~\ref{prop:ordering}).
NIB's non-parametric bound $\widehat I \ge I(X;M)$ is $y$-unconditioned and tight only
when the components share a covariance and form well-separated clusters (Kolchinsky et
al.\ 2017, Eq.~10 and the discussion following Eq.~11); when the noise is small relative
to the cluster spread it saturates at $\ln N$ and its gradient vanishes (their
Section~IVA) --- loose exactly in the regime where the parametric FGIB bound is
well-behaved. The two bounds are complementary: NIB certifies the full compression
non-parametrically when the geometry is favorable; FGIB bounds the residual of a
\emph{noisy copy} of the deployed representation parametrically, with a geometry of its
own choosing --- subject to the conditioning caveat of
Proposition~\ref{prop:conditioning}(iii).
\end{remark}

\begin{remark}[Degenerate anchors]\label{rem:a0}
At $a = 0$ the anchors collapse to $N(0, \tau^2 I)$ for all classes and the gap loses its
class
structure: $\Delta = \mathbb{E}_y[\operatorname{KL}(q(z|y)\|N(0,\tau^2 I))]$ --- a
VIB-style
marginal gap; the structural separation of Proposition~\ref{prop:object}(b) still
holds, but the compressed endpoint becomes the trivial collapse
(Remark~\ref{rem:endpoint}).
\end{remark}

% ---------------------------------------------------------------
\subsection{What is proved and what is open}\label{sec:limits}
% ---------------------------------------------------------------

\begin{remark}[Scope of the results]
\textbf{Proved.} (1) The dead-knob proposition for CEB and the MNI reading of it
(Prop.~\ref{prop:dead}, Remark~\ref{rem:mni}); (2) the gap identity, Gaussian moment
matching, and the value ordering of CEB's bound over FGIB's at a fixed encoder
(Lemmas~\ref{lem:gap}, \ref{lem:mm}; Prop.~\ref{prop:ordering}); (3) the blind
certificate on the $T_\alpha$ manifold and the $1{:}\beta$ exchange rate
(Prop.~\ref{prop:flat}); (4) KL invariance under invertible linear maps and the
pointwise equivalence of the side path to a CEB bottleneck at a fixed side head
(Lemmas~\ref{lem:kl-invariance}, \ref{lem:mi-invariance}; Prop.~\ref{prop:fgib-ceb});
(5) the variance decoupling $\sigma^{2*} = \tau^2$ and the non-collapsed compressed endpoint
(Lemma~\ref{lem:lagrangian}, with $\sigma^{2*} = \tau^2$;
Remark~\ref{rem:endpoint}); (6) the one-dimensional sweep
and its strict concavity \emph{at fixed classifier scale}, as a KL--CE surrogate curve
on the aligned family (Props.~\ref{prop:sweep}, \ref{prop:curve};
Remark~\ref{rem:surrogate} --- not a recovered IB curve); (7) the well-conditioning of the fixed
reference side (Prop.~\ref{prop:conditioning}(i)--(ii)).

\textbf{Open or conditional.} (1) The side head $A$ is a trainable, inference-unused
party in the matching (Remark~\ref{rem:A-channel}); combined with the
counterexample of Prop.~\ref{prop:conditioning}(iii), the perturbation in
Prop.~\ref{prop:inverse} can be arbitrarily large, so the bound's object may be far from
the deployed representation exactly where the residual lives. A constraint on $A$ or a
measurement of $\kappa(A^{\top}A)$ over training would close or quantify this channel.
(2) The global sweep requires a bounded classifier scale
(Remark~\ref{rem:c-scale}); the current implementation leaves the scale free and the
proposition conditional. A fixed-temperature cosine classifier would satisfy the
fixed-$c$ assumption by construction; weight decay only mitigates the scale divergence
and replaces the problem with a new joint optimization over $(s,c)$ that must be
re-derived. (3) Neither CEB's nor FGIB's certificate upper-bounds the deployed
deterministic residual (Prop.~\ref{prop:object}); claims of inference-representation
compression are hypotheses, not conclusions. (4) Caveat~2 is not resolved
(Remark~\ref{rem:endpoint}): the compressed endpoint is a chosen label-code geometry,
not certified semantic compression. (5) The two-stage variance dynamics of
Prop.~\ref{prop:variance}(ii)--(iii) await the measurement of
Remark~\ref{rem:p3-status}. (6) The regression RFF anchor analysis inherits the same
fixed-scale and conditioning caveats and is presently a transfer argument, not a proof.
\end{remark}
。
% 记号沿用 Fischer (2020)、Kolchinsky & Tracey (2019, Caveats)、Kolchinsky et al. (2017, NIB)、
% Alemi et al. (2017, VIB)。
%
% 核心结论（FGIB 论文问题陈述与解法）：
% 第一部分（四条失效通道）：
%   P1 退化继承（死旋钮）：CEB ≡ IB(β=1+γ>1)，确定性场景下所有 γ>0 的唯一最优解钉在角点，
%      γ 扫不出任何中间点（SIB Caveat 1 的退化半区，经 β'=1/(1+γ)∈(0,1) 映射）。
%      该角点恰是 CEB 自己的 MNI 目标，rem:mni 正面处理这一反驳；
%   P2 证书盲视 + 1:β 交换率：T_α 流形上 ReX 恒为 0（证书对标签保留量完全盲视），
%      且目标以 1 nat 标签信息 ↔ β nat 认证残差的比率交换；
%   P3 方差竞赛：KL 两侧皆可训练，分类项压 σ→0、钳位处退化为权重 e^{10} 的均值匹配。
%      (ii)(iii) 是受力方向分析而非收敛证明，且可实测（rem:p3-status）；
%   P4 对象错位：两种方法的界都只界住各自的含噪码，都不界住确定性连续推理表示的残差
%      （后者为 +∞）。差别是结构性的（噪声在不在推理路径上），不是紧度排序。
% 第二部分（FGIB 机制，以及它没能关闭的通道）：
%   M1 等价性：旁路 = 固定 A 下作用在 h+已知噪声上的 CEB 瓶颈（逐步成立，非全程固定参考）；
%   M2 扫参：固定分类器尺度 c 时 L_β(s)=CE(s)+(β/2)(s−a)² 严格凸、每 β 唯一内点、曲线严格凹；
%      c 自由时下确界对所有 β 均为 0（rem:c-scale，给出权重归一化分类器的架构修法）；
%   M3 证书：gap 恒等式 + 矩匹配 + 固定编码器下的值排序；固定参考的力持续性成立，
%      但侧头 A 是推理无关的第二个自由方（rem:A-channel），构成未关闭的通道。
%
% 修订记录（第二轮）：理论/实现分层后统一引入锚点方差 τ²（先验 N(a v_y, τ²I)，默认 τ²=1，
% 后验方差最优 σ²*=τ²，σ_p 为公共方差记号；point-anchor 极限统一写为 τ→0）；理论旁路固定为
% 可逆线性头 T(h)=Ah（实现用带 bias 的 nn.Linear 近似，见理论-代码分层 caveat）；
% 明确 sweep 证明的是对齐族上 KL–CE surrogate 曲线（rem:surrogate），不是 IB 互信息曲线；
% 无穷互信息结论统一挂靠分布性 standing assumption；rem:c-scale 的架构修法收紧为
% 固定温度 cosine 分类器（weight decay 只缓解不发散）；a=0 退化点、条件数表述、
% 梯度消失条件、row-space 表述等小项修正。
%
% 修订记录（第一轮）：删除三处错误论断——(1) P3 原称 ReX≈0 时 I(X;Z|Y) 发散，与 gap 恒等式矛盾；
% (2) prop:object(c) 原称 FGIB 的界"更紧/唯一有效"，由 DPI 不成立；(3) prop:conditioning(iii)
% 原从 Gram 恒等式推出条件数为 1，有反例。修正 P2 的换算、Caveat 2 的主张、rem:frame 的
% A^T/V^T、mu_head 的仿射与非推理路径身份。未证结论统一收入 sec:limits。

\newtheorem{lemma}{Lemma}
\newtheorem{proposition}{Proposition}
\newtheorem{remark}{Remark}

% =====================================================================
\section*{Part I. The four failure channels of the conditional bottleneck}
% =====================================================================

\textbf{Notation.}
$X$ is the input and $Y$ the label with $K$ classes; the encoder induces the Markov chain
$Y - X - Z$. The \emph{deterministic scenario} is $Y = f(X)$ ($H(Y|X) = 0$). Following
Fischer (2020), we write the MNI point $I(X;Y) = I(X;Z) = I(Y;Z)$ (Eq.~1), the
chain-rule identity $I(X;Z|Y) = I(X;Z) - I(Y;Z)$ (Eq.~4), the CEB objective
$\mathrm{CEB} \equiv \min_Z I(X;Z|Y) - \gamma I(Y;Z)$ (Eq.~5), the IB objective
$\mathrm{IB} \equiv \min_Z I(X;Z) - \beta_{\mathrm{IB}} I(Y;Z)$ (Eq.~16), the
variational objective
$\mathrm{VCEB} \equiv \min_{e,b,c} \langle \log e(z|x)\rangle - \langle \log b(z|y)\rangle
- \gamma \langle \log c(y|z)\rangle$ (Eq.~15), and the residual certificate
$\mathrm{ReX} \equiv \langle \log e(z|x)\rangle - \langle \log b(z|y)\rangle
\ge I(X;Z|Y)$ (Eq.~20), where $b(z|y)$ is a \emph{trainable} backward encoder (we write
$q(z|x)$ for Fischer's encoder $e(z|x)$; the two notations are interchangeable
throughout). In the loss convention $\mathrm{CE} + \beta\,\mathrm{KL}$ used throughout, the weights are
$\beta = 1/\gamma$ and $\beta_{\mathrm{IB}} = 1 + \gamma = 1 + 1/\beta$. Following
Kolchinsky \& Tracey (2019), we write the IB curve $F(r)$ (their Eq.~1), the IB
Lagrangian $L^{\mathrm{IB}}_\beta = I(Y;T) - \beta I(X;T)$ (their Eq.~2), the
DPI-saturating manifold $T_\alpha := B_\alpha \cdot Y$ (their Eq.~6), the bound
$L^{\mathrm{IB}}_\beta \le (1-\beta)H(Y)$ (their Eq.~7), and their footnote~4 (the
$\beta > 1$ regime of trivial maximizers). The variational objects are
\begin{align*}
R &= \mathbb{E}\big[\operatorname{KL}(q(z|x)\,\|\,m(z))\big] \ge I(X;Z)
&&\text{(VIB rate, marginal prior $m$; Fischer Eq.~19, Alemi et al.\ Eq.~14)} \\
\mathrm{ReX} &= \mathbb{E}\big[\operatorname{KL}(q(z|x)\,\|\,b_\theta(z|y))\big]
\ge I(X;Z|Y)
&&\text{(CEB residual; Fischer Eq.~20)} \\
\widehat I &= -\frac{1}{N}\sum_i \ln\Big[\frac{1}{N}\sum_j
e^{-\|f_i-f_j\|^2/(2\sigma^2)}\Big] \ge I(X;M)
&&\text{(NIB, non-parametric; Eq.~16, the homoscedastic case of Eq.~10)} \\
\mathrm{ReX}_F &= \mathbb{E}\big[\operatorname{KL}(q(z|x)\,\|\,r(z|y))\big]
\ge I(X;Z|Y)
&&\text{(FGIB, fixed anchor)}
\end{align*}
(Fischer 2020, Eqs.~18--20; Alemi et al.\ 2017, Eq.~14; Kolchinsky et al.\ 2017,
Eqs.~10 and 16). Fischer's comparison of the first two: at the respective prior optima
($m = q(z)$, $b = q(z|y)$), $\mathrm{ReX} = R - I(Y;Z)$ --- the conditional prior removes
exactly the retained label information, so CEB's bound is tighter than VIB's by
$I(Y;Z)$. We ask the same question one level down: where does FGIB's fixed-anchor bound
sit relative to CEB's learnable-backward bound?

\begin{proposition}[Channel P1, dead knob: every $\gamma > 0$ selects the same corner]\label{prop:dead}
Assume $Y = f(X)$. For every $\gamma > 0$,
\begin{equation}
\min_Z L_\gamma(Z) = -\gamma H(Y),
\qquad
L_\gamma(Z) := I(X;Z|Y) - \gamma I(Y;Z),
\label{eq:min}
\end{equation}
attained exactly on the corner family
$\{Z : I(Y;Z) = H(Y),\; I(X;Z|Y) = 0\}$ (e.g.\ $Z = Y$, or $Z = (Y, N)$ with
$N \perp X \mid Y$): the corner is unique in the information plane, while at the
representation level it is an entire family. Consequently the map
$\gamma \mapsto \big(I(X;Z^*|Y),\, I(Y;Z^*)\big)$ is constant: no interior point of the
IB curve is ever a Lagrangian minimizer, and CEB inherits the deterministic degeneration
of the Caveats paper (their Issue~1) over its \emph{entire} parameter range, since
$\beta_{\mathrm{IB}} = 1 + \gamma > 1$ always.
\end{proposition}

\begin{proof}
By Fischer's Eq.~4, $L_\gamma(Z) = I(X;Z) - (1+\gamma)I(Y;Z)$. The DPI
($I(Y;Z) \le I(X;Z)$, by the Markov chain $Y - X - Z$) gives
\begin{equation}
L_\gamma(Z) \ge I(Y;Z) - (1+\gamma)\,I(Y;Z)
= -\gamma\,I(Y;Z) \ge -\gamma H(Y),
\label{eq:chain}
\end{equation}
with equality iff both inequalities bind: $I(X;Z) = I(Y;Z)$ and $I(Y;Z) = H(Y)$.
$Z = Y$ attains the bound. Fischer (2020, p.~7) states the underlying equivalence
explicitly: CEB and IB are the same objective for $\gamma = \beta_{\mathrm{IB}} - 1$.
In the Caveats paper's maximizing convention, minimizing $L_\gamma$ is maximizing
$L^{\mathrm{IB}}$ with $\beta' = 1/\beta_{\mathrm{IB}} = 1/(1+\gamma) \in (0,1)$ ---
precisely the range in which their Eq.~7 pins every maximizer at the copy corner (their
Issue~1); their footnote~4 covers the complementary regime $\beta' > 1$ (i.e.\
$\gamma < 0$, outside CEB's range), whose maximizers are the trivial collapse. \hfill\qed
\end{proof}

\begin{remark}[The corner is the MNI point: why that does not rescue the knob]\label{rem:mni}
In the deterministic scenario the corner of Proposition~\ref{prop:dead} \emph{is} the MNI
point: $I(Y;Z) = H(Y) = I(X;Y)$ together with $I(X;Z|Y) = 0$ gives
$I(X;Y) = I(X;Z) = I(Y;Z)$, which is Fischer's Eq.~1. CEB's declared target and the
degenerate attractor therefore coincide, and Kolchinsky \& Tracey concede the point
themselves (their Section~4): ``if $Y = f(X)$ and one's goal is precisely to find the corner
point $\langle H(Y), H(Y)\rangle$ (maximum compression at no prediction loss), then our
results suggest that the IB Lagrangian provides a robust way to do so, being invariant to
the particular choice of $\beta$.'' Read that way, Proposition~\ref{prop:dead} states
hyperparameter robustness rather than failure. Three things keep it from closing the matter.
\begin{enumerate}
\item[(i)] \emph{The corner is one point in the information plane and an entire equivalence
class of representations.} The objective is exactly constant on
$\{Z : I(Y;Z) = H(Y),\ I(X;Z|Y) = 0\}$, which contains $Z = Y$, $Z = (Y,N)$ with
$N \perp X \mid Y$, and every bijective recoding of these. Adversarial robustness, OoD
detection and calibration --- the properties CEB is argued for --- are properties of the
\emph{geometry} of $Z$, on which the objective has no opinion; which member of the class is
reached is settled by the optimizer. That is channels P3 and P4 below.
\item[(ii)] \emph{CEB's own reading of the knob is unavailable at the exact optimum.}
Fischer (2020, p.~7) writes ``$\rho < 0$ may capture less than the MNI. $\rho > 0$ may
capture more'', whereas Proposition~\ref{prop:dead} places every $\gamma > 0$ at the MNI
point exactly. His argument for preferring $\gamma = 1$ is correspondingly not about
minimizers but about the optimizer: weighting the two appearances of $I(Y;Z)$ differently
``forces the optimization procedure to count bits of $I(Y;Z)$ in two different ways'' (his
Eq.~8 and surrounding discussion).
\item[(iii)] \emph{The knob does move trained models, which locates what it acts on.} On
Fashion MNIST --- a deterministically labelled dataset --- Fischer's own Figure~4 reports
the rate lower bound $R_X$ rising from $\approx 2.0$ to $\approx 5.0$ nats as $\rho$ goes
from $-1$ to $5$, with $\rho = 0$ landing on $I(X;Y) = \log 10 \approx 2.3$ nats. Since
every $\gamma > 0$ shares one exact optimum, that axis is not a path along the IB curve: it
is the distance between the optimum of the variational objective over a restricted family
and the point at which training was stopped. A knob acting through the approximation error
is precisely the knob that Propositions~\ref{prop:flat} and~\ref{prop:variance} show cannot
be read off the certificate.
\end{enumerate}
Kolchinsky \& Tracey state that their results ``are based on analytically-provable
properties of the IB curve, i.e., global optima'' and ``do not concern practical issues of
optimization''; Proposition~\ref{prop:flat} supplies the finite-$\beta$ variational
counterpart. The account is owed by any method reporting a $\beta$-sweep on
deterministically labelled data --- FGIB and the experiments of this paper included
(Remark~\ref{rem:c-scale}).
\end{remark}

\begin{proposition}[Channel P2, a blind certificate and a $1{:}\beta$ exchange rate]\label{prop:flat}
On the manifold $T_\alpha = B_\alpha \cdot Y$ of the Caveats paper (their Eq.~6),
$I(X;T_\alpha|Y) = 0$ for every $\alpha$, and with the optimal backward encoder and
classifier,
\begin{equation}
\mathrm{ReX}(T_\alpha) = 0,
\qquad
\mathrm{CE}(T_\alpha) = H(Y) - I(Y;T_\alpha) =: \delta_\alpha,
\qquad
\mathrm{VCEB}(T_\alpha) = \gamma\,\delta_\alpha .
\label{eq:vceb}
\end{equation}
Two consequences.
\emph{(a) The certificate is blind to label retention.} $\mathrm{ReX} = 0$ holds from the
total collapse ($\alpha = 0$, no label information) to the full copy ($\alpha = 1$), so its
tightness is compatible with retaining all, part, or none of $Y$. Fischer's reading of the
certificate as a meter --- ``observing ReX tells us exactly how many bits we are from the
MNI point'' (p.~7) --- therefore holds only along the $I(X;Z|Y)$ axis; on the necessity axis
the meter reads $0$ at chance level.
\emph{(b) Exchange rate.} With the optimal classifier the objective is
$L = \mathrm{CE} + \beta\,\mathrm{ReX} = \delta + \beta\,\mathrm{ReX}$, so it is indifferent
between giving up $\delta$ nats of label information and saving $\delta/\beta$ nats of
certified residual. In the heavy-regularization regime the bottleneck is supposed to act in,
the whole label axis --- up to $H(Y)$ nats --- is purchasable for $H(Y)/\beta$ nats of
certified compression.
\end{proposition}

\begin{proof}
$T_\alpha$ is a function of $Y$ alone, so $T_\alpha \perp X \mid Y$ and
$I(X;T_\alpha|Y) = 0$; the Caveats paper shows $I(Y;T_\alpha)$ sweeps $[0, H(Y)]$.
With the true backward encoder ($b = p(z|y)$) the certificate is tight:
$\mathrm{ReX} = I(X;T_\alpha|Y) = 0$ (Fischer Eqs.~9--11, 20). The optimal classifier
achieves $\langle \log c(y|z)\rangle = -H(Y|T_\alpha)$ (Fischer Eq.~12), i.e.\
$\mathrm{CE} = H(Y|T_\alpha) = \delta_\alpha$, and the VCEB value follows from Eq.~15.
Part~(b) is the definition of the objective evaluated at $\mathrm{CE} = \delta$.
\hfill\qed
\end{proof}

\begin{remark}[What the flatness claim is, and what it is not]\label{rem:convention}
An earlier form of this argument read the near-optimality band off the \emph{value}
$\mathrm{VCEB} = \gamma\delta \to 0$. That reading is a scaling artifact:
$\mathrm{VCEB} = \gamma\,(\mathrm{CE} + \beta\,\mathrm{ReX}) = \gamma L$, so the two
conventions differ by a positive factor and share their minimizers, and on $T_\alpha$ itself
the objective is $L = \delta$ in either convention and strictly prefers $\alpha = 1$. What
is convention-free is the \emph{ratio} of Proposition~\ref{prop:flat}(b): the prediction
term is weaker than the compression term by exactly the factor $\beta$, so in units of the
compression term the entire manifold lies within $H(Y)/\beta$ of optimal. Under a
scale-invariant optimizer (Adam, used throughout the experiments) only the ratio matters;
under plain SGD at fixed learning rate the convention additionally rescales the effective
step size. Combined with Proposition~\ref{prop:flat}(a) --- no certified compression is
bought anywhere on the manifold --- the operative statement is that the label axis is cheap
and unmetered, not that the objective is literally flat along it.
Pushing $\beta$ up also drives $\beta_{\mathrm{IB}} = 1 + 1/\beta \to 1^{+}$: the objective
approaches the Caveats paper's $\beta = 1$ degeneracy (their Eq.~7 discussion), at which
\emph{every} point of the diagonal is simultaneously optimal, from above. Fischer (2020,
p.~7) anticipates the direction --- ``$\rho < 0$ may capture less than the MNI'' --- but by
Proposition~\ref{prop:dead} that cannot happen at the exact optimum;
Proposition~\ref{prop:flat} gives the rate at which the variational objective is willing to
let it happen.
\end{remark}

\begin{remark}[Approximately deterministic labels]\label{rem:approx}
The Caveats paper (Appendix~C, Issue~1) shows that when $Y$ is only $\epsilon$-close to
deterministic, all optimizers fall within $O(-\epsilon\log\epsilon)$ of a single corner:
the dead-knob statement of Proposition~\ref{prop:dead} degrades gracefully, and the
near-optimality band of Proposition~\ref{prop:flat} is perturbed by the same order.
\end{remark}

\begin{proposition}[Channel P3, collusion collapse: the variance race and the overfit channel]\label{prop:variance}
For the diagonal Gaussian family $e(z|x) = \mathcal{N}(m(x), \sigma^2 I)$,
$b(z|y) = \mathcal{N}(g(y), \tau^2 I)$ (the implementation's setup, with the
\texttt{logvar} clamp $[-10, 10]$ of \texttt{kl\_divergence}), and a \emph{linear}
classifier $c(z) = Wz + b$ (the implementation's classifier is \texttt{nn.Linear}):
\begin{enumerate}
\item[(i)] for fixed means, $\operatorname{KL}$ is minimized over $\sigma^2$ at
$\sigma^2 = \tau^2$, with value $\tfrac{1}{2\tau^2}\|m - g\|^2$ --- a pure mean-matching
term;
\item[(ii)] the classifier term favors $\sigma^2 \to 0$:
$\mathrm{CE}(c(z), y) = -\log \mathrm{softmax}_y(Wz+b)$ is convex in $z$ (log-sum-exp
of an affine function of $z$, minus a linear term --- the convexity claim requires $c$
affine), so Jensen gives
$\mathbb{E}_z[\mathrm{CE}(c(z), y)] \ge \mathrm{CE}(c(m), y)$. The KL's derivative in
$\sigma^2$ vanishes at $\sigma^2 = \tau^2$, but along the matched line it \emph{resists}
shrinking $\tau^2$ while the means are unmatched
($\partial \operatorname{KL}/\partial \tau^2|_{\sigma^2=\tau^2}
= -\|m-g\|^2/(2\tau^4) < 0$). The joint dynamics therefore runs in two stages: the mean
matching $m \to g$ completes first (with weight $1/\tau^2$, which grows as $\tau$
shrinks), and once $\|m-g\| \to 0$ the variance race is unopposed --- the classifier
drives $\sigma^2 \to 0$ with $\tau^2$ following the moment-matched value --- until the
clamp at $s = e^{-10}$, where the encoder collapses toward the deterministic code
$z \approx g(y)$ and
$\mathrm{ReX} = \tfrac{1}{2s}\mathbb{E}\|m(x) - g(y)\|^2$: a mean-matching force of
weight $1/s = e^{10} \approx 2.2\times10^{4}$ between two \emph{trainable} maps;
\item[(iii)] the certified object and the deployed object separate. By the gap lemma
(Lemma~\ref{lem:gap}) $\mathrm{ReX} = I(X;Z|Y) + \Delta \ge I(X;Z|Y)$, so a small
$\mathrm{ReX}$ \emph{entails} a small residual for the stochastic training code --- the
certificate cannot under-report its own object. What it does not control is the relation
between the training code and the deployed representation. Exact mean matching
($m(x) = g(y)$) makes both objects class-constant, so both residuals vanish; the
scenario of interest is the \emph{near}-match regime, where
$\mathbb{E}\,\|m(X) - g(Y)\|^2 = o(s)$ keeps $\mathrm{ReX}$ small while the deployed
continuous $m$ retains within-class variation, and any non-class-constant continuous
$m$ has $I(X;m(X)|Y) = +\infty$ under the standing non-atomic assumption (the
pathology of deterministic networks, Amjad \&
Geiger 2018; Kolchinsky et al.\ 2017). The variance race therefore moves the
\emph{certified} object toward a noise-dominated class-constant code while the
\emph{deployed} object can keep diverging residual --- and the encoder is fitted to the
backward encoder's per-class targets, memorizing the label code $g(y)$ rather than
being told what residual to erase, while the certificate reports $\mathrm{ReX} \approx
0$ on a code that inference does not use.
\end{enumerate}
\end{proposition}

\begin{proof}
The closed form
$\operatorname{KL} = \tfrac{d}{2}\big[\sigma^2/\tau^2 + \|m-g\|^2/(d\tau^2)
- 1 + \ln(\tau^2/\sigma^2)\big]$ has derivative $\tfrac{d}{2}(1/\tau^2 - 1/\sigma^2)$
in $\sigma^2$, so the variance optimum is $\sigma^2 = \tau^2$, giving (i); at
$\sigma^2 = \tau^2 = s$ the value is $\|m-g\|^2/(2s)$. For (ii): $\mathrm{CE}(c(z), y)$ has Hessian
$\operatorname{diag}(p) - pp^{\top} \succeq 0$ (the categorical Fisher information), so
it is convex in $z$ and Jensen gives the stated inequality; the same Hessian identity makes
$\sigma^2 \mapsto \mathbb{E}_z[\mathrm{CE}(c(z),y)]$ non-decreasing (Gaussians of larger
variance dominate in convex order), strictly increasing in $\sigma^2$ wherever the
classifier is unsaturated, so that term pushes $\sigma^2$ down. The KL's partial derivatives at the
matched point are $\partial_{\sigma^2}\operatorname{KL} = 0$ and
$\partial_{\tau^2}\operatorname{KL} = -\|m-g\|^2/(2\tau^4) < 0$, so the KL pulls
$\tau^2$ up until the means match, while the mean-matching gradient on the encoder is
$(m - g)/\tau^2$; once $\|m-g\| \to 0$ the race to the clamp is unopposed. For (iii): the
gap identity (Lemma~\ref{lem:gap}) forces the certified residual and the certificate to
vanish together; the divergence attaches to the deployed mean $m(x)$ whenever it is not
class-constant (Amjad \& Geiger; Kolchinsky et al.). \hfill\qed
\end{proof}

\begin{remark}[Status of the two-stage account]\label{rem:p3-status}
Parts~(ii) and~(iii) are a stationarity-and-direction-of-force argument, not a convergence
proof: the ``two stages'' presume a time-scale separation (mean matching completes before
the variance race) that is not established here. The account is directly falsifiable in
this codebase: log the mean posterior log-variance and the fraction of units at the clamp
$s = e^{-10}$ during CEB training. \emph{We ran the test.} On MNIST, across all six
$\beta$ settings and 5 seeds (Section~\ref{sec:ablation} of the main paper), the trained
posterior and prior log-variances of CEB remain near their initialization ($\sigma^2 =
\tau^2 = 1$, mean $\log\sigma^2$ within $0.74$ of zero, clamp fraction $0.0$ at every
epoch of every run). At the early-stopping time scale of the experimental protocol, the
variance race therefore does not materialize, and
Proposition~\ref{prop:variance}(iii) is empirically void for those runs: Part~(iii)
stands as a sharp hypothesis about what longer training or different schedules would
produce, the force analysis of (i)--(ii) as the only proved content, and the
measurement itself as the finding.
\end{remark}

\begin{remark}[Train--evaluation mismatch]\label{rem:mismatch}
The classifier is trained on the sampled code $z$ and evaluated on the mean $m(x)$
(Fischer 2020, Appendix~A.1). Proposition~\ref{prop:variance}(ii) is the force that
closes this gap by destroying the noise altogether ($\sigma \to 0$), at which point the
bottleneck's stochasticity --- the object the bound is defined on --- disappears.
\end{remark}

\begin{proposition}[Channel P4, certificate-object mismatch: neither bound is an upper bound on the deployed representation]\label{prop:object}
(a)~\emph{CEB}. The classifier consumes the sampled code $Z$ at training and the mean
$\mu(X)$ at evaluation (Fischer 2020, Appendix~A.1). At the variance-optimal posterior,
$\sigma(x) \equiv \sqrt{S_\theta(y)}$ (the backward encoder's per-class standard
deviation) is class-constant, so $Z = \mu(X) + \sqrt{S_\theta(Y)}\,\varepsilon$ is a
noisy function of $\mu(X)$ given $Y$; by the DPI,
\begin{equation}
I(X; Z | Y) \;\le\; I(X;\, \mu(X) | Y) ,
\label{eq:dpi}
\end{equation}
strictly so whenever the noise destroys within-class detail. Hence CEB's certificate
$\mathrm{ReX} \ge I(X;Z|Y)$ is a bound on the \emph{training code}, not on the residual
of the representation deployed at inference. The gap is not a small correction: whenever
$\mu$ is a continuous map that is not class-constant, $I(X;\mu(X)|Y) = +\infty$ under
the standing non-atomic assumption (Amjad \& Geiger 2020; Kolchinsky et al.\ 2017), so
no finite quantity --- CEB's certificate included --- bounds the deployed residual from
above. The CE term presses the noise
toward $0$ (by Jensen, $\mathbb{E}_z[\mathrm{CE}(c(z),y)] \ge \mathrm{CE}(c(\mu),y)$),
but no finite $\sigma$ closes the gap.
(b)~\emph{FGIB}. The inference path is $h$ at both training and evaluation; the side
noise of Proposition~\ref{prop:inverse} is not part of any inference path. By the DPI
applied to the same perturbed-variable construction,
\begin{equation}
I(X;\, h + \tau (A^{\top}A)^{-1/2}\varepsilon \,|\, Y)
\;\le\; I(X; H | Y) ,
\label{eq:dpi-f}
\end{equation}
so the direction of the bound is the same as CEB's: $\mathrm{ReX}_F$ upper-bounds the
noisy side-path object, and \emph{no} finite value upper-bounds the deployed residual.
The point-anchor limit $\tau \to 0$ does not repair this: it is a change of model whose
bound value $\|\mu - a v_y\|^2/(2\tau^2)$ typically diverges
(Remark~\ref{rem:limit}), and even the limiting object needs $A$ well-conditioned
(Proposition~\ref{prop:conditioning}(iii)).
(c)~\emph{Conclusion}. In the pure variational value sense (same encoder, same Gaussian
family), CEB's bound is tighter (Proposition~\ref{prop:ordering}). In the object sense
the two are not ordered: both are valid upper bounds on their respective noisy codes, and
neither bounds the deterministic deployed representation. What remains structural --- not
a tightness claim --- is \emph{where the gap lives}: CEB's perturbation is inside the
inference path and is driven by the prediction term (so removing it changes the model),
while FGIB's perturbation lives on a side path that inference never reads and is
controlled by side-path parameters alone. The latter is a cleaner separation of
regularizer from model, not a tighter bound.
\end{proposition}

\begin{remark}[The bound is a training instrument, not an inference one]\label{rem:training-role}
No method computes its bound at inference: the bound's roles are \emph{training-time}
--- a regularizer (a differentiable surrogate for the intractable information term) and
a meter (Fischer 2020, p.~7: ``observing ReX tells us exactly how many bits we are from
the MNI point''). In the $Z \perp X \mid Y$ limit CEB's mean is also class-constant
($m(x) = \mathbb{E}[Z|x] = \mathbb{E}[Z|y]$), so the mismatch of
Proposition~\ref{prop:object}(a) is not a limit failure: it is a finite-$\beta$
\emph{under-reporting} --- the regularizer controls, and the meter reports, the training
code, which is strictly more compressed than the deployed representation $\mu(X)$, by a
positive, trained deficit that cannot be measured from the certificate. FGIB's
side-path noise is a known side-path parameter and is removable in principle
($\tau \to 0$ changes the side path, never the deployed model), but for finite
$\tau$ its regularizer and meter act on a \emph{perturbed copy} of the deployed
object, with the perturbation size governed by the conditioning of $A$
(Proposition~\ref{prop:conditioning}(iii)).
The one genuine inference-time use of a certificate value is as a detection score
(Fischer 2020, RG3, uses the marginal rate $R$): CEB's label-conditioned
$\mathrm{ReX}$ is not even computable at inference (the label is unknown --- Fischer
trains a separate marginal for this purpose), while FGIB's geometry gives a label-free
score, the distance to the nearest anchor
$\min_y \operatorname{KL}(q(z|x)\|\mathcal{N}(a v_y, \tau^2 I))$, with no extra network
(the minimum is an internal enumeration over the $K$ fixed anchor buffers --- the single
QR orthonormal frame --- and consults no ground-truth label, no more than softmax
consults labels when it enumerates the $K$ logits). Whether this score is useful for OoD
detection is a hypothesis; nothing here proves it.
\end{remark}

\textbf{The problem addressed by FGIB.}
The four channels are addressed by the FGIB design as follows; two are closed by
construction, one holds conditionally, and one is only partially addressed. The precise
scope of each claim is documented in Part~II and the open items are collected in
Section~\ref{sec:limits}.
\begin{itemize}
\item \emph{Dead knob and cheap label axis} (Props.~\ref{prop:dead}, \ref{prop:flat}):
FGIB's compression term is the anchor gap $\tfrac{1}{2\tau^2}(s-a)^2$, quadratic in the
alignment scalar, and \emph{at fixed classifier scale} the Lagrangian
$\mathrm{CE}(s) + \tfrac{\beta}{2\tau^2}(s-a)^2$ is strictly convex with a unique interior
minimizer per $\beta$ --- the $\beta$-sweep is one-to-one onto the whole trade-off curve
(Props.~\ref{prop:sweep}, \ref{prop:curve}). With the scale free the infimum is $0$ for
every $\beta$; Remark~\ref{rem:c-scale} states the condition and the architectural fix.
\item \emph{Blind certificate} (Prop.~\ref{prop:flat}): with the fixed anchor
$r(z|y) = \mathcal{N}(a v_y, \tau^2 I)$ the gap
$\mathbb{E}_y[\operatorname{KL}(q(z|y)\,\|\,r(z|y))]$ can shrink only if the
\emph{posterior} moves toward the anchors --- the reference itself cannot chase the
encoder, so the force is persistent (Prop.~\ref{prop:ordering} and
Remark~\ref{rem:force}). The side head $A$, however, is a second trainable party in that
matching, and it is not part of the inference path
(Remark~\ref{rem:A-channel}).
\item \emph{Variance collapse} (Prop.~\ref{prop:variance}): the classifier consumes the
deterministic $h$ at both training and evaluation, so $\sigma^{2*} = \tau^2$ for every $\beta$
(the variance decouples from the prediction term; Lemma~\ref{lem:lagrangian}), the
side path keeps a finite covariance and the bound stays well-defined
(Remark~\ref{rem:limit}), and the trivial collapse $\mu \equiv 0$ pays
the fixed positive price $\operatorname{KL} = \tfrac12 a^2 > 0$ for every class --- the
degenerate corner is excluded by geometry rather than by the optimizer's good behavior.
This channel is closed cleanly, by construction.
\item \emph{Certificate-object mismatch} (Prop.~\ref{prop:object}): the classifier
consumes the deterministic $h$ at both training and evaluation --- there is no
training-code/evaluation-mean split --- but the bound is an upper bound on a
\emph{noisy side-path copy} of $h$, not on $I(X;H|Y)$, which no finite certificate
bounds. This channel is \emph{not} closed; what FGIB buys is that the perturbation is
side-path-only, and its size is governed by the conditioning of $A$
(Prop.~\ref{prop:conditioning}(iii)).
\end{itemize}

% =====================================================================
\section*{Part II. The FGIB theory}
% =====================================================================

\textbf{FGIB notation.}
The deterministic main path computes $h = f_{\mathrm{enc}}(X)$ and the classifier
consumes $h$ directly; the side path parameterizes $q(z|x) = \mathcal{N}(\mu(x),
\Sigma(x))$ with \emph{linear} mean head $\mu(x) = A\,h(x)$ for
$A \in \mathbb{R}^{z \times z}$ (the theory assumes $z = d$ and $A$ invertible; the
inverse transforms of Propositions~\ref{prop:fgib-ceb} and~\ref{prop:inverse} are
defined under this assumption), and a separate logvar head for $\Sigma(x)$. The
side path is \emph{not part of inference}: at evaluation the model returns the classifier
output and never evaluates the mean head (so $A$ is a trainable module that exists only
inside the regularizer --- structurally the same position CEB's backward encoder
occupies on the other side of the KL). The label-conditional prior is the
\emph{fixed} anchor
$r(z|y) = \mathcal{N}(a\,v_y, \tau^2 I)$ with \emph{anchor variance} $\tau^2 > 0$
(default $1.0$): the $v_y$ are orthonormal directions obtained by
QR-decomposing a random matrix ($K \le z$; the orthonormal factor of the random matrix),
scaled by the anchor scale $a$ (default $4.0$), stored as
non-trainable buffers. At the variance optimum of the KL the posterior variance matches
the anchor variance, $\sigma^{2*} = \tau^2$; we write $\sigma_p^2$ for this common
variance where the distinction does not matter. The orthogonality of the method lies in
this fixed frame and in the equilibrium it induces --- the KL drives the
class-conditional means $\mu_y = A\,\mathbb{E}[h|y]$ toward the orthonormal anchors
$a v_y$, so the transformed class representations become mutually orthogonal (the
aligned family below) --- while $A$ itself remains an unconstrained linear map. For
regression, fixed random-Fourier anchors
$\mu_p(y) = a\sqrt{2/z}\,\cos(2\pi\omega y + b)$ with frozen
$\omega \sim \mathcal{N}(0, 2^2)$, $b \sim U[0, 2\pi)$ replace the per-class table
($\|\mu_p(y)\| \approx a$). The loss is
$L_{\mathrm{FGIB}} = \mathbb{E}[\mathrm{CE}(c(h), y)]
+ \beta\,\mathbb{E}[\operatorname{KL}(q(z|x)\,\|\,r(z|y))]$, with gradients reaching the
main path only through the shared encoder; the zero-initialized logvar head fixes
$\sigma^2 = 1$ at initialization, which equals $\tau^2$ in the default configuration
$\tau^2 = 1$ (for general $\tau$ the theory's variance optimum $\sigma^2 = \tau^2$ is
reached by the logvar head, whose initialization is a scheduling choice), so training
starts at $\operatorname{KL} \approx \tfrac{1}{2\tau^2} a^2$.

\emph{Implementation caveat (theory--code layering).} The theory works with a fixed
orthonormal factor, an invertible linear side head $T(h) = Ah$, and per-class Gaussian
priors with anchor variance $\tau^2$. The code implements an \emph{affine} head
(\texttt{nn.Linear} with bias) --- which the theory absorbs as an idealization, the bias
being immaterial when $A$ is learned jointly --- and does not enforce orthogonality,
invertibility, or conditioning constraints on $A$. The guarantees below (orthonormal
frames, invariance, conditioning) are properties of the idealized model; the
implementation approximates them, and the guarantees do not transfer automatically
(Remarks~\ref{rem:A-channel}, \ref{rem:rank}).
\textbf{Distributional standing assumption.} Several results state that a deterministic
continuous representation carries infinite conditional information. Those statements
assume the standard continuous setting: $X \mid Y$ is non-atomic, and the induced
representation $H(X) \mid Y$ is a non-atomic continuous variable. For the finite
empirical distribution of the training data no such conclusion is claimed; there the
quantities are finite and only their growth with sample size is governed by these
results.

% ---------------------------------------------------------------
\subsection{The equivalence: the side path is a CEB bottleneck on the main path}
% ---------------------------------------------------------------

\textbf{Motivation.}
Stochastic bottlenecks such as VIB and CEB are \emph{stochastic at training time but
deterministic at evaluation time}: the classifier is trained on sampled codes
$z \sim q(z|x)$ yet evaluated on the mean $\mu_q(x)$ (Fischer 2020, Appendix A.1), so
its training and evaluation distributions differ. In the same models the cross-entropy
gradient and the KL gradient act on the same heads: the classification term can exploit
the variance (by Jensen, $\mathbb{E}_z[\mathrm{CE}(c(z), y)] \ge \mathrm{CE}(c(\mu_q),
y)$, so loosening the noise relaxes the upper bound), while the KL term pushes the
posterior back toward the prior. Finally, the regularizer constrains the noisy code $z$,
whereas inference consumes the noise-free mean --- the object being regularized is not
the object being used. FGIB resolves all three mismatches by design: the classifier
consumes the deterministic representation $h$ at both training and evaluation time, the
two objectives act on disjoint parameter sets (the classification gradient does not
enter the side heads, and the KL gradient reaches the main path only through the shared
encoder), and, by Proposition~\ref{prop:fgib-ceb} below, the side-path regularizer is,
\emph{at each fixed value of the side head}, a CEB bottleneck read in the coordinates in
which the posterior is centered at the representation that inference actually uses
(Remark~\ref{rem:A-channel} states why this is weaker than a fixed constraint on $h$).
Confining the stochasticity to a side path also keeps the variational objective
well-defined: for continuous deterministic representations the residual information
diverges (Amjad \& Geiger; Kolchinsky et al.), so a purely deterministic bottleneck
admits no finite variational bound, while the finite side-path covariance of FGIB does
(Remark~\ref{rem:limit}). Whether the robustness benefits of the stochastic bottleneck
survive the deterministic inference path is the empirical question studied in the
experiments.

\begin{lemma}[Linear invariance of the Gaussian KL]\label{lem:kl-invariance}
Let $q = \mathcal{N}(\mu_q, \Sigma_q)$ and $p = \mathcal{N}(\mu_p, \Sigma_p)$ on
$\mathbb{R}^{z}$, and let $A \in \mathbb{R}^{z \times z}$ be invertible. Define the
common pushforwards $q' = A_\# q = \mathcal{N}(A\mu_q, A\Sigma_q A^\top)$ and
$p' = A_\# p$. Then $\operatorname{KL}(q'\,\|\,p') = \operatorname{KL}(q\,\|\,p)$.
\end{lemma}

\begin{proof}
With $d = z$, the closed form reads
$\operatorname{KL}(q\,\|\,p)
= \tfrac{1}{2}\big[\operatorname{tr}(\Sigma_p^{-1}\Sigma_q)
+ (\mu_q - \mu_p)^\top \Sigma_p^{-1}(\mu_q - \mu_p)
- d + \ln|\Sigma_p| - \ln|\Sigma_q|\big]$.
Under the common map $A$ the three non-constant terms are unchanged:
\begin{align*}
\operatorname{tr}\big((A\Sigma_p A^\top)^{-1}(A\Sigma_q A^\top)\big)
&= \operatorname{tr}(A^{-\top}\Sigma_p^{-1}A^{-1} A\,\Sigma_q A^\top)
 = \operatorname{tr}(A^{-\top}\Sigma_p^{-1}\Sigma_q A^\top) \\
&= \operatorname{tr}(\Sigma_p^{-1}\Sigma_q A^\top A^{-\top})
 = \operatorname{tr}(\Sigma_p^{-1}\Sigma_q), \\[2pt]
(A\Delta\mu)^\top (A\Sigma_p A^\top)^{-1} (A\Delta\mu)
&= \Delta\mu^\top A^\top A^{-\top}\Sigma_p^{-1}A^{-1}A\,\Delta\mu
 = \Delta\mu^\top \Sigma_p^{-1}\Delta\mu, \\[2pt]
\ln|A\Sigma_p A^\top| - \ln|A\Sigma_q A^\top|
&= \big(\ln|\Sigma_p| + 2\ln|A|\big) - \big(\ln|\Sigma_q| + 2\ln|A|\big)
 = \ln|\Sigma_p| - \ln|\Sigma_q|,
\end{align*}
where $\Delta\mu = \mu_q - \mu_p$ and the trace step uses the cyclic property. \hfill\qed
\end{proof}

\begin{lemma}[Invariance of the CEB objective and the MNI point]\label{lem:mi-invariance}
For any bijective measurable map $\phi$,
$I(X;\phi(Z)|Y) = I(X;Z|Y)$ and $I(Y;\phi(Z)) = I(Y;Z)$.
Consequently the CEB objective (Eq.~5) is invariant under bijections of $Z$, and
$Z$ attains the MNI point (Eq.~1) if and only if $\phi(Z)$ does.
\end{lemma}

\begin{proof}
Mutual information is invariant under bijective transformations of either argument
(data processing inequality applied in both directions). \hfill\qed
\end{proof}

\begin{proposition}[The FGIB side path is a CEB bottleneck on the main path, at each fixed side head]\label{prop:fgib-ceb}
Let $z = d$ and let the learned linear head $T(h) = A h$ be invertible (unconstrained,
so a.s.\ so at initialization). The \emph{fixed}
prior $r(z|y) = \mathcal{N}(\mu_p(y), \Sigma_p(y))$ --- for FGIB, the orthonormal anchor
table $\mu_p(y) = a\,v_y$, $\Sigma_p = \tau^2 I$ --- is, by construction, the image
under $T$ of the canonical prior
$r'(\tilde{z}|y) = \mathcal{N}(A^{-1}\mu_p(y),\, A^{-1}\Sigma_p(y) A^{-\top})$
(since $A A^{-1} = I$), and the posterior $q(z|x) = \mathcal{N}(A\,h(x), \Sigma_q(x))$
is the image of
$q'(\tilde{z}|x) = \mathcal{N}(h(x), A^{-1}\Sigma_q(x) A^{-\top})$.
\emph{At a fixed value of $A$}, the prior is never transformed by training:
$A^{-1}$ only reads the fixed anchor geometry in the coordinates in which the posterior
is centered at $h$; during training $A$ moves, so the pulled-back prior moves with
it (Remark~\ref{rem:A-channel}).
By KL invariance under the common invertible linear map $T^{-1} = A^{-1}$
(Lemma~\ref{lem:kl-invariance}), pointwise for every $(x, y)$,
\begin{equation}
\operatorname{KL}\big(q(z|x)\,\|\,r(z|y)\big)
= \operatorname{KL}\big(q'(\tilde{z}|x)\,\|\,r'(\tilde{z}|y)\big),
\label{eq:fgib-ceb}
\end{equation}
That is, at each fixed side head the side-path regularizer is pointwise identical to a
CEB bottleneck with a backward encoder \emph{frozen at the current head}, applied to
a noisy version of the main-path representation,
$\tilde{z} = h + \tilde{\epsilon}$ with
$\tilde{\epsilon} \sim \mathcal{N}(0, A^{-1}\Sigma_q A^{-\top})$,
using the canonical prior $r'$ as that backward encoder. For FGIB the canonical backward
encoder is the frame $\mathcal{N}(a\,A^{-1} v_y, \tau^2 A^{-1}A^{-\top})$, whose
mean anchors $A^{-1}(a v_y)$ are orthonormal at scale $a$ in the induced metric
$A^{\top}A$
($\langle A^{-1}v_i, A^{-1}v_j\rangle_{A^{\top}A} = \delta_{ij}$): the orthogonality
enters through the fixed frame $\{v_y\}$ and the equilibrium it induces, not through a
constraint on $A$.
Hence the regularizer's gradient with respect to $h$ (through the shared encoder)
coincides, up to the fixed linear map $A^{\top}$, with the gradient that a CEB model
applies to its representation mean.
\end{proposition}

\begin{proof}
The map $z \mapsto \tilde{z} = A^{-1}z$ is a bijection, and
$q' = (A^{-1})_\# q$, $r' = (A^{-1})_\# r$:
$A^{-1}A\,h = h$ and $A^{-1}\Sigma_q A^{-\top}$ give the canonical posterior, and
$A^{-1}\mu_p$, $A^{-1}\Sigma_p A^{-\top}$ give the canonical prior, with
$A\,(A^{-1}\mu_p) = \mu_p$ and $A\,(A^{-1}\Sigma_p A^{-\top})\,A^{\top}
= \Sigma_p$, so $r$ is exactly the image of $r'$ under $T$.
Equation~\eqref{eq:fgib-ceb} is Lemma~\ref{lem:kl-invariance} applied to $A^{-1}$.
Sampling $\tilde{z} \sim q'$ is equivalent to
$\tilde{z} = h + \tilde{\epsilon}$ with the stated noise.
Combined with the variational bound of Fischer (2020, Eqs.~9--11), the regularizer
upper-bounds the residual information of the noisy main-path representation:
$I(X;\tilde{Z}|Y) \le \mathbb{E}[\operatorname{KL}(q'(\tilde{z}|x)\,\|\,r'(\tilde{z}|y))]$.
\hfill\qed
\end{proof}

\begin{remark}[Trainable-implementation caveat: the side head is a second trainable party in the matching]\label{rem:A-channel}
The identity of Proposition~\ref{prop:fgib-ceb} holds pointwise at each value of the
side head, so the asymmetry advertised against CEB --- one trainable side versus two ---
is weaker than it looks. In the idealized theory $A$ may be treated as fixed
(orthonormal or any fixed invertible map); in the \emph{trainable implementation} $A$ is
free and \emph{unused at inference}:
like CEB's backward encoder $b(z|y)$, it lives inside the KL and can lower it without
moving the deployed representation $h$. Since $h$ is linear-separable by the jointly
trained classifier, the linear fit of the anchor target $a v_y$ on $h$ fits well, and
much of that fit can be absorbed by $A$ alone. The variance term
$\tfrac{z}{2}(\sigma^2/\tau^2 - 1 - \ln(\sigma^2/\tau^2))$ does not depend on $A$ at
all, so nothing in
the regularizer resists $A$ drifting toward a projector onto the $K$-dimensional
between-class subspace with a vanishing component on within-class directions. Full
collapse is blocked --- with invertible $A$, $\operatorname{KL} = 0$ forces $h$ itself
to be class-constant --- but the approach to the rank-$K$ regime is not blocked, and
that regime is exactly where the conditioning failure and the perturbation blow-up of
Propositions~\ref{prop:conditioning} and~\ref{prop:inverse} live. This is the one
channel of Part~I that FGIB has not closed; Section~\ref{sec:limits} collects it.
\end{remark}

\begin{remark}[Family closure]\label{rem:closure}
Lemma~\ref{lem:kl-invariance} holds for the underlying Gaussian distributions. Within a
restricted variational family the equality of Proposition~\ref{prop:fgib-ceb} is exact
only if the family is closed under the inverse map $A^{-1}$: the full-covariance family
is closed (and is the parameterization used by Fischer 2020); the diagonal family is not
--- for diagonal $q, r$ the transformed pair is diagonal only if $A^{-1}$ preserves
diagonality (e.g., permutations and sign flips) and approximate otherwise. The diagonal
implementation of the codebase is therefore covered at the distribution level, with the
family approximation confined to the noise covariance $A^{-1}\Sigma_q A^{-\top}$
(Remark~\ref{rem:orth}).
\end{remark}

\begin{remark}[Rank deficiency]\label{rem:rank}
If $z < d$ or $A$ is rank-deficient, only the projection of $h$ onto the row space of
$A$ (the image of $A^{\top}$ in $h$-space) receives gradient from the regularizer; the
null-space component is unconstrained.
In the default configuration $z = d$ and a randomly initialized $A$ is almost surely
invertible, but invertibility carries no conditioning margin --- near-singularity is
enough to produce the effects of Proposition~\ref{prop:conditioning}(iii), and it is the
direction the gap gradient points in (Remark~\ref{rem:A-channel}). If $A \in
\mathbb{R}^{z \times d}$ has full column rank ($z \ge d$), a left inverse $B$ with
$BA = I$ exists and the data-processing inequality for KL gives, per $(x,y)$,
$\operatorname{KL}(q\|r) \ge \operatorname{KL}(B_\# q\,\|\,B_\# r)
\ge I(X;\,BZ\,|\,Y)$ --- a valid lower-level bound, but one whose equality
(invariance) is lost, and whose pulled-back prior again depends on $B$. The exact
invariance argument of Proposition~\ref{prop:fgib-ceb} genuinely requires the invertible
linear map $T(h) = Ah$ with $z = d$.
\end{remark}

\begin{remark}[Noise-free limit]\label{rem:limit}
The equivalence of Proposition~\ref{prop:fgib-ceb} should be understood at the level of
the variational objective, not as a mutual-information limit: as the side-path noise
vanishes, $I(X;Z|Y)$ diverges for continuous deterministic $H$ under the standing
non-atomic assumption (the pathology of
deterministic networks discussed by Amjad \& Geiger and Kolchinsky et al.), and the
variational KL bound becomes loose. Keeping the side-path covariance finite --- without
injecting noise into the inference path --- is precisely what keeps the regularizer
well-defined; this is the design rationale of FGIB.
\end{remark}

% ---------------------------------------------------------------
\subsection{The sweep}
% ---------------------------------------------------------------

\textbf{The degeneracy to be addressed (restated in our terms).}
When $Y = f(X)$, the IB curve saturates the data-processing bound and is piecewise
linear, $F(r) = \min\{r, H(Y)\}$ (Caveats paper, Section~3). Two consequences (their
Caveat~1): (i)~by the DPI, $L^{\mathrm{IB}}_\beta(T) = I(Y;T) - \beta I(X;T)
\le (1-\beta)I(Y;T) \le (1-\beta)H(Y)$, and $T_{\mathrm{copy}} := f(X) = Y$ attains the
bound for \emph{every} $\beta \in [0,1]$: all $\beta \in (0,1)$ are pinned at the single
corner point $\langle H(Y), H(Y)\rangle$ of the curve; (ii)~conversely, at $\beta = 1$
every point of the increasing part simultaneously maximizes $L^{\mathrm{IB}}_1$. The
mechanism is exposed by writing $\max_T L^{\mathrm{IB}}_\beta = \max_r [F(r) - \beta r]$
(the Legendre transform of $-F$): on the linear increasing part every point has the
\emph{same} derivative $F'(r) = 1$, hence the same supporting slope --- a Lagrangian of
the form ``prediction $-\beta\cdot{}$compression'' sweeps a Pareto curve only if each
point of the curve has a \emph{unique} supporting slope, i.e.\ the curve is strictly
concave. The dIB variant fails identically: on the hard-clustering family $T = g(Y)$ one
has $H(T|Y) = 0$, hence $I(Y;T) = H(T)$, so $L^{\mathrm{dIB}}_\beta = (1-\beta)H(T)$
--- a \emph{monotone} function of the single compression scalar $H(T)$, maximized by the
finest partition for every $\beta < 1$. The squared fixes repair this by making the
objective strictly concave in the compression scalar: $F'(r)/(2r) = \beta$ is strictly
decreasing, so each $\beta$ selects a unique $r$ (Eq.~8 argument), and squared-dIB has
the unique maximizer $H(T) = 1/(2\beta)$ (their Eq.~A20). By Proposition~\ref{prop:dead}
this degeneracy is inherited verbatim by CEB ($\beta_{\mathrm{IB}} = 1+\gamma > 1$),
which is the first channel of Part~I.

\textbf{Reduced model of FGIB.}
The KL depends on the side mean $\mu$ alone, and the classifier term is invariant under
the coordinate change $\tilde w_j = A^{-\top} w_j$ (since $w_j \cdot h =
\tilde w_j \cdot \mu$): with $A$ invertible ($z = d$ in the default configuration, so a
random $A$ is a.s.\ invertible), every equilibrium in $(h, w)$ coordinates maps
bijectively to one in $(\mu, \tilde w)$ coordinates. We therefore analyze the reduced
model $\mu(x) = h(x)$ without loss of generality.
Consider the \emph{class-symmetric aligned family}: the encoder maps every input of
class $y$ to $h_y = s\,v_y$ ($s \ge 0$), the classifier is held at the symmetric
configuration $w_j = c\,v_j$, $b_j = 0$ with fixed scale $c > 0$, and the posterior is
$q_y = \mathcal{N}(s v_y, \sigma^2 I)$. The family is the natural equilibrium manifold of
the two forces: the KL gradient
$\nabla_h \operatorname{KL} = (\mu - a v_y)/\tau^2$ acts exactly
along $v_y$, and the classification force on the class center acts along $v_y$ at the
symmetric configuration; the ansatz is a \emph{working approximation}: it captures the
direction of both forces, and is motivated as exact in the $\beta$-dominant regime, but
that exactness is not proved here --- the cross terms are of order $O(c/\beta)$ and are
idealized away (see Remark~\ref{rem:caveats}).

\begin{lemma}[FGIB Lagrangian on the aligned family]\label{lem:lagrangian}
On the aligned family,
\begin{equation}
\operatorname{KL}(s, \sigma^2)
= \tfrac{z}{2}\big(\sigma^2/\tau^2 - 1 - \ln(\sigma^2/\tau^2)\big)
+ \tfrac{1}{2\tau^2}(s-a)^2,
\qquad
\mathrm{CE}(s) = \ln\!\big(1 + (K-1)e^{-cs}\big),
\label{eq:kl-ce}
\end{equation}
and the optimizer over $\sigma^2$ is $\sigma^2 = \tau^2$ for \emph{every} $\beta$, since
the CE term does not depend on $\sigma^2$ (the classifier consumes the deterministic
$h$, not $z$).
\end{lemma}

\begin{proof}
With $h_y = s v_y$ and $w_j = c v_j$, the logits are $c\,\langle v_j, s v_y\rangle
= cs\,\delta_{jy}$, so $p_y = e^{cs}/(e^{cs} + K - 1)$ and
$\mathrm{CE} = -\ln p_y = \ln(1 + (K-1)e^{-cs})$.
For the KL, the diagonal closed form with $\Sigma_q = \sigma^2 I$,
$\Sigma_p = \tau^2 I$ reads
$\tfrac{1}{2}\sum_d\big[\sigma^2/\tau^2 + (\mu_d - a v_{yd})^2/\tau^2 - 1
- \ln(\sigma^2/\tau^2)\big]$, giving
Eq.~\eqref{eq:kl-ce} since $\|v_y\| = 1$. The variance term is non-negative with a
unique minimum at $\sigma^2 = \tau^2$, and $\mathrm{CE}$ is independent of $\sigma^2$.
\hfill\qed
\end{proof}

\begin{proposition}[Unique interior optimum per $\beta$ at fixed classifier scale --- the sweep]\label{prop:sweep}
\emph{At a fixed classifier scale $c$}, for every $\beta > 0$, the FGIB Lagrangian on
the aligned family,
\begin{equation}
L_\beta(s) := \ln\!\big(1 + (K-1)e^{-cs}\big) + \tfrac{\beta}{2\tau^2}(s-a)^2,
\label{eq:fgib-lag}
\end{equation}
is strictly convex in $s$ and has a unique minimizer $s^*(\beta) > a$, given implicitly
by
\begin{equation}
\frac{\beta}{\tau^2}\,(s-a) = c\,\varphi(cs),
\qquad
\varphi(u) := \frac{(K-1)e^{-u}}{1 + (K-1)e^{-u}} \in (0,1) .
\label{eq:stationarity}
\end{equation}
The map $\beta \mapsto s^*(\beta)$ is continuous and strictly decreasing, with
$s^*(\beta) \to a$ as $\beta \to \infty$ and $s^*(\beta) \to \infty$ as $\beta \to 0$.
Consequently the map $\beta \mapsto \big(\operatorname{KL}(s^*), \mathrm{CE}(s^*)\big)$
is one-to-one onto the entire curve
$\{\big(\tfrac{1}{2\tau^2}(s-a)^2,\; \ln(1+(K-1)e^{-cs})\big) : s \in (a, \infty)\}$:
no two
$\beta$ share the same optimum, and every interior point of the trade-off curve is
attained --- the exact failure mode of $L^{\mathrm{IB}}_\beta$ (all $\beta$ pinned at
$T_{\mathrm{copy}}$) is absent. The fixed-$c$ condition is load-bearing and is analyzed
in Remark~\ref{rem:c-scale}.
\end{proposition}

\begin{proof}
Compute $\mathrm{CE}'(s) = -c\varphi(cs)$ and
$\mathrm{CE}''(s) = c^2 (K-1) e^{-cs}/(1 + (K-1)e^{-cs})^2 > 0$, so CE is strictly
convex and strictly decreasing in $s$; the compression term $\tfrac{1}{2\tau^2}(s-a)^2$
is strictly convex. Hence $L_\beta''(s) = \mathrm{CE}''(s) + \beta/\tau^2 > 0$:
$L_\beta$ is
strictly convex with at most one minimizer, characterized by
$0 = L_\beta'(s) = -c\varphi(cs) + \beta(s-a)/\tau^2$, i.e.\
Eq.~\eqref{eq:stationarity}.
The left side $\beta(s-a)/\tau^2$ is strictly increasing from $-\beta a/\tau^2$
(at $s=0$) through $0$
(at $s=a$) to $+\infty$; the right side $c\varphi(cs)$ is strictly decreasing in $s$
with $c\varphi(cs) \in (0, c)$. The unique crossing therefore satisfies $s^* > a$.
Differentiating Eq.~\eqref{eq:stationarity} with respect to $\beta$ gives
$(s-a)/\tau^2 + (\beta/\tau^2) s' = c^2\varphi'(cs)\,s'$ with $\varphi' < 0$, so
$s' = -(s-a)/(\beta - c^2\tau^2\varphi'(cs)) < 0$: $s^*$ is strictly decreasing.
The limits follow by monotonicity: as $\beta \to \infty$ the left side can only match
the bounded right side when $s \to a$; as $\beta \to 0$ the crossing requires
$\varphi(cs)\to 0$, i.e.\ $s \to \infty$. \hfill\qed
\end{proof}

\begin{proposition}[The trade-off curve is strictly concave in the compression--prediction plane]\label{prop:curve}
On the curve of Proposition~\ref{prop:sweep} (fixed classifier scale $c$),
\begin{equation}
\frac{d\,\mathrm{CE}}{d\,\operatorname{KL}}
= \frac{\mathrm{CE}'(s)}{(s-a)/\tau^2}
= -\frac{\tau^2\,c\varphi(cs)}{s-a},
\label{eq:slope}
\end{equation}
which is strictly increasing from $-\infty$ (as $s \to a^+$) to $0^-$ (as $s \to
\infty$). Thus the curve $\operatorname{KL} \mapsto -\mathrm{CE}$ (compression versus
achievable prediction, on the KL--CE surrogate plane of
Remark~\ref{rem:surrogate}) is strictly concave, and each of its points has a \emph{unique}
supporting slope $-\beta$. By the criterion of the Caveats paper (their Fig.~2 right
panel and the $F'(r)/(2r)$ argument of Eq.~8), this is precisely the property that makes
the Lagrangian $-\mathrm{CE} - \beta\,\operatorname{KL}$ sweep the full \emph{surrogate}
curve on the aligned family at fixed scale: the convexity of the compression term
supplies the strict-concavity condition \emph{without squaring it}. This is a statement
about the KL--CE surrogate of Proposition~\ref{prop:sweep}, not about the Caveats
information plane (Remark~\ref{rem:surrogate}).
\end{proposition}

\begin{proof}
$\varphi$ is positive and strictly decreasing and $s-a$ is positive and strictly
increasing, so $\tau^2 c\varphi(cs)/(s-a)$ is strictly decreasing in $s$; the signs give
the stated limits. The stationarity condition Eq.~\eqref{eq:stationarity} is exactly
$-d\mathrm{CE}/d\operatorname{KL} = \beta$, i.e.\ $\beta$ is the supporting slope at the
optimum. \hfill\qed
\end{proof}

\begin{remark}[What the sweep sweeps: a KL--CE surrogate curve, not the IB curve]\label{rem:surrogate}
Propositions~\ref{prop:sweep} and~\ref{prop:curve} prove strict convexity of the
FGIB Lagrangian and one-to-one parameterization by $\beta$ \emph{on the class-symmetric
aligned family, at fixed classifier scale}, in the plane
$\big(\operatorname{KL}_{\mathrm{anchor}},\, \mathrm{CE}\big)$. That plane is not the
Caveats paper's IB plane $\big(I(X;T),\, I(Y;T)\big)$, and the claim does not extend to
it. On the aligned family the posterior is class-constant by construction, so
$I(X;Z|Y) = 0$ for every $s$, while the deterministic class prototypes $h_y = s v_y$
already retain the label completely for any $s > 0$ ($I(Y;H) = H(Y)$); what varies
along the sweep is the anchor mismatch and the classifier margin, not a proved
compression of mutual information. Whether the deployed representation's residual
$I(X;H|Y)$ moves monotonically with $\beta$ is a separate, empirical question; the
theory guarantees a strictly concave surrogate, not a recovered IB curve. This
qualification is carried into Section~\ref{sec:limits}.
\end{remark}

\begin{remark}[The classifier scale is load-bearing, and the current implementation leaves it free]\label{rem:c-scale}
If $c$ is free, the pair $(s, c) = (a, c)$ with $c \to \infty$ sends
$\operatorname{KL} = \tfrac{1}{2\tau^2}(s-a)^2 = 0$ and
$\mathrm{CE} = \ln(1 + (K-1)e^{-ca}) \to 0$ simultaneously (for $a > 0$): the infimum
of $L_\beta$ over the aligned family is $0$ for \emph{every} $\beta$, no minimizer is
attained, and the one-to-one sweep of Proposition~\ref{prop:sweep} has no global
counterpart. (At $a = 0$ the margin growth alone cannot drive CE to zero.) This is
the FGIB analogue of the dead knob --- the same one-corner attractor, now arrived at by
growing the margin rather than by the linearity of the deterministic IB curve.
Three consequences.
\emph{(a) The present implementation is exposed.} The classifier is
\texttt{nn.Linear} trained by Adam with no weight decay, so $c$ is unconstrained. On
separable data the CE gradient drives $\|w_y\|$ upward; it does not diverge in finite
time (the gradient vanishes as $p_y \to 1$, giving logistic growth), but at any finite
scale $c$ the equilibrium $s^*(\beta; c)$ sits on the curve of
Proposition~\ref{prop:sweep} \emph{for that $c$}, and $s^*$ drifts toward $a$ as $c$
grows.
\emph{(b) Testable prediction.} The $\beta$-dependence of the equilibrium should decay
with training length; since the experiments use early stopping with $\texttt{patience} =
10$, part of any observed $\beta$-effect is plausibly an early-stopping effect. This can
be tested by training well past saturation and comparing the converged KL across
$\beta$.
\emph{(c) Architectural fix.} The device that satisfies the fixed-$c$ assumption
\emph{by construction} is a cosine classifier: normalized features, normalized rows,
and a fixed temperature (scale). Weight normalization with a learnable gain does not
guarantee a bounded scale, and weight decay only replaces the free-scale problem with a
joint optimization problem over $(s, c)$ that must be re-derived --- it mitigates the
divergence but does not make Proposition~\ref{prop:sweep} exact. Under a fixed-norm
classifier with fixed temperature the proposition holds as stated, at the price of one
architecture change.
\end{remark}

\textbf{Where the convexity comes from (mechanism).}
Compare the three objectives on their natural solution families, as functions of the
single compression scalar:
\begin{itemize}
\item \emph{dIB on hard clusterings} (Caveats Appendix~B): $L^{\mathrm{dIB}}_\beta
= (1-\beta)H(T)$ --- compression $H(T)$ and prediction $I(Y;T) = H(T)$ are the
\emph{same} scalar, so the objective is monotone in it and the optimum sits at the
endpoint for every $\beta < 1$. Squared-dIB restores interior optima by hand,
$H(T) - \beta H(T)^2$ with unique maximizer $H(T) = 1/(2\beta)$.
\item \emph{Standard IB} (deterministic): the curve is linear ($F'(r) \equiv 1$), every
point shares one supporting slope, and the corner $\langle H(Y), H(Y)\rangle$ is
selected for all $\beta \in (0,1)$. Squared-IB convexifies via $F'(r)/(2r)$ strictly
decreasing.
\item \emph{FGIB}: the compression term is $\operatorname{KL}(q_y \| r_y) =
\tfrac{1}{2\tau^2}\|\mu_y - a v_y\|^2 + {}$(variance terms), equal to
$\tfrac{1}{2\tau^2}(s-a)^2$
at the variance optimum --- \emph{quadratic} in the alignment scalar $s$, hence strictly
convex, because the KL is measured against a \emph{fixed} reference $a v_y$:
$\operatorname{KL}(\mathcal{N}(\mu,\cdot)\|\mathcal{N}(\mu_p,\cdot))
= \tfrac{1}{2\tau^2}\|\mu - \mu_p\|^2 + {}$(variance terms) is quadratic in the
displacement
$\mu - \mu_p$ whenever $\mu_p$ does not move. Entropy-based compression ($H(T)$,
$I(X;T)$) is self-referential --- it depends only on the representation's own
distribution and enters linearly in the clustering probabilities; the fixed anchor
supplies an absolute reference frame and turns compression into a squared distance. The
convexification that squared-IB installs by squaring the compression term is thus built
into FGIB's fixed-anchor KL. Together with the strictly convex prediction term
(Lemma~\ref{lem:lagrangian}), the Lagrangian is strictly convex with a unique interior
minimizer per $\beta$ \emph{at fixed classifier scale} (Proposition~\ref{prop:sweep}),
and the resulting curve is strictly concave (Proposition~\ref{prop:curve}). With the
scale free, the same quadratic gap is driven to zero by margin growth alone
(Remark~\ref{rem:c-scale}) --- the convexification is necessary but not sufficient
without the scale condition.
\end{itemize}

\begin{remark}[The compressed endpoint is non-collapsed but still a label code --- Caveat 2 remains]\label{rem:endpoint}
VIB's fixed marginal $m(z) = \mathcal{N}(0, I)$ places the $\operatorname{KL}=0$ point
at $\mu = 0$ for \emph{all} classes: the fully compressed solution is the trivial
collapse with $\mathrm{CE} = \ln K$ (chance level), and squaring (SVIB) leaves this
corner in place since $\operatorname{KL}^2$ is stationary at $\operatorname{KL}=0$.
FGIB's per-class orthonormal anchors make the $\operatorname{KL}=0$ point a
\emph{class-separated} configuration: $\mu_y = a v_y$ with
$\langle v_y, v_{y'}\rangle = \delta_{yy'}$, so at the corner
$\mathrm{CE}(a) = \ln(1 + (K-1)e^{-ca})$, which is bounded and tends to $0$ as the
anchor scale $a$ grows, and $a = 0$ recovers the VIB corner. But class-separation does
not buy non-triviality in the sense of the Caveats paper: their Caveat~2 identifies as
uninteresting the variables on the manifold $T_\alpha$ that mix a constant map with an
exact copy of $Y$, compressing by ``forgetting the input with some probability'' rather
than by merging perceptually similar classes --- and FGIB's compressed corner
$\mu_y = a v_y$ is a deterministic function of $Y$ alone, i.e.\ a label code of exactly
that kind. The orthonormal frame is in fact \emph{opposed} to the kind of compression
Caveat~2 holds up as interesting: it places every pair of classes at the same distance
and can never express ``merge collie with retriever, keep both far from teapot''. What
the anchors do is choose \emph{one specific} label-code geometry --- class means on a
scaled simplex, a neural-collapse-like configuration --- instead of leaving it to the
optimizer, and refuse the chance-level corner. The compressed endpoint is thus designed,
not certified, and any claim that the method compresses \emph{input} semantics at that
endpoint is exactly the claim Caveat~2 rules out.
\end{remark}

\begin{remark}[Residual-information interpretation (CEB)]\label{rem:residual}
For any backward encoder, $\mathbb{E}[\operatorname{KL}(q(z|x)\|r(z|y))]
= I(X;Z|Y) + \mathbb{E}_y[\operatorname{KL}(q(z|y)\|r(z|y))]
\ge I(X;Z|Y)$ (Fischer 2020, Eqs.~9--11, 20; Lemma~\ref{lem:gap}). FGIB is therefore a
CEB bottleneck with a \emph{fixed} backward encoder: the regularizer upper-bounds the
residual information, and on the aligned family the residual vanishes identically
($q(z|x)$ depends on $y$ alone), so the regularizer equals the anchor-matching gap
$\tfrac{1}{2\tau^2}(s-a)^2$. The quadratic mechanism of Proposition~\ref{prop:sweep} is
inherited from this gap term, which exists only because the anchor is not allowed to
follow the encoder.
\end{remark}

\begin{remark}[Variance decoupling and training--evaluation consistency]
Since the classifier consumes the deterministic $h$ (and the KL does not see the
classifier), the posterior variance decouples: $\sigma^{2*} = \tau^2$ for all $\beta$, and
the trade-off is purely geometric (mean alignment). This contrasts with VIB, where the
classifier consumes the sampled $z$: the noise couples into the prediction term (by
Jensen, $\mathbb{E}_z[\mathrm{CE}(c(z),y)] \ge \mathrm{CE}(c(\mu),y)$), the rate
$R \ge I(X;Z)$ over-counts compression, and training-time sampling versus
evaluation-time mean creates a distribution mismatch (Fischer 2020, Appendix~A.1).
\end{remark}

\begin{remark}[Idealizations]\label{rem:caveats}
The class-symmetric ansatz idealizes the equilibrium geometry; it is exact in the
$\beta$-dominant regime and captures the direction of the KL force in general. For
regression, the continuous RFF anchors satisfy $\|\mu_p(y)\| \approx a$, so the KL again
contains the quadratic displacement term
$\tfrac{1}{2}\|\mu - \mu_p(y)\|^2 + {}$(variance terms) --- the same convexity mechanism
carries over, again subject to the fixed-scale condition of Remark~\ref{rem:c-scale}
(the regression classifier has no margin scale; its norm is likewise unconstrained in
the implementation). Training starts at $\operatorname{KL} \approx a^2/2$
(zero-initialized logvar head, $\sigma = 1$, $\mu \approx 0$), so the anchor scale sets
the initial regularization pressure; larger $a$ delays the approach to $s^*$. The
finite-scale caveat formerly in this remark is stated, with its fix, in
Remark~\ref{rem:c-scale}.
\end{remark}

% ---------------------------------------------------------------
\subsection{The certificate}
% ---------------------------------------------------------------

\begin{lemma}[Gap identity]\label{lem:gap}
For any conditional $r(z|y)$,
\begin{equation}
\mathbb{E}\big[\operatorname{KL}(q(z|x)\,\|\,r(z|y))\big]
= I(X;Z|Y) + \mathbb{E}_y\big[\operatorname{KL}(q(z|y)\,\|\,r(z|y))\big] .
\label{eq:gap}
\end{equation}
\end{lemma}

\begin{proof}
$\mathbb{E}[\operatorname{KL}(q(z|x)\|r(z|y))] = \mathbb{E}[\ln q(z|x) - \ln r(z|y)]$
and $I(X;Z|Y) = \mathbb{E}[\ln q(z|x) - \ln q(z|y)]$ (Fischer 2020, Eqs.~9--10), where
both expectations are over $p(x,y)q(z|x)$; subtracting,
$\mathbb{E}[\ln q(z|y) - \ln r(z|y)] = \mathbb{E}_y[\operatorname{KL}(q(z|y)\|r(z|y))]$
since the argument depends on $z, y$ only. \hfill\qed
\end{proof}

\begin{lemma}[Moment matching]\label{lem:mm}
Within the per-class Gaussian family with \emph{diagonal covariance}, the gap of
Lemma~\ref{lem:gap} is
minimized by the moment-matched Gaussian $b^*_y = \mathcal{N}(m_y, S_y)$ with
$m_y = \mathbb{E}\,q(z|y)$, $S_y = \operatorname{diag}\operatorname{Var} q(z|y)$.
\end{lemma}

\begin{proof}
$\operatorname{KL}(p\|q)$ as a function of $q$ is minimized by matching the first two
moments of $p$ (standard; for Gaussians the minimizer is unique). \hfill\qed
\end{proof}

\begin{proposition}[Variational ordering: CEB's bound is tighter in value]\label{prop:ordering}
For every fixed encoder $q(z|x)$,
$\mathrm{ReX}(\theta^*) \le \mathrm{ReX}_F$ where $\theta^*$ is CEB's optimal backward
encoder, and the excess is the \emph{anchor gap}
\begin{equation}
\Delta := \mathrm{ReX}_F - \mathrm{ReX}(\theta^*)
= \mathbb{E}_y\big[\operatorname{KL}(q(z|y)\,\|\,N(a v_y, \tau^2 I))\big]
- \mathbb{E}_y\big[\operatorname{KL}(q(z|y)\,\|\,N(m_y, S_y))\big] \ge 0 ,
\label{eq:anchor-gap}
\end{equation}
with equality iff the anchors equal the moment-matched marginals. Here $\theta^*$ ranges
over the \emph{ideal} CEB backward encoder family --- per-class Gaussians with free mean
and \emph{free diagonal covariance} (the family of Lemma~\ref{lem:mm}, matching the
diagonal implementation), i.e.\ any conditional $r(z|y)$ of that family the theory can
express; the code's
single affine \texttt{prior\_net} is only an approximation of this family, and the
ordering is a statement about the idealized CEB, not about the implemented
parameterization.
On the aligned family of the sweep analysis ($q(z|y) = \mathcal{N}(s v_y, \sigma^2 I)$,
so $I(X;Z|Y) = 0$ and $b^* = q(z|y)$ exactly),
\begin{equation}
\Delta = \tfrac{1}{2\tau^2}(s-a)^2 + \tfrac{z}{2}\big(\sigma^2/\tau^2 - 1 - \ln(\sigma^2/\tau^2)\big) ,
\label{eq:aligned-gap}
\end{equation}
i.e.\ $\mathrm{ReX}_F = \Delta$ and $\mathrm{ReX}(\theta^*) = 0$: FGIB's bound value
\emph{equals} the sweep-driving compression term of Proposition~\ref{prop:sweep}, while
CEB's bound collapses to the (zero) residual.
\end{proposition}

\begin{proof}
The anchor $N(a v_y, \tau^2 I)$ is a member of CEB's per-class Gaussian family, so the
minimum over the family is no larger than the value at the anchor: the ordering and
Eq.~\eqref{eq:anchor-gap} follow from Lemmas~\ref{lem:gap} and~\ref{lem:mm}.
On the aligned family the marginal is itself Gaussian, so $b^* = q(z|y)$ gives gap $0$,
$I(X;Z|Y) = 0$ (the posterior is class-constant), and
$\operatorname{KL}(\mathcal{N}(s v_y, \sigma^2 I)\|\mathcal{N}(a v_y, \tau^2 I))
= \tfrac{1}{2\tau^2}(s-a)^2 + \tfrac{z}{2}(\sigma^2/\tau^2-1-\ln(\sigma^2/\tau^2))$.
\hfill\qed
\end{proof}

\begin{remark}[Tightness--force trade-off]\label{rem:force}
Proposition~\ref{prop:ordering} is the price of the fixed anchor, and it is paid for a
reason: in CEB the gap is a function of the \emph{backward encoder alone} ---
$b_\theta \to q(z|y)$ eliminates it with the encoder held fixed, so CEB's bound can
become tight while the representation stays uncompressed: its tightness certifies
nothing about compression, and the gap is an unobservable optimization residual that
tracks the forward encoder. FGIB's gap can shrink only if the \emph{posterior} moves
toward the anchors (the anchor itself cannot move), so the force is persistent and, on
the aligned family, the bound value \emph{is} the anchor-matching gap
$\tfrac{1}{2\tau^2}(s-a)^2$, the sweep-driving compression term
(Eq.~\eqref{eq:aligned-gap}). The qualifier is necessary: ``the posterior'' includes
the side head $A$, which is trainable and unused at inference, so part of the movement
can be absorbed without the deployed representation $h$ moving
(Remark~\ref{rem:A-channel}). On the aligned family, where $A$ is fixed by the ansatz,
the equivalence is exact; off it, bound tightness is monotonically equivalent to
alignment of the posterior, not to compression of the deployed $h$.
\end{remark}

\begin{proposition}[The inverse linear transform: the object of the bound]\label{prop:inverse}
Let the side head $T(h) = Ah$ be \emph{invertible} ($z = d$, unconstrained), and let
the trained posterior be isotropic, $\Sigma(x) = \sigma_p^2 I$ with $\sigma_p^2 = \tau^2$
the anchor variance (the KL variance term is minimized at $\sigma^2 = \tau^2$, so this
holds at the variance optimum).
Then $\tilde z := A^{-1} z$ gives
$q'(\tilde z|x) = \mathcal{N}(h(x), \tau^2 (A^{\top}A)^{-1})$ --- the posterior is
centered at the \emph{inference representation} $h$ with side noise
$\tau^2 (A^{\top}A)^{-1}$ --- and the transformed anchors
$r'(\tilde z|y) = \mathcal{N}(a\, A^{-1} v_y, \tau^2 (A^{\top}A)^{-1})$ form an
orthonormal frame at scale $a$ in the induced metric $A^{\top}A$
($\langle A^{-1}v_i, A^{-1}v_j\rangle_{A^{\top}A} = \delta_{ij}$: the fixed frame
$\{v_y\}$ read in canonical coordinates). By KL invariance under invertible linear maps
(Lemma~\ref{lem:kl-invariance}),
\begin{equation}
\mathrm{ReX}_F
= I(X;\, h + \tau\,(A^{\top}A)^{-1/2}\varepsilon\,|Y)
+ \mathbb{E}_y\big[\operatorname{KL}\big(q'(\tilde z|y)\,\|\,r'(\tilde z|y)\big)\big],
\qquad \varepsilon \sim \mathcal{N}(0, I) .
\label{eq:object}
\end{equation}
The bound's object is the inference representation perturbed by \emph{known,
side-path-only} noise (its covariance is a model parameter), and the anchor geometry is
preserved exactly under the transform, in the induced metric. The perturbation,
however, is small only if the side head is well-conditioned: its magnitude is governed
by $A^{\top}A$, which is unconstrained
(Proposition~\ref{prop:conditioning}(iii)).
\end{proposition}

\begin{proof}
$A^{-1}A = I$ and $A^{-1} z = h + A^{-1}\varepsilon$ for $\varepsilon \sim
\mathcal{N}(0, \Sigma)$; $A^{-1}\Sigma A^{-\top} = \tau^2 A^{-1}A^{-\top}
= \tau^2 (A^{\top}A)^{-1}$ for $\Sigma = \tau^2 I$; the prior covariance maps to
$A^{-1} (\tau^2 I) A^{-\top} = \tau^2 (A^{\top}A)^{-1}$; and
$\langle A^{-1}v_i, A^{-1}v_j\rangle_{A^{\top}A}
= v_i^{\top} A^{-\top} A^{\top} A\, A^{-1} v_j = \delta_{ij}$.
Equation~\eqref{eq:object} is Lemma~\ref{lem:gap} applied to the transformed pair.
\hfill\qed
\end{proof}

\begin{proposition}[Orthogonal-factor conditioning: the reference is fixed and well-conditioned]\label{prop:conditioning}
FGIB's anchor frame $V = [v_1, \dots, v_K]$ is the orthonormal factor of a random matrix
(QR: $G = QR$, $V := Q$, $V^{\top}V = I_K$), and the anchor covariance is
$\tau^2 I$ with $\tau^2 > 0$.
Then, for every anchor scale $a \ge 0$ and throughout training:
\begin{enumerate}
\item[(i)] the \emph{reference} of the matching --- the \emph{unscaled} mean frame $V$
and the covariance $\tau^2 I$ --- is fixed (non-trainable buffers) and has condition
number $1$: the certificate is a distance to a compact, non-degenerate,
class-separated (for $a > 0$) target. (At $a = 0$ the mean matrix $a V$ is zero and
has no condition number; the condition-$1$ statement attaches to $V$ and $\tau^2 I$,
not to $a V$.)
\item[(ii)] in side-mean coordinates the certificate's displacement metric is a
scaled identity: the KL's displacement term is $\tfrac{1}{2\tau^2}\|\mu - a v_y\|^2$
with gradient $(\mu - a v_y)/\tau^2$;
\item[(iii)] on the shared-encoder side the metric is $A^{\top}A$, where $A$ is a
learnable unconstrained map, so its \emph{global} condition number is not controlled by
construction --- and the aligned family does not control it either. An earlier version
inferred condition number $1$ on the class-mean subspace from the identity
$h_i^{\top} A^{\top} A\, h_j = s^2\delta_{ij}$; that inference is invalid, since
$\{h_i\}$ is not a Euclidean orthonormal basis. Counterexample: with
$A = \operatorname{diag}(\varepsilon, 1)$, $h_1 = (1/\varepsilon, 0)$,
$h_2 = (0, 1)$ and $s = 1$, one has $A h_i = v_i$ yet
$\kappa(A^{\top}A) = \varepsilon^{-2}$ --- and the inverse-transform perturbation of
Proposition~\ref{prop:inverse} is correspondingly $\varepsilon^{-2}$-large on the first
direction. The counterexample is not adversarial: it is the direction the gap gradient
points in, since the KL is minimized by fitting the anchor target on $h$ in least
squares, which suppresses within-class directions
(Remark~\ref{rem:A-channel}). Any guarantee on this side therefore requires an explicit
constraint on $A$ (orthogonal parameterization, weight norm, or the like) or an
empirical measurement of $\kappa(A^{\top}A)$ over training.
\end{enumerate}
CEB's backward encoder $b_\theta(z|y) = \mathcal{N}(m_\theta(y), S_\theta(y))$ is an
unconstrained linear map of the label: \emph{both} sides of its matching are trainable,
its covariance diagonal can lie anywhere in the clamp $[e^{-10}, e^{10}]$ (condition
number up to $e^{20}$), its means can be rank-deficient (class-merged codes, or the
collapse solution $m_\theta \equiv 0$), and the variance race drives $S_\theta$ to the
clamp, degenerating the certificate into a mean-matching force of weight $1/s = e^{10}$
(Prop.~\ref{prop:variance}, subject to Remark~\ref{rem:p3-status}). What the fixed QR
reference guarantees for FGIB is therefore the \emph{reference} side only: a compact,
non-degenerate, class-separated target whose conditioning cannot degrade. On the
encoder side FGIB's bound is exposed to the same degeneration mechanism CEB's is, with
$A$ playing the role of the second trainable party.
\end{proposition}

\begin{proof}
$V^{\top}V = I_K$: all singular values of $V$ are $1$; the covariance $\tau^2 I$ has
all eigenvalues $\tau^2$; both are frozen buffers, so (i) holds. The diagonal closed
form of
$\operatorname{KL}(\mathcal{N}(\mu, \sigma^2 I)\,\|\,\mathcal{N}(a v_y, \tau^2 I))$
has displacement term $\tfrac{1}{2\tau^2}\|\mu - a v_y\|^2$, giving (ii). For (iii):
the identity
$h_i^{\top} A^{\top} A\, h_j = \mu_i^{\top}\mu_j = s^2 v_i^{\top} v_j = s^2\delta_{ij}$
holds on the aligned family and only says that $\{h_i\}$ is $A^{\top}A$-orthogonal by
construction; it implies nothing about the spectrum of $A^{\top}A$, and the displayed
example verifies the failure. The KL gradient on $A$ is
$(\mu - a v_y)\,h^{\top}/\tau^2$, an alignment force toward the orthonormal frame that
vanishes only at perfect alignment \emph{on non-degenerate samples} (it also vanishes
pointwise at $h = 0$, or by cancellation across a batch) --- but nothing in that force
resists a vanishing within-class component of $A$. For CEB, the displacement term is
$\tfrac12(m - g)^{\top} S_\theta^{-1}(m - g)$: the metric $S_\theta^{-1}$ has condition
number up to $e^{20}$ under the logvar clamp, and the gradient weight $1/s$ at the
collapse equilibrium is Prop.~\ref{prop:variance}. \hfill\qed
\end{proof}

\begin{remark}[Two senses of tightness]\label{rem:senses}
In the pure variational \emph{value} sense (same encoder, same Gaussian family), CEB's
bound is tighter: $\mathrm{ReX}(\theta^*) \le \mathrm{ReX}_F$
(Prop.~\ref{prop:ordering}), the excess being the anchor gap --- the price of the fixed
anchor. The \emph{object} sense does not reverse this ordering: both bounds are valid
upper bounds on their own noisy codes, and neither bounds the residual of a deployed
deterministic continuous representation, which is infinite under the standing
non-atomic assumption (Prop.~\ref{prop:object}). The structural difference is where the certified-versus-
deployed gap lives --- in CEB it is driven by the prediction term inside the inference
path; in FGIB it is a side-path quantity whose size is governed by the conditioning of
$A$ (Prop.~\ref{prop:conditioning}(iii)) --- not which bound is ``tighter'' or
``valid''. The learnability of the side head $A$ bears on that difference: it is a
second trainable party in the matching, unused at inference
(Remark~\ref{rem:A-channel}), so it can absorb part of the anchor gap without moving
the deployed representation. On the aligned family, where $A$ is fixed by the ansatz,
its optimization drives the gap to the equilibrium value $\tfrac12(s^*-a)^2$ at fixed
classifier scale (Prop.~\ref{prop:sweep}); off the ansatz this is an open question
(Section~\ref{sec:limits}).
\end{remark}

\begin{remark}[Unconstrained side head (the implementation)]\label{rem:orth}
Proposition~\ref{prop:inverse} holds for any invertible linear head: the KL invariance
is exact, and the object claim needs only that the perturbation covariance
$\tau^2(A^{\top}A)^{-1}$ be \emph{known} ($A$ is a model parameter) --- but the
object is informative about $h$ only insofar as that covariance is moderate, which the
construction does not guarantee (Prop.~\ref{prop:conditioning}(iii)). The diagonal
variational family is not closed under $A^{-1}$ (Remark~\ref{rem:closure}), so within
the diagonal parameterization of the codebase the equivalence is exact at the
distribution level and approximate as a variational identity; a full-covariance side
head would make it exact within the family. No orthogonality constraint on $A$ is
required --- the orthogonality of the method lives in the fixed anchor frame and in the
equilibrium it induces (the aligned family).
\end{remark}

\begin{remark}[Anchor-frame reading]\label{rem:frame}
With the orthonormal anchor matrix $V$ (columns $v_y$, $K \le z$), left multiplication
by $V^{\top}$ expresses the anchor-space components of $z$ in class coordinates: the
anchors become $a\,e_y$ and the gap of Lemma~\ref{lem:gap} decomposes into $K$
per-class terms plus, when $K < z$, an orthogonal-complement term that the reading
ignores ($V^{\top}$ is a projection, not a coordinate change). The orthogonal frame is
the geometric content of the prior, and it is preserved by the invertible linear
pull-back $A^{-1}$ of Proposition~\ref{prop:inverse} --- not by $A^{\top}$, which is
the Jacobian factor that appears in gradients and equals $A^{-1}$ only when $A$ is
orthogonal. An earlier version conflated the two.
\end{remark}

\begin{remark}[Positioning against VIB and NIB]\label{rem:position}
VIB's rate bounds the \emph{full} mutual information $I(X;Z)$ with a marginal prior;
CEB removes the retained label information ($\mathrm{ReX} = R - I(Y;Z)$ at the optima),
and FGIB inherits the residual-level object while replacing the learnable backward
encoder with a fixed geometric prior (Prop.~\ref{prop:ordering}).
NIB's non-parametric bound $\widehat I \ge I(X;M)$ is $y$-unconditioned and tight only
when the components share a covariance and form well-separated clusters (Kolchinsky et
al.\ 2017, Eq.~10 and the discussion following Eq.~11); when the noise is small relative
to the cluster spread it saturates at $\ln N$ and its gradient vanishes (their
Section~IVA) --- loose exactly in the regime where the parametric FGIB bound is
well-behaved. The two bounds are complementary: NIB certifies the full compression
non-parametrically when the geometry is favorable; FGIB bounds the residual of a
\emph{noisy copy} of the deployed representation parametrically, with a geometry of its
own choosing --- subject to the conditioning caveat of
Proposition~\ref{prop:conditioning}(iii).
\end{remark}

\begin{remark}[Degenerate anchors]\label{rem:a0}
At $a = 0$ the anchors collapse to $N(0, \tau^2 I)$ for all classes and the gap loses its
class
structure: $\Delta = \mathbb{E}_y[\operatorname{KL}(q(z|y)\|N(0,\tau^2 I))]$ --- a
VIB-style
marginal gap; the structural separation of Proposition~\ref{prop:object}(b) still
holds, but the compressed endpoint becomes the trivial collapse
(Remark~\ref{rem:endpoint}).
\end{remark}

% ---------------------------------------------------------------
\subsection{What is proved and what is open}\label{sec:limits}
% ---------------------------------------------------------------

\begin{remark}[Scope of the results]
\textbf{Proved.} (1) The dead-knob proposition for CEB and the MNI reading of it
(Prop.~\ref{prop:dead}, Remark~\ref{rem:mni}); (2) the gap identity, Gaussian moment
matching, and the value ordering of CEB's bound over FGIB's at a fixed encoder
(Lemmas~\ref{lem:gap}, \ref{lem:mm}; Prop.~\ref{prop:ordering}); (3) the blind
certificate on the $T_\alpha$ manifold and the $1{:}\beta$ exchange rate
(Prop.~\ref{prop:flat}); (4) KL invariance under invertible linear maps and the
pointwise equivalence of the side path to a CEB bottleneck at a fixed side head
(Lemmas~\ref{lem:kl-invariance}, \ref{lem:mi-invariance}; Prop.~\ref{prop:fgib-ceb});
(5) the variance decoupling $\sigma^{2*} = \tau^2$ and the non-collapsed compressed endpoint
(Lemma~\ref{lem:lagrangian}, with $\sigma^{2*} = \tau^2$;
Remark~\ref{rem:endpoint}); (6) the one-dimensional sweep
and its strict concavity \emph{at fixed classifier scale}, as a KL--CE surrogate curve
on the aligned family (Props.~\ref{prop:sweep}, \ref{prop:curve};
Remark~\ref{rem:surrogate} --- not a recovered IB curve); (7) the well-conditioning of the fixed
reference side (Prop.~\ref{prop:conditioning}(i)--(ii)).

\textbf{Open or conditional.} (1) The side head $A$ is a trainable, inference-unused
party in the matching (Remark~\ref{rem:A-channel}); combined with the
counterexample of Prop.~\ref{prop:conditioning}(iii), the perturbation in
Prop.~\ref{prop:inverse} can be arbitrarily large, so the bound's object may be far from
the deployed representation exactly where the residual lives. A constraint on $A$ or a
measurement of $\kappa(A^{\top}A)$ over training would close or quantify this channel.
(2) The global sweep requires a bounded classifier scale
(Remark~\ref{rem:c-scale}); the current implementation leaves the scale free and the
proposition conditional. A fixed-temperature cosine classifier would satisfy the
fixed-$c$ assumption by construction; weight decay only mitigates the scale divergence
and replaces the problem with a new joint optimization over $(s,c)$ that must be
re-derived. (3) Neither CEB's nor FGIB's certificate upper-bounds the deployed
deterministic residual (Prop.~\ref{prop:object}); claims of inference-representation
compression are hypotheses, not conclusions. (4) Caveat~2 is not resolved
(Remark~\ref{rem:endpoint}): the compressed endpoint is a chosen label-code geometry,
not certified semantic compression. (5) The two-stage variance dynamics of
Prop.~\ref{prop:variance}(ii)--(iii) await the measurement of
Remark~\ref{rem:p3-status}. (6) The regression RFF anchor analysis inherits the same
fixed-scale and conditioning caveats and is presently a transfer argument, not a proof.
\end{remark}
。
% 记号沿用 Fischer (2020)、Kolchinsky & Tracey (2019, Caveats)、Kolchinsky et al. (2017, NIB)、
% Alemi et al. (2017, VIB)。
%
% 核心结论（FGIB 论文问题陈述与解法）：
% 第一部分（四条失效通道）：
%   P1 退化继承（死旋钮）：CEB ≡ IB(β=1+γ>1)，确定性场景下所有 γ>0 的唯一最优解钉在角点，
%      γ 扫不出任何中间点（SIB Caveat 1 的退化半区，经 β'=1/(1+γ)∈(0,1) 映射）。
%      该角点恰是 CEB 自己的 MNI 目标，rem:mni 正面处理这一反驳；
%   P2 证书盲视 + 1:β 交换率：T_α 流形上 ReX 恒为 0（证书对标签保留量完全盲视），
%      且目标以 1 nat 标签信息 ↔ β nat 认证残差的比率交换；
%   P3 方差竞赛：KL 两侧皆可训练，分类项压 σ→0、钳位处退化为权重 e^{10} 的均值匹配。
%      (ii)(iii) 是受力方向分析而非收敛证明，且可实测（rem:p3-status）；
%   P4 对象错位：两种方法的界都只界住各自的含噪码，都不界住确定性连续推理表示的残差
%      （后者为 +∞）。差别是结构性的（噪声在不在推理路径上），不是紧度排序。
% 第二部分（FGIB 机制，以及它没能关闭的通道）：
%   M1 等价性：旁路 = 固定 A 下作用在 h+已知噪声上的 CEB 瓶颈（逐步成立，非全程固定参考）；
%   M2 扫参：固定分类器尺度 c 时 L_β(s)=CE(s)+(β/2)(s−a)² 严格凸、每 β 唯一内点、曲线严格凹；
%      c 自由时下确界对所有 β 均为 0（rem:c-scale，给出权重归一化分类器的架构修法）；
%   M3 证书：gap 恒等式 + 矩匹配 + 固定编码器下的值排序；固定参考的力持续性成立，
%      但侧头 A 是推理无关的第二个自由方（rem:A-channel），构成未关闭的通道。
%
% 修订记录（第二轮）：理论/实现分层后统一引入锚点方差 τ²（先验 N(a v_y, τ²I)，默认 τ²=1，
% 后验方差最优 σ²*=τ²，σ_p 为公共方差记号；point-anchor 极限统一写为 τ→0）；理论旁路固定为
% 可逆线性头 T(h)=Ah（实现用带 bias 的 nn.Linear 近似，见理论-代码分层 caveat）；
% 明确 sweep 证明的是对齐族上 KL–CE surrogate 曲线（rem:surrogate），不是 IB 互信息曲线；
% 无穷互信息结论统一挂靠分布性 standing assumption；rem:c-scale 的架构修法收紧为
% 固定温度 cosine 分类器（weight decay 只缓解不发散）；a=0 退化点、条件数表述、
% 梯度消失条件、row-space 表述等小项修正。
%
% 修订记录（第一轮）：删除三处错误论断——(1) P3 原称 ReX≈0 时 I(X;Z|Y) 发散，与 gap 恒等式矛盾；
% (2) prop:object(c) 原称 FGIB 的界"更紧/唯一有效"，由 DPI 不成立；(3) prop:conditioning(iii)
% 原从 Gram 恒等式推出条件数为 1，有反例。修正 P2 的换算、Caveat 2 的主张、rem:frame 的
% A^T/V^T、mu_head 的仿射与非推理路径身份。未证结论统一收入 sec:limits。

\newtheorem{lemma}{Lemma}
\newtheorem{proposition}{Proposition}
\newtheorem{remark}{Remark}

% =====================================================================
\section*{Part I. The four failure channels of the conditional bottleneck}
% =====================================================================

\textbf{Notation.}
$X$ is the input and $Y$ the label with $K$ classes; the encoder induces the Markov chain
$Y - X - Z$. The \emph{deterministic scenario} is $Y = f(X)$ ($H(Y|X) = 0$). Following
Fischer (2020), we write the MNI point $I(X;Y) = I(X;Z) = I(Y;Z)$ (Eq.~1), the
chain-rule identity $I(X;Z|Y) = I(X;Z) - I(Y;Z)$ (Eq.~4), the CEB objective
$\mathrm{CEB} \equiv \min_Z I(X;Z|Y) - \gamma I(Y;Z)$ (Eq.~5), the IB objective
$\mathrm{IB} \equiv \min_Z I(X;Z) - \beta_{\mathrm{IB}} I(Y;Z)$ (Eq.~16), the
variational objective
$\mathrm{VCEB} \equiv \min_{e,b,c} \langle \log e(z|x)\rangle - \langle \log b(z|y)\rangle
- \gamma \langle \log c(y|z)\rangle$ (Eq.~15), and the residual certificate
$\mathrm{ReX} \equiv \langle \log e(z|x)\rangle - \langle \log b(z|y)\rangle
\ge I(X;Z|Y)$ (Eq.~20), where $b(z|y)$ is a \emph{trainable} backward encoder (we write
$q(z|x)$ for Fischer's encoder $e(z|x)$; the two notations are interchangeable
throughout). In the loss convention $\mathrm{CE} + \beta\,\mathrm{KL}$ used throughout, the weights are
$\beta = 1/\gamma$ and $\beta_{\mathrm{IB}} = 1 + \gamma = 1 + 1/\beta$. Following
Kolchinsky \& Tracey (2019), we write the IB curve $F(r)$ (their Eq.~1), the IB
Lagrangian $L^{\mathrm{IB}}_\beta = I(Y;T) - \beta I(X;T)$ (their Eq.~2), the
DPI-saturating manifold $T_\alpha := B_\alpha \cdot Y$ (their Eq.~6), the bound
$L^{\mathrm{IB}}_\beta \le (1-\beta)H(Y)$ (their Eq.~7), and their footnote~4 (the
$\beta > 1$ regime of trivial maximizers). The variational objects are
\begin{align*}
R &= \mathbb{E}\big[\operatorname{KL}(q(z|x)\,\|\,m(z))\big] \ge I(X;Z)
&&\text{(VIB rate, marginal prior $m$; Fischer Eq.~19, Alemi et al.\ Eq.~14)} \\
\mathrm{ReX} &= \mathbb{E}\big[\operatorname{KL}(q(z|x)\,\|\,b_\theta(z|y))\big]
\ge I(X;Z|Y)
&&\text{(CEB residual; Fischer Eq.~20)} \\
\widehat I &= -\frac{1}{N}\sum_i \ln\Big[\frac{1}{N}\sum_j
e^{-\|f_i-f_j\|^2/(2\sigma^2)}\Big] \ge I(X;M)
&&\text{(NIB, non-parametric; Eq.~16, the homoscedastic case of Eq.~10)} \\
\mathrm{ReX}_F &= \mathbb{E}\big[\operatorname{KL}(q(z|x)\,\|\,r(z|y))\big]
\ge I(X;Z|Y)
&&\text{(FGIB, fixed anchor)}
\end{align*}
(Fischer 2020, Eqs.~18--20; Alemi et al.\ 2017, Eq.~14; Kolchinsky et al.\ 2017,
Eqs.~10 and 16). Fischer's comparison of the first two: at the respective prior optima
($m = q(z)$, $b = q(z|y)$), $\mathrm{ReX} = R - I(Y;Z)$ --- the conditional prior removes
exactly the retained label information, so CEB's bound is tighter than VIB's by
$I(Y;Z)$. We ask the same question one level down: where does FGIB's fixed-anchor bound
sit relative to CEB's learnable-backward bound?

\begin{proposition}[Channel P1, dead knob: every $\gamma > 0$ selects the same corner]\label{prop:dead}
Assume $Y = f(X)$. For every $\gamma > 0$,
\begin{equation}
\min_Z L_\gamma(Z) = -\gamma H(Y),
\qquad
L_\gamma(Z) := I(X;Z|Y) - \gamma I(Y;Z),
\label{eq:min}
\end{equation}
attained exactly on the corner family
$\{Z : I(Y;Z) = H(Y),\; I(X;Z|Y) = 0\}$ (e.g.\ $Z = Y$, or $Z = (Y, N)$ with
$N \perp X \mid Y$): the corner is unique in the information plane, while at the
representation level it is an entire family. Consequently the map
$\gamma \mapsto \big(I(X;Z^*|Y),\, I(Y;Z^*)\big)$ is constant: no interior point of the
IB curve is ever a Lagrangian minimizer, and CEB inherits the deterministic degeneration
of the Caveats paper (their Issue~1) over its \emph{entire} parameter range, since
$\beta_{\mathrm{IB}} = 1 + \gamma > 1$ always.
\end{proposition}

\begin{proof}
By Fischer's Eq.~4, $L_\gamma(Z) = I(X;Z) - (1+\gamma)I(Y;Z)$. The DPI
($I(Y;Z) \le I(X;Z)$, by the Markov chain $Y - X - Z$) gives
\begin{equation}
L_\gamma(Z) \ge I(Y;Z) - (1+\gamma)\,I(Y;Z)
= -\gamma\,I(Y;Z) \ge -\gamma H(Y),
\label{eq:chain}
\end{equation}
with equality iff both inequalities bind: $I(X;Z) = I(Y;Z)$ and $I(Y;Z) = H(Y)$.
$Z = Y$ attains the bound. Fischer (2020, p.~7) states the underlying equivalence
explicitly: CEB and IB are the same objective for $\gamma = \beta_{\mathrm{IB}} - 1$.
In the Caveats paper's maximizing convention, minimizing $L_\gamma$ is maximizing
$L^{\mathrm{IB}}$ with $\beta' = 1/\beta_{\mathrm{IB}} = 1/(1+\gamma) \in (0,1)$ ---
precisely the range in which their Eq.~7 pins every maximizer at the copy corner (their
Issue~1); their footnote~4 covers the complementary regime $\beta' > 1$ (i.e.\
$\gamma < 0$, outside CEB's range), whose maximizers are the trivial collapse. \hfill\qed
\end{proof}

\begin{remark}[The corner is the MNI point: why that does not rescue the knob]\label{rem:mni}
In the deterministic scenario the corner of Proposition~\ref{prop:dead} \emph{is} the MNI
point: $I(Y;Z) = H(Y) = I(X;Y)$ together with $I(X;Z|Y) = 0$ gives
$I(X;Y) = I(X;Z) = I(Y;Z)$, which is Fischer's Eq.~1. CEB's declared target and the
degenerate attractor therefore coincide, and Kolchinsky \& Tracey concede the point
themselves (their Section~4): ``if $Y = f(X)$ and one's goal is precisely to find the corner
point $\langle H(Y), H(Y)\rangle$ (maximum compression at no prediction loss), then our
results suggest that the IB Lagrangian provides a robust way to do so, being invariant to
the particular choice of $\beta$.'' Read that way, Proposition~\ref{prop:dead} states
hyperparameter robustness rather than failure. Three things keep it from closing the matter.
\begin{enumerate}
\item[(i)] \emph{The corner is one point in the information plane and an entire equivalence
class of representations.} The objective is exactly constant on
$\{Z : I(Y;Z) = H(Y),\ I(X;Z|Y) = 0\}$, which contains $Z = Y$, $Z = (Y,N)$ with
$N \perp X \mid Y$, and every bijective recoding of these. Adversarial robustness, OoD
detection and calibration --- the properties CEB is argued for --- are properties of the
\emph{geometry} of $Z$, on which the objective has no opinion; which member of the class is
reached is settled by the optimizer. That is channels P3 and P4 below.
\item[(ii)] \emph{CEB's own reading of the knob is unavailable at the exact optimum.}
Fischer (2020, p.~7) writes ``$\rho < 0$ may capture less than the MNI. $\rho > 0$ may
capture more'', whereas Proposition~\ref{prop:dead} places every $\gamma > 0$ at the MNI
point exactly. His argument for preferring $\gamma = 1$ is correspondingly not about
minimizers but about the optimizer: weighting the two appearances of $I(Y;Z)$ differently
``forces the optimization procedure to count bits of $I(Y;Z)$ in two different ways'' (his
Eq.~8 and surrounding discussion).
\item[(iii)] \emph{The knob does move trained models, which locates what it acts on.} On
Fashion MNIST --- a deterministically labelled dataset --- Fischer's own Figure~4 reports
the rate lower bound $R_X$ rising from $\approx 2.0$ to $\approx 5.0$ nats as $\rho$ goes
from $-1$ to $5$, with $\rho = 0$ landing on $I(X;Y) = \log 10 \approx 2.3$ nats. Since
every $\gamma > 0$ shares one exact optimum, that axis is not a path along the IB curve: it
is the distance between the optimum of the variational objective over a restricted family
and the point at which training was stopped. A knob acting through the approximation error
is precisely the knob that Propositions~\ref{prop:flat} and~\ref{prop:variance} show cannot
be read off the certificate.
\end{enumerate}
Kolchinsky \& Tracey state that their results ``are based on analytically-provable
properties of the IB curve, i.e., global optima'' and ``do not concern practical issues of
optimization''; Proposition~\ref{prop:flat} supplies the finite-$\beta$ variational
counterpart. The account is owed by any method reporting a $\beta$-sweep on
deterministically labelled data --- FGIB and the experiments of this paper included
(Remark~\ref{rem:c-scale}).
\end{remark}

\begin{proposition}[Channel P2, a blind certificate and a $1{:}\beta$ exchange rate]\label{prop:flat}
On the manifold $T_\alpha = B_\alpha \cdot Y$ of the Caveats paper (their Eq.~6),
$I(X;T_\alpha|Y) = 0$ for every $\alpha$, and with the optimal backward encoder and
classifier,
\begin{equation}
\mathrm{ReX}(T_\alpha) = 0,
\qquad
\mathrm{CE}(T_\alpha) = H(Y) - I(Y;T_\alpha) =: \delta_\alpha,
\qquad
\mathrm{VCEB}(T_\alpha) = \gamma\,\delta_\alpha .
\label{eq:vceb}
\end{equation}
Two consequences.
\emph{(a) The certificate is blind to label retention.} $\mathrm{ReX} = 0$ holds from the
total collapse ($\alpha = 0$, no label information) to the full copy ($\alpha = 1$), so its
tightness is compatible with retaining all, part, or none of $Y$. Fischer's reading of the
certificate as a meter --- ``observing ReX tells us exactly how many bits we are from the
MNI point'' (p.~7) --- therefore holds only along the $I(X;Z|Y)$ axis; on the necessity axis
the meter reads $0$ at chance level.
\emph{(b) Exchange rate.} With the optimal classifier the objective is
$L = \mathrm{CE} + \beta\,\mathrm{ReX} = \delta + \beta\,\mathrm{ReX}$, so it is indifferent
between giving up $\delta$ nats of label information and saving $\delta/\beta$ nats of
certified residual. In the heavy-regularization regime the bottleneck is supposed to act in,
the whole label axis --- up to $H(Y)$ nats --- is purchasable for $H(Y)/\beta$ nats of
certified compression.
\end{proposition}

\begin{proof}
$T_\alpha$ is a function of $Y$ alone, so $T_\alpha \perp X \mid Y$ and
$I(X;T_\alpha|Y) = 0$; the Caveats paper shows $I(Y;T_\alpha)$ sweeps $[0, H(Y)]$.
With the true backward encoder ($b = p(z|y)$) the certificate is tight:
$\mathrm{ReX} = I(X;T_\alpha|Y) = 0$ (Fischer Eqs.~9--11, 20). The optimal classifier
achieves $\langle \log c(y|z)\rangle = -H(Y|T_\alpha)$ (Fischer Eq.~12), i.e.\
$\mathrm{CE} = H(Y|T_\alpha) = \delta_\alpha$, and the VCEB value follows from Eq.~15.
Part~(b) is the definition of the objective evaluated at $\mathrm{CE} = \delta$.
\hfill\qed
\end{proof}

\begin{remark}[What the flatness claim is, and what it is not]\label{rem:convention}
An earlier form of this argument read the near-optimality band off the \emph{value}
$\mathrm{VCEB} = \gamma\delta \to 0$. That reading is a scaling artifact:
$\mathrm{VCEB} = \gamma\,(\mathrm{CE} + \beta\,\mathrm{ReX}) = \gamma L$, so the two
conventions differ by a positive factor and share their minimizers, and on $T_\alpha$ itself
the objective is $L = \delta$ in either convention and strictly prefers $\alpha = 1$. What
is convention-free is the \emph{ratio} of Proposition~\ref{prop:flat}(b): the prediction
term is weaker than the compression term by exactly the factor $\beta$, so in units of the
compression term the entire manifold lies within $H(Y)/\beta$ of optimal. Under a
scale-invariant optimizer (Adam, used throughout the experiments) only the ratio matters;
under plain SGD at fixed learning rate the convention additionally rescales the effective
step size. Combined with Proposition~\ref{prop:flat}(a) --- no certified compression is
bought anywhere on the manifold --- the operative statement is that the label axis is cheap
and unmetered, not that the objective is literally flat along it.
Pushing $\beta$ up also drives $\beta_{\mathrm{IB}} = 1 + 1/\beta \to 1^{+}$: the objective
approaches the Caveats paper's $\beta = 1$ degeneracy (their Eq.~7 discussion), at which
\emph{every} point of the diagonal is simultaneously optimal, from above. Fischer (2020,
p.~7) anticipates the direction --- ``$\rho < 0$ may capture less than the MNI'' --- but by
Proposition~\ref{prop:dead} that cannot happen at the exact optimum;
Proposition~\ref{prop:flat} gives the rate at which the variational objective is willing to
let it happen.
\end{remark}

\begin{remark}[Approximately deterministic labels]\label{rem:approx}
The Caveats paper (Appendix~C, Issue~1) shows that when $Y$ is only $\epsilon$-close to
deterministic, all optimizers fall within $O(-\epsilon\log\epsilon)$ of a single corner:
the dead-knob statement of Proposition~\ref{prop:dead} degrades gracefully, and the
near-optimality band of Proposition~\ref{prop:flat} is perturbed by the same order.
\end{remark}

\begin{proposition}[Channel P3, collusion collapse: the variance race and the overfit channel]\label{prop:variance}
For the diagonal Gaussian family $e(z|x) = \mathcal{N}(m(x), \sigma^2 I)$,
$b(z|y) = \mathcal{N}(g(y), \tau^2 I)$ (the implementation's setup, with the
\texttt{logvar} clamp $[-10, 10]$ of \texttt{kl\_divergence}), and a \emph{linear}
classifier $c(z) = Wz + b$ (the implementation's classifier is \texttt{nn.Linear}):
\begin{enumerate}
\item[(i)] for fixed means, $\operatorname{KL}$ is minimized over $\sigma^2$ at
$\sigma^2 = \tau^2$, with value $\tfrac{1}{2\tau^2}\|m - g\|^2$ --- a pure mean-matching
term;
\item[(ii)] the classifier term favors $\sigma^2 \to 0$:
$\mathrm{CE}(c(z), y) = -\log \mathrm{softmax}_y(Wz+b)$ is convex in $z$ (log-sum-exp
of an affine function of $z$, minus a linear term --- the convexity claim requires $c$
affine), so Jensen gives
$\mathbb{E}_z[\mathrm{CE}(c(z), y)] \ge \mathrm{CE}(c(m), y)$. The KL's derivative in
$\sigma^2$ vanishes at $\sigma^2 = \tau^2$, but along the matched line it \emph{resists}
shrinking $\tau^2$ while the means are unmatched
($\partial \operatorname{KL}/\partial \tau^2|_{\sigma^2=\tau^2}
= -\|m-g\|^2/(2\tau^4) < 0$). The joint dynamics therefore runs in two stages: the mean
matching $m \to g$ completes first (with weight $1/\tau^2$, which grows as $\tau$
shrinks), and once $\|m-g\| \to 0$ the variance race is unopposed --- the classifier
drives $\sigma^2 \to 0$ with $\tau^2$ following the moment-matched value --- until the
clamp at $s = e^{-10}$, where the encoder collapses toward the deterministic code
$z \approx g(y)$ and
$\mathrm{ReX} = \tfrac{1}{2s}\mathbb{E}\|m(x) - g(y)\|^2$: a mean-matching force of
weight $1/s = e^{10} \approx 2.2\times10^{4}$ between two \emph{trainable} maps;
\item[(iii)] the certified object and the deployed object separate. By the gap lemma
(Lemma~\ref{lem:gap}) $\mathrm{ReX} = I(X;Z|Y) + \Delta \ge I(X;Z|Y)$, so a small
$\mathrm{ReX}$ \emph{entails} a small residual for the stochastic training code --- the
certificate cannot under-report its own object. What it does not control is the relation
between the training code and the deployed representation. Exact mean matching
($m(x) = g(y)$) makes both objects class-constant, so both residuals vanish; the
scenario of interest is the \emph{near}-match regime, where
$\mathbb{E}\,\|m(X) - g(Y)\|^2 = o(s)$ keeps $\mathrm{ReX}$ small while the deployed
continuous $m$ retains within-class variation, and any non-class-constant continuous
$m$ has $I(X;m(X)|Y) = +\infty$ under the standing non-atomic assumption (the
pathology of deterministic networks, Amjad \&
Geiger 2018; Kolchinsky et al.\ 2017). The variance race therefore moves the
\emph{certified} object toward a noise-dominated class-constant code while the
\emph{deployed} object can keep diverging residual --- and the encoder is fitted to the
backward encoder's per-class targets, memorizing the label code $g(y)$ rather than
being told what residual to erase, while the certificate reports $\mathrm{ReX} \approx
0$ on a code that inference does not use.
\end{enumerate}
\end{proposition}

\begin{proof}
The closed form
$\operatorname{KL} = \tfrac{d}{2}\big[\sigma^2/\tau^2 + \|m-g\|^2/(d\tau^2)
- 1 + \ln(\tau^2/\sigma^2)\big]$ has derivative $\tfrac{d}{2}(1/\tau^2 - 1/\sigma^2)$
in $\sigma^2$, so the variance optimum is $\sigma^2 = \tau^2$, giving (i); at
$\sigma^2 = \tau^2 = s$ the value is $\|m-g\|^2/(2s)$. For (ii): $\mathrm{CE}(c(z), y)$ has Hessian
$\operatorname{diag}(p) - pp^{\top} \succeq 0$ (the categorical Fisher information), so
it is convex in $z$ and Jensen gives the stated inequality; the same Hessian identity makes
$\sigma^2 \mapsto \mathbb{E}_z[\mathrm{CE}(c(z),y)]$ non-decreasing (Gaussians of larger
variance dominate in convex order), strictly increasing in $\sigma^2$ wherever the
classifier is unsaturated, so that term pushes $\sigma^2$ down. The KL's partial derivatives at the
matched point are $\partial_{\sigma^2}\operatorname{KL} = 0$ and
$\partial_{\tau^2}\operatorname{KL} = -\|m-g\|^2/(2\tau^4) < 0$, so the KL pulls
$\tau^2$ up until the means match, while the mean-matching gradient on the encoder is
$(m - g)/\tau^2$; once $\|m-g\| \to 0$ the race to the clamp is unopposed. For (iii): the
gap identity (Lemma~\ref{lem:gap}) forces the certified residual and the certificate to
vanish together; the divergence attaches to the deployed mean $m(x)$ whenever it is not
class-constant (Amjad \& Geiger; Kolchinsky et al.). \hfill\qed
\end{proof}

\begin{remark}[Status of the two-stage account]\label{rem:p3-status}
Parts~(ii) and~(iii) are a stationarity-and-direction-of-force argument, not a convergence
proof: the ``two stages'' presume a time-scale separation (mean matching completes before
the variance race) that is not established here. The account is directly falsifiable in
this codebase: log the mean posterior log-variance and the fraction of units at the clamp
$s = e^{-10}$ during CEB training. \emph{We ran the test.} On MNIST, across all six
$\beta$ settings and 5 seeds (Section~\ref{sec:ablation} of the main paper), the trained
posterior and prior log-variances of CEB remain near their initialization ($\sigma^2 =
\tau^2 = 1$, mean $\log\sigma^2$ within $0.74$ of zero, clamp fraction $0.0$ at every
epoch of every run). At the early-stopping time scale of the experimental protocol, the
variance race therefore does not materialize, and
Proposition~\ref{prop:variance}(iii) is empirically void for those runs: Part~(iii)
stands as a sharp hypothesis about what longer training or different schedules would
produce, the force analysis of (i)--(ii) as the only proved content, and the
measurement itself as the finding.
\end{remark}

\begin{remark}[Train--evaluation mismatch]\label{rem:mismatch}
The classifier is trained on the sampled code $z$ and evaluated on the mean $m(x)$
(Fischer 2020, Appendix~A.1). Proposition~\ref{prop:variance}(ii) is the force that
closes this gap by destroying the noise altogether ($\sigma \to 0$), at which point the
bottleneck's stochasticity --- the object the bound is defined on --- disappears.
\end{remark}

\begin{proposition}[Channel P4, certificate-object mismatch: neither bound is an upper bound on the deployed representation]\label{prop:object}
(a)~\emph{CEB}. The classifier consumes the sampled code $Z$ at training and the mean
$\mu(X)$ at evaluation (Fischer 2020, Appendix~A.1). At the variance-optimal posterior,
$\sigma(x) \equiv \sqrt{S_\theta(y)}$ (the backward encoder's per-class standard
deviation) is class-constant, so $Z = \mu(X) + \sqrt{S_\theta(Y)}\,\varepsilon$ is a
noisy function of $\mu(X)$ given $Y$; by the DPI,
\begin{equation}
I(X; Z | Y) \;\le\; I(X;\, \mu(X) | Y) ,
\label{eq:dpi}
\end{equation}
strictly so whenever the noise destroys within-class detail. Hence CEB's certificate
$\mathrm{ReX} \ge I(X;Z|Y)$ is a bound on the \emph{training code}, not on the residual
of the representation deployed at inference. The gap is not a small correction: whenever
$\mu$ is a continuous map that is not class-constant, $I(X;\mu(X)|Y) = +\infty$ under
the standing non-atomic assumption (Amjad \& Geiger 2020; Kolchinsky et al.\ 2017), so
no finite quantity --- CEB's certificate included --- bounds the deployed residual from
above. The CE term presses the noise
toward $0$ (by Jensen, $\mathbb{E}_z[\mathrm{CE}(c(z),y)] \ge \mathrm{CE}(c(\mu),y)$),
but no finite $\sigma$ closes the gap.
(b)~\emph{FGIB}. The inference path is $h$ at both training and evaluation; the side
noise of Proposition~\ref{prop:inverse} is not part of any inference path. By the DPI
applied to the same perturbed-variable construction,
\begin{equation}
I(X;\, h + \tau (A^{\top}A)^{-1/2}\varepsilon \,|\, Y)
\;\le\; I(X; H | Y) ,
\label{eq:dpi-f}
\end{equation}
so the direction of the bound is the same as CEB's: $\mathrm{ReX}_F$ upper-bounds the
noisy side-path object, and \emph{no} finite value upper-bounds the deployed residual.
The point-anchor limit $\tau \to 0$ does not repair this: it is a change of model whose
bound value $\|\mu - a v_y\|^2/(2\tau^2)$ typically diverges
(Remark~\ref{rem:limit}), and even the limiting object needs $A$ well-conditioned
(Proposition~\ref{prop:conditioning}(iii)).
(c)~\emph{Conclusion}. In the pure variational value sense (same encoder, same Gaussian
family), CEB's bound is tighter (Proposition~\ref{prop:ordering}). In the object sense
the two are not ordered: both are valid upper bounds on their respective noisy codes, and
neither bounds the deterministic deployed representation. What remains structural --- not
a tightness claim --- is \emph{where the gap lives}: CEB's perturbation is inside the
inference path and is driven by the prediction term (so removing it changes the model),
while FGIB's perturbation lives on a side path that inference never reads and is
controlled by side-path parameters alone. The latter is a cleaner separation of
regularizer from model, not a tighter bound.
\end{proposition}

\begin{remark}[The bound is a training instrument, not an inference one]\label{rem:training-role}
No method computes its bound at inference: the bound's roles are \emph{training-time}
--- a regularizer (a differentiable surrogate for the intractable information term) and
a meter (Fischer 2020, p.~7: ``observing ReX tells us exactly how many bits we are from
the MNI point''). In the $Z \perp X \mid Y$ limit CEB's mean is also class-constant
($m(x) = \mathbb{E}[Z|x] = \mathbb{E}[Z|y]$), so the mismatch of
Proposition~\ref{prop:object}(a) is not a limit failure: it is a finite-$\beta$
\emph{under-reporting} --- the regularizer controls, and the meter reports, the training
code, which is strictly more compressed than the deployed representation $\mu(X)$, by a
positive, trained deficit that cannot be measured from the certificate. FGIB's
side-path noise is a known side-path parameter and is removable in principle
($\tau \to 0$ changes the side path, never the deployed model), but for finite
$\tau$ its regularizer and meter act on a \emph{perturbed copy} of the deployed
object, with the perturbation size governed by the conditioning of $A$
(Proposition~\ref{prop:conditioning}(iii)).
The one genuine inference-time use of a certificate value is as a detection score
(Fischer 2020, RG3, uses the marginal rate $R$): CEB's label-conditioned
$\mathrm{ReX}$ is not even computable at inference (the label is unknown --- Fischer
trains a separate marginal for this purpose), while FGIB's geometry gives a label-free
score, the distance to the nearest anchor
$\min_y \operatorname{KL}(q(z|x)\|\mathcal{N}(a v_y, \tau^2 I))$, with no extra network
(the minimum is an internal enumeration over the $K$ fixed anchor buffers --- the single
QR orthonormal frame --- and consults no ground-truth label, no more than softmax
consults labels when it enumerates the $K$ logits). Whether this score is useful for OoD
detection is a hypothesis; nothing here proves it.
\end{remark}

\textbf{The problem addressed by FGIB.}
The four channels are addressed by the FGIB design as follows; two are closed by
construction, one holds conditionally, and one is only partially addressed. The precise
scope of each claim is documented in Part~II and the open items are collected in
Section~\ref{sec:limits}.
\begin{itemize}
\item \emph{Dead knob and cheap label axis} (Props.~\ref{prop:dead}, \ref{prop:flat}):
FGIB's compression term is the anchor gap $\tfrac{1}{2\tau^2}(s-a)^2$, quadratic in the
alignment scalar, and \emph{at fixed classifier scale} the Lagrangian
$\mathrm{CE}(s) + \tfrac{\beta}{2\tau^2}(s-a)^2$ is strictly convex with a unique interior
minimizer per $\beta$ --- the $\beta$-sweep is one-to-one onto the whole trade-off curve
(Props.~\ref{prop:sweep}, \ref{prop:curve}). With the scale free the infimum is $0$ for
every $\beta$; Remark~\ref{rem:c-scale} states the condition and the architectural fix.
\item \emph{Blind certificate} (Prop.~\ref{prop:flat}): with the fixed anchor
$r(z|y) = \mathcal{N}(a v_y, \tau^2 I)$ the gap
$\mathbb{E}_y[\operatorname{KL}(q(z|y)\,\|\,r(z|y))]$ can shrink only if the
\emph{posterior} moves toward the anchors --- the reference itself cannot chase the
encoder, so the force is persistent (Prop.~\ref{prop:ordering} and
Remark~\ref{rem:force}). The side head $A$, however, is a second trainable party in that
matching, and it is not part of the inference path
(Remark~\ref{rem:A-channel}).
\item \emph{Variance collapse} (Prop.~\ref{prop:variance}): the classifier consumes the
deterministic $h$ at both training and evaluation, so $\sigma^{2*} = \tau^2$ for every $\beta$
(the variance decouples from the prediction term; Lemma~\ref{lem:lagrangian}), the
side path keeps a finite covariance and the bound stays well-defined
(Remark~\ref{rem:limit}), and the trivial collapse $\mu \equiv 0$ pays
the fixed positive price $\operatorname{KL} = \tfrac12 a^2 > 0$ for every class --- the
degenerate corner is excluded by geometry rather than by the optimizer's good behavior.
This channel is closed cleanly, by construction.
\item \emph{Certificate-object mismatch} (Prop.~\ref{prop:object}): the classifier
consumes the deterministic $h$ at both training and evaluation --- there is no
training-code/evaluation-mean split --- but the bound is an upper bound on a
\emph{noisy side-path copy} of $h$, not on $I(X;H|Y)$, which no finite certificate
bounds. This channel is \emph{not} closed; what FGIB buys is that the perturbation is
side-path-only, and its size is governed by the conditioning of $A$
(Prop.~\ref{prop:conditioning}(iii)).
\end{itemize}

% =====================================================================
\section*{Part II. The FGIB theory}
% =====================================================================

\textbf{FGIB notation.}
The deterministic main path computes $h = f_{\mathrm{enc}}(X)$ and the classifier
consumes $h$ directly; the side path parameterizes $q(z|x) = \mathcal{N}(\mu(x),
\Sigma(x))$ with \emph{linear} mean head $\mu(x) = A\,h(x)$ for
$A \in \mathbb{R}^{z \times z}$ (the theory assumes $z = d$ and $A$ invertible; the
inverse transforms of Propositions~\ref{prop:fgib-ceb} and~\ref{prop:inverse} are
defined under this assumption), and a separate logvar head for $\Sigma(x)$. The
side path is \emph{not part of inference}: at evaluation the model returns the classifier
output and never evaluates the mean head (so $A$ is a trainable module that exists only
inside the regularizer --- structurally the same position CEB's backward encoder
occupies on the other side of the KL). The label-conditional prior is the
\emph{fixed} anchor
$r(z|y) = \mathcal{N}(a\,v_y, \tau^2 I)$ with \emph{anchor variance} $\tau^2 > 0$
(default $1.0$): the $v_y$ are orthonormal directions obtained by
QR-decomposing a random matrix ($K \le z$; the orthonormal factor of the random matrix),
scaled by the anchor scale $a$ (default $4.0$), stored as
non-trainable buffers. At the variance optimum of the KL the posterior variance matches
the anchor variance, $\sigma^{2*} = \tau^2$; we write $\sigma_p^2$ for this common
variance where the distinction does not matter. The orthogonality of the method lies in
this fixed frame and in the equilibrium it induces --- the KL drives the
class-conditional means $\mu_y = A\,\mathbb{E}[h|y]$ toward the orthonormal anchors
$a v_y$, so the transformed class representations become mutually orthogonal (the
aligned family below) --- while $A$ itself remains an unconstrained linear map. For
regression, fixed random-Fourier anchors
$\mu_p(y) = a\sqrt{2/z}\,\cos(2\pi\omega y + b)$ with frozen
$\omega \sim \mathcal{N}(0, 2^2)$, $b \sim U[0, 2\pi)$ replace the per-class table
($\|\mu_p(y)\| \approx a$). The loss is
$L_{\mathrm{FGIB}} = \mathbb{E}[\mathrm{CE}(c(h), y)]
+ \beta\,\mathbb{E}[\operatorname{KL}(q(z|x)\,\|\,r(z|y))]$, with gradients reaching the
main path only through the shared encoder; the zero-initialized logvar head fixes
$\sigma^2 = 1$ at initialization, which equals $\tau^2$ in the default configuration
$\tau^2 = 1$ (for general $\tau$ the theory's variance optimum $\sigma^2 = \tau^2$ is
reached by the logvar head, whose initialization is a scheduling choice), so training
starts at $\operatorname{KL} \approx \tfrac{1}{2\tau^2} a^2$.

\emph{Implementation caveat (theory--code layering).} The theory works with a fixed
orthonormal factor, an invertible linear side head $T(h) = Ah$, and per-class Gaussian
priors with anchor variance $\tau^2$. The code implements an \emph{affine} head
(\texttt{nn.Linear} with bias) --- which the theory absorbs as an idealization, the bias
being immaterial when $A$ is learned jointly --- and does not enforce orthogonality,
invertibility, or conditioning constraints on $A$. The guarantees below (orthonormal
frames, invariance, conditioning) are properties of the idealized model; the
implementation approximates them, and the guarantees do not transfer automatically
(Remarks~\ref{rem:A-channel}, \ref{rem:rank}).
\textbf{Distributional standing assumption.} Several results state that a deterministic
continuous representation carries infinite conditional information. Those statements
assume the standard continuous setting: $X \mid Y$ is non-atomic, and the induced
representation $H(X) \mid Y$ is a non-atomic continuous variable. For the finite
empirical distribution of the training data no such conclusion is claimed; there the
quantities are finite and only their growth with sample size is governed by these
results.

% ---------------------------------------------------------------
\subsection{The equivalence: the side path is a CEB bottleneck on the main path}
% ---------------------------------------------------------------

\textbf{Motivation.}
Stochastic bottlenecks such as VIB and CEB are \emph{stochastic at training time but
deterministic at evaluation time}: the classifier is trained on sampled codes
$z \sim q(z|x)$ yet evaluated on the mean $\mu_q(x)$ (Fischer 2020, Appendix A.1), so
its training and evaluation distributions differ. In the same models the cross-entropy
gradient and the KL gradient act on the same heads: the classification term can exploit
the variance (by Jensen, $\mathbb{E}_z[\mathrm{CE}(c(z), y)] \ge \mathrm{CE}(c(\mu_q),
y)$, so loosening the noise relaxes the upper bound), while the KL term pushes the
posterior back toward the prior. Finally, the regularizer constrains the noisy code $z$,
whereas inference consumes the noise-free mean --- the object being regularized is not
the object being used. FGIB resolves all three mismatches by design: the classifier
consumes the deterministic representation $h$ at both training and evaluation time, the
two objectives act on disjoint parameter sets (the classification gradient does not
enter the side heads, and the KL gradient reaches the main path only through the shared
encoder), and, by Proposition~\ref{prop:fgib-ceb} below, the side-path regularizer is,
\emph{at each fixed value of the side head}, a CEB bottleneck read in the coordinates in
which the posterior is centered at the representation that inference actually uses
(Remark~\ref{rem:A-channel} states why this is weaker than a fixed constraint on $h$).
Confining the stochasticity to a side path also keeps the variational objective
well-defined: for continuous deterministic representations the residual information
diverges (Amjad \& Geiger; Kolchinsky et al.), so a purely deterministic bottleneck
admits no finite variational bound, while the finite side-path covariance of FGIB does
(Remark~\ref{rem:limit}). Whether the robustness benefits of the stochastic bottleneck
survive the deterministic inference path is the empirical question studied in the
experiments.

\begin{lemma}[Linear invariance of the Gaussian KL]\label{lem:kl-invariance}
Let $q = \mathcal{N}(\mu_q, \Sigma_q)$ and $p = \mathcal{N}(\mu_p, \Sigma_p)$ on
$\mathbb{R}^{z}$, and let $A \in \mathbb{R}^{z \times z}$ be invertible. Define the
common pushforwards $q' = A_\# q = \mathcal{N}(A\mu_q, A\Sigma_q A^\top)$ and
$p' = A_\# p$. Then $\operatorname{KL}(q'\,\|\,p') = \operatorname{KL}(q\,\|\,p)$.
\end{lemma}

\begin{proof}
With $d = z$, the closed form reads
$\operatorname{KL}(q\,\|\,p)
= \tfrac{1}{2}\big[\operatorname{tr}(\Sigma_p^{-1}\Sigma_q)
+ (\mu_q - \mu_p)^\top \Sigma_p^{-1}(\mu_q - \mu_p)
- d + \ln|\Sigma_p| - \ln|\Sigma_q|\big]$.
Under the common map $A$ the three non-constant terms are unchanged:
\begin{align*}
\operatorname{tr}\big((A\Sigma_p A^\top)^{-1}(A\Sigma_q A^\top)\big)
&= \operatorname{tr}(A^{-\top}\Sigma_p^{-1}A^{-1} A\,\Sigma_q A^\top)
 = \operatorname{tr}(A^{-\top}\Sigma_p^{-1}\Sigma_q A^\top) \\
&= \operatorname{tr}(\Sigma_p^{-1}\Sigma_q A^\top A^{-\top})
 = \operatorname{tr}(\Sigma_p^{-1}\Sigma_q), \\[2pt]
(A\Delta\mu)^\top (A\Sigma_p A^\top)^{-1} (A\Delta\mu)
&= \Delta\mu^\top A^\top A^{-\top}\Sigma_p^{-1}A^{-1}A\,\Delta\mu
 = \Delta\mu^\top \Sigma_p^{-1}\Delta\mu, \\[2pt]
\ln|A\Sigma_p A^\top| - \ln|A\Sigma_q A^\top|
&= \big(\ln|\Sigma_p| + 2\ln|A|\big) - \big(\ln|\Sigma_q| + 2\ln|A|\big)
 = \ln|\Sigma_p| - \ln|\Sigma_q|,
\end{align*}
where $\Delta\mu = \mu_q - \mu_p$ and the trace step uses the cyclic property. \hfill\qed
\end{proof}

\begin{lemma}[Invariance of the CEB objective and the MNI point]\label{lem:mi-invariance}
For any bijective measurable map $\phi$,
$I(X;\phi(Z)|Y) = I(X;Z|Y)$ and $I(Y;\phi(Z)) = I(Y;Z)$.
Consequently the CEB objective (Eq.~5) is invariant under bijections of $Z$, and
$Z$ attains the MNI point (Eq.~1) if and only if $\phi(Z)$ does.
\end{lemma}

\begin{proof}
Mutual information is invariant under bijective transformations of either argument
(data processing inequality applied in both directions). \hfill\qed
\end{proof}

\begin{proposition}[The FGIB side path is a CEB bottleneck on the main path, at each fixed side head]\label{prop:fgib-ceb}
Let $z = d$ and let the learned linear head $T(h) = A h$ be invertible (unconstrained,
so a.s.\ so at initialization). The \emph{fixed}
prior $r(z|y) = \mathcal{N}(\mu_p(y), \Sigma_p(y))$ --- for FGIB, the orthonormal anchor
table $\mu_p(y) = a\,v_y$, $\Sigma_p = \tau^2 I$ --- is, by construction, the image
under $T$ of the canonical prior
$r'(\tilde{z}|y) = \mathcal{N}(A^{-1}\mu_p(y),\, A^{-1}\Sigma_p(y) A^{-\top})$
(since $A A^{-1} = I$), and the posterior $q(z|x) = \mathcal{N}(A\,h(x), \Sigma_q(x))$
is the image of
$q'(\tilde{z}|x) = \mathcal{N}(h(x), A^{-1}\Sigma_q(x) A^{-\top})$.
\emph{At a fixed value of $A$}, the prior is never transformed by training:
$A^{-1}$ only reads the fixed anchor geometry in the coordinates in which the posterior
is centered at $h$; during training $A$ moves, so the pulled-back prior moves with
it (Remark~\ref{rem:A-channel}).
By KL invariance under the common invertible linear map $T^{-1} = A^{-1}$
(Lemma~\ref{lem:kl-invariance}), pointwise for every $(x, y)$,
\begin{equation}
\operatorname{KL}\big(q(z|x)\,\|\,r(z|y)\big)
= \operatorname{KL}\big(q'(\tilde{z}|x)\,\|\,r'(\tilde{z}|y)\big),
\label{eq:fgib-ceb}
\end{equation}
That is, at each fixed side head the side-path regularizer is pointwise identical to a
CEB bottleneck with a backward encoder \emph{frozen at the current head}, applied to
a noisy version of the main-path representation,
$\tilde{z} = h + \tilde{\epsilon}$ with
$\tilde{\epsilon} \sim \mathcal{N}(0, A^{-1}\Sigma_q A^{-\top})$,
using the canonical prior $r'$ as that backward encoder. For FGIB the canonical backward
encoder is the frame $\mathcal{N}(a\,A^{-1} v_y, \tau^2 A^{-1}A^{-\top})$, whose
mean anchors $A^{-1}(a v_y)$ are orthonormal at scale $a$ in the induced metric
$A^{\top}A$
($\langle A^{-1}v_i, A^{-1}v_j\rangle_{A^{\top}A} = \delta_{ij}$): the orthogonality
enters through the fixed frame $\{v_y\}$ and the equilibrium it induces, not through a
constraint on $A$.
Hence the regularizer's gradient with respect to $h$ (through the shared encoder)
coincides, up to the fixed linear map $A^{\top}$, with the gradient that a CEB model
applies to its representation mean.
\end{proposition}

\begin{proof}
The map $z \mapsto \tilde{z} = A^{-1}z$ is a bijection, and
$q' = (A^{-1})_\# q$, $r' = (A^{-1})_\# r$:
$A^{-1}A\,h = h$ and $A^{-1}\Sigma_q A^{-\top}$ give the canonical posterior, and
$A^{-1}\mu_p$, $A^{-1}\Sigma_p A^{-\top}$ give the canonical prior, with
$A\,(A^{-1}\mu_p) = \mu_p$ and $A\,(A^{-1}\Sigma_p A^{-\top})\,A^{\top}
= \Sigma_p$, so $r$ is exactly the image of $r'$ under $T$.
Equation~\eqref{eq:fgib-ceb} is Lemma~\ref{lem:kl-invariance} applied to $A^{-1}$.
Sampling $\tilde{z} \sim q'$ is equivalent to
$\tilde{z} = h + \tilde{\epsilon}$ with the stated noise.
Combined with the variational bound of Fischer (2020, Eqs.~9--11), the regularizer
upper-bounds the residual information of the noisy main-path representation:
$I(X;\tilde{Z}|Y) \le \mathbb{E}[\operatorname{KL}(q'(\tilde{z}|x)\,\|\,r'(\tilde{z}|y))]$.
\hfill\qed
\end{proof}

\begin{remark}[Trainable-implementation caveat: the side head is a second trainable party in the matching]\label{rem:A-channel}
The identity of Proposition~\ref{prop:fgib-ceb} holds pointwise at each value of the
side head, so the asymmetry advertised against CEB --- one trainable side versus two ---
is weaker than it looks. In the idealized theory $A$ may be treated as fixed
(orthonormal or any fixed invertible map); in the \emph{trainable implementation} $A$ is
free and \emph{unused at inference}:
like CEB's backward encoder $b(z|y)$, it lives inside the KL and can lower it without
moving the deployed representation $h$. Since $h$ is linear-separable by the jointly
trained classifier, the linear fit of the anchor target $a v_y$ on $h$ fits well, and
much of that fit can be absorbed by $A$ alone. The variance term
$\tfrac{z}{2}(\sigma^2/\tau^2 - 1 - \ln(\sigma^2/\tau^2))$ does not depend on $A$ at
all, so nothing in
the regularizer resists $A$ drifting toward a projector onto the $K$-dimensional
between-class subspace with a vanishing component on within-class directions. Full
collapse is blocked --- with invertible $A$, $\operatorname{KL} = 0$ forces $h$ itself
to be class-constant --- but the approach to the rank-$K$ regime is not blocked, and
that regime is exactly where the conditioning failure and the perturbation blow-up of
Propositions~\ref{prop:conditioning} and~\ref{prop:inverse} live. This is the one
channel of Part~I that FGIB has not closed; Section~\ref{sec:limits} collects it.
\end{remark}

\begin{remark}[Family closure]\label{rem:closure}
Lemma~\ref{lem:kl-invariance} holds for the underlying Gaussian distributions. Within a
restricted variational family the equality of Proposition~\ref{prop:fgib-ceb} is exact
only if the family is closed under the inverse map $A^{-1}$: the full-covariance family
is closed (and is the parameterization used by Fischer 2020); the diagonal family is not
--- for diagonal $q, r$ the transformed pair is diagonal only if $A^{-1}$ preserves
diagonality (e.g., permutations and sign flips) and approximate otherwise. The diagonal
implementation of the codebase is therefore covered at the distribution level, with the
family approximation confined to the noise covariance $A^{-1}\Sigma_q A^{-\top}$
(Remark~\ref{rem:orth}).
\end{remark}

\begin{remark}[Rank deficiency]\label{rem:rank}
If $z < d$ or $A$ is rank-deficient, only the projection of $h$ onto the row space of
$A$ (the image of $A^{\top}$ in $h$-space) receives gradient from the regularizer; the
null-space component is unconstrained.
In the default configuration $z = d$ and a randomly initialized $A$ is almost surely
invertible, but invertibility carries no conditioning margin --- near-singularity is
enough to produce the effects of Proposition~\ref{prop:conditioning}(iii), and it is the
direction the gap gradient points in (Remark~\ref{rem:A-channel}). If $A \in
\mathbb{R}^{z \times d}$ has full column rank ($z \ge d$), a left inverse $B$ with
$BA = I$ exists and the data-processing inequality for KL gives, per $(x,y)$,
$\operatorname{KL}(q\|r) \ge \operatorname{KL}(B_\# q\,\|\,B_\# r)
\ge I(X;\,BZ\,|\,Y)$ --- a valid lower-level bound, but one whose equality
(invariance) is lost, and whose pulled-back prior again depends on $B$. The exact
invariance argument of Proposition~\ref{prop:fgib-ceb} genuinely requires the invertible
linear map $T(h) = Ah$ with $z = d$.
\end{remark}

\begin{remark}[Noise-free limit]\label{rem:limit}
The equivalence of Proposition~\ref{prop:fgib-ceb} should be understood at the level of
the variational objective, not as a mutual-information limit: as the side-path noise
vanishes, $I(X;Z|Y)$ diverges for continuous deterministic $H$ under the standing
non-atomic assumption (the pathology of
deterministic networks discussed by Amjad \& Geiger and Kolchinsky et al.), and the
variational KL bound becomes loose. Keeping the side-path covariance finite --- without
injecting noise into the inference path --- is precisely what keeps the regularizer
well-defined; this is the design rationale of FGIB.
\end{remark}

% ---------------------------------------------------------------
\subsection{The sweep}
% ---------------------------------------------------------------

\textbf{The degeneracy to be addressed (restated in our terms).}
When $Y = f(X)$, the IB curve saturates the data-processing bound and is piecewise
linear, $F(r) = \min\{r, H(Y)\}$ (Caveats paper, Section~3). Two consequences (their
Caveat~1): (i)~by the DPI, $L^{\mathrm{IB}}_\beta(T) = I(Y;T) - \beta I(X;T)
\le (1-\beta)I(Y;T) \le (1-\beta)H(Y)$, and $T_{\mathrm{copy}} := f(X) = Y$ attains the
bound for \emph{every} $\beta \in [0,1]$: all $\beta \in (0,1)$ are pinned at the single
corner point $\langle H(Y), H(Y)\rangle$ of the curve; (ii)~conversely, at $\beta = 1$
every point of the increasing part simultaneously maximizes $L^{\mathrm{IB}}_1$. The
mechanism is exposed by writing $\max_T L^{\mathrm{IB}}_\beta = \max_r [F(r) - \beta r]$
(the Legendre transform of $-F$): on the linear increasing part every point has the
\emph{same} derivative $F'(r) = 1$, hence the same supporting slope --- a Lagrangian of
the form ``prediction $-\beta\cdot{}$compression'' sweeps a Pareto curve only if each
point of the curve has a \emph{unique} supporting slope, i.e.\ the curve is strictly
concave. The dIB variant fails identically: on the hard-clustering family $T = g(Y)$ one
has $H(T|Y) = 0$, hence $I(Y;T) = H(T)$, so $L^{\mathrm{dIB}}_\beta = (1-\beta)H(T)$
--- a \emph{monotone} function of the single compression scalar $H(T)$, maximized by the
finest partition for every $\beta < 1$. The squared fixes repair this by making the
objective strictly concave in the compression scalar: $F'(r)/(2r) = \beta$ is strictly
decreasing, so each $\beta$ selects a unique $r$ (Eq.~8 argument), and squared-dIB has
the unique maximizer $H(T) = 1/(2\beta)$ (their Eq.~A20). By Proposition~\ref{prop:dead}
this degeneracy is inherited verbatim by CEB ($\beta_{\mathrm{IB}} = 1+\gamma > 1$),
which is the first channel of Part~I.

\textbf{Reduced model of FGIB.}
The KL depends on the side mean $\mu$ alone, and the classifier term is invariant under
the coordinate change $\tilde w_j = A^{-\top} w_j$ (since $w_j \cdot h =
\tilde w_j \cdot \mu$): with $A$ invertible ($z = d$ in the default configuration, so a
random $A$ is a.s.\ invertible), every equilibrium in $(h, w)$ coordinates maps
bijectively to one in $(\mu, \tilde w)$ coordinates. We therefore analyze the reduced
model $\mu(x) = h(x)$ without loss of generality.
Consider the \emph{class-symmetric aligned family}: the encoder maps every input of
class $y$ to $h_y = s\,v_y$ ($s \ge 0$), the classifier is held at the symmetric
configuration $w_j = c\,v_j$, $b_j = 0$ with fixed scale $c > 0$, and the posterior is
$q_y = \mathcal{N}(s v_y, \sigma^2 I)$. The family is the natural equilibrium manifold of
the two forces: the KL gradient
$\nabla_h \operatorname{KL} = (\mu - a v_y)/\tau^2$ acts exactly
along $v_y$, and the classification force on the class center acts along $v_y$ at the
symmetric configuration; the ansatz is a \emph{working approximation}: it captures the
direction of both forces, and is motivated as exact in the $\beta$-dominant regime, but
that exactness is not proved here --- the cross terms are of order $O(c/\beta)$ and are
idealized away (see Remark~\ref{rem:caveats}).

\begin{lemma}[FGIB Lagrangian on the aligned family]\label{lem:lagrangian}
On the aligned family,
\begin{equation}
\operatorname{KL}(s, \sigma^2)
= \tfrac{z}{2}\big(\sigma^2/\tau^2 - 1 - \ln(\sigma^2/\tau^2)\big)
+ \tfrac{1}{2\tau^2}(s-a)^2,
\qquad
\mathrm{CE}(s) = \ln\!\big(1 + (K-1)e^{-cs}\big),
\label{eq:kl-ce}
\end{equation}
and the optimizer over $\sigma^2$ is $\sigma^2 = \tau^2$ for \emph{every} $\beta$, since
the CE term does not depend on $\sigma^2$ (the classifier consumes the deterministic
$h$, not $z$).
\end{lemma}

\begin{proof}
With $h_y = s v_y$ and $w_j = c v_j$, the logits are $c\,\langle v_j, s v_y\rangle
= cs\,\delta_{jy}$, so $p_y = e^{cs}/(e^{cs} + K - 1)$ and
$\mathrm{CE} = -\ln p_y = \ln(1 + (K-1)e^{-cs})$.
For the KL, the diagonal closed form with $\Sigma_q = \sigma^2 I$,
$\Sigma_p = \tau^2 I$ reads
$\tfrac{1}{2}\sum_d\big[\sigma^2/\tau^2 + (\mu_d - a v_{yd})^2/\tau^2 - 1
- \ln(\sigma^2/\tau^2)\big]$, giving
Eq.~\eqref{eq:kl-ce} since $\|v_y\| = 1$. The variance term is non-negative with a
unique minimum at $\sigma^2 = \tau^2$, and $\mathrm{CE}$ is independent of $\sigma^2$.
\hfill\qed
\end{proof}

\begin{proposition}[Unique interior optimum per $\beta$ at fixed classifier scale --- the sweep]\label{prop:sweep}
\emph{At a fixed classifier scale $c$}, for every $\beta > 0$, the FGIB Lagrangian on
the aligned family,
\begin{equation}
L_\beta(s) := \ln\!\big(1 + (K-1)e^{-cs}\big) + \tfrac{\beta}{2\tau^2}(s-a)^2,
\label{eq:fgib-lag}
\end{equation}
is strictly convex in $s$ and has a unique minimizer $s^*(\beta) > a$, given implicitly
by
\begin{equation}
\frac{\beta}{\tau^2}\,(s-a) = c\,\varphi(cs),
\qquad
\varphi(u) := \frac{(K-1)e^{-u}}{1 + (K-1)e^{-u}} \in (0,1) .
\label{eq:stationarity}
\end{equation}
The map $\beta \mapsto s^*(\beta)$ is continuous and strictly decreasing, with
$s^*(\beta) \to a$ as $\beta \to \infty$ and $s^*(\beta) \to \infty$ as $\beta \to 0$.
Consequently the map $\beta \mapsto \big(\operatorname{KL}(s^*), \mathrm{CE}(s^*)\big)$
is one-to-one onto the entire curve
$\{\big(\tfrac{1}{2\tau^2}(s-a)^2,\; \ln(1+(K-1)e^{-cs})\big) : s \in (a, \infty)\}$:
no two
$\beta$ share the same optimum, and every interior point of the trade-off curve is
attained --- the exact failure mode of $L^{\mathrm{IB}}_\beta$ (all $\beta$ pinned at
$T_{\mathrm{copy}}$) is absent. The fixed-$c$ condition is load-bearing and is analyzed
in Remark~\ref{rem:c-scale}.
\end{proposition}

\begin{proof}
Compute $\mathrm{CE}'(s) = -c\varphi(cs)$ and
$\mathrm{CE}''(s) = c^2 (K-1) e^{-cs}/(1 + (K-1)e^{-cs})^2 > 0$, so CE is strictly
convex and strictly decreasing in $s$; the compression term $\tfrac{1}{2\tau^2}(s-a)^2$
is strictly convex. Hence $L_\beta''(s) = \mathrm{CE}''(s) + \beta/\tau^2 > 0$:
$L_\beta$ is
strictly convex with at most one minimizer, characterized by
$0 = L_\beta'(s) = -c\varphi(cs) + \beta(s-a)/\tau^2$, i.e.\
Eq.~\eqref{eq:stationarity}.
The left side $\beta(s-a)/\tau^2$ is strictly increasing from $-\beta a/\tau^2$
(at $s=0$) through $0$
(at $s=a$) to $+\infty$; the right side $c\varphi(cs)$ is strictly decreasing in $s$
with $c\varphi(cs) \in (0, c)$. The unique crossing therefore satisfies $s^* > a$.
Differentiating Eq.~\eqref{eq:stationarity} with respect to $\beta$ gives
$(s-a)/\tau^2 + (\beta/\tau^2) s' = c^2\varphi'(cs)\,s'$ with $\varphi' < 0$, so
$s' = -(s-a)/(\beta - c^2\tau^2\varphi'(cs)) < 0$: $s^*$ is strictly decreasing.
The limits follow by monotonicity: as $\beta \to \infty$ the left side can only match
the bounded right side when $s \to a$; as $\beta \to 0$ the crossing requires
$\varphi(cs)\to 0$, i.e.\ $s \to \infty$. \hfill\qed
\end{proof}

\begin{proposition}[The trade-off curve is strictly concave in the compression--prediction plane]\label{prop:curve}
On the curve of Proposition~\ref{prop:sweep} (fixed classifier scale $c$),
\begin{equation}
\frac{d\,\mathrm{CE}}{d\,\operatorname{KL}}
= \frac{\mathrm{CE}'(s)}{(s-a)/\tau^2}
= -\frac{\tau^2\,c\varphi(cs)}{s-a},
\label{eq:slope}
\end{equation}
which is strictly increasing from $-\infty$ (as $s \to a^+$) to $0^-$ (as $s \to
\infty$). Thus the curve $\operatorname{KL} \mapsto -\mathrm{CE}$ (compression versus
achievable prediction, on the KL--CE surrogate plane of
Remark~\ref{rem:surrogate}) is strictly concave, and each of its points has a \emph{unique}
supporting slope $-\beta$. By the criterion of the Caveats paper (their Fig.~2 right
panel and the $F'(r)/(2r)$ argument of Eq.~8), this is precisely the property that makes
the Lagrangian $-\mathrm{CE} - \beta\,\operatorname{KL}$ sweep the full \emph{surrogate}
curve on the aligned family at fixed scale: the convexity of the compression term
supplies the strict-concavity condition \emph{without squaring it}. This is a statement
about the KL--CE surrogate of Proposition~\ref{prop:sweep}, not about the Caveats
information plane (Remark~\ref{rem:surrogate}).
\end{proposition}

\begin{proof}
$\varphi$ is positive and strictly decreasing and $s-a$ is positive and strictly
increasing, so $\tau^2 c\varphi(cs)/(s-a)$ is strictly decreasing in $s$; the signs give
the stated limits. The stationarity condition Eq.~\eqref{eq:stationarity} is exactly
$-d\mathrm{CE}/d\operatorname{KL} = \beta$, i.e.\ $\beta$ is the supporting slope at the
optimum. \hfill\qed
\end{proof}

\begin{remark}[What the sweep sweeps: a KL--CE surrogate curve, not the IB curve]\label{rem:surrogate}
Propositions~\ref{prop:sweep} and~\ref{prop:curve} prove strict convexity of the
FGIB Lagrangian and one-to-one parameterization by $\beta$ \emph{on the class-symmetric
aligned family, at fixed classifier scale}, in the plane
$\big(\operatorname{KL}_{\mathrm{anchor}},\, \mathrm{CE}\big)$. That plane is not the
Caveats paper's IB plane $\big(I(X;T),\, I(Y;T)\big)$, and the claim does not extend to
it. On the aligned family the posterior is class-constant by construction, so
$I(X;Z|Y) = 0$ for every $s$, while the deterministic class prototypes $h_y = s v_y$
already retain the label completely for any $s > 0$ ($I(Y;H) = H(Y)$); what varies
along the sweep is the anchor mismatch and the classifier margin, not a proved
compression of mutual information. Whether the deployed representation's residual
$I(X;H|Y)$ moves monotonically with $\beta$ is a separate, empirical question; the
theory guarantees a strictly concave surrogate, not a recovered IB curve. This
qualification is carried into Section~\ref{sec:limits}.
\end{remark}

\begin{remark}[The classifier scale is load-bearing, and the current implementation leaves it free]\label{rem:c-scale}
If $c$ is free, the pair $(s, c) = (a, c)$ with $c \to \infty$ sends
$\operatorname{KL} = \tfrac{1}{2\tau^2}(s-a)^2 = 0$ and
$\mathrm{CE} = \ln(1 + (K-1)e^{-ca}) \to 0$ simultaneously (for $a > 0$): the infimum
of $L_\beta$ over the aligned family is $0$ for \emph{every} $\beta$, no minimizer is
attained, and the one-to-one sweep of Proposition~\ref{prop:sweep} has no global
counterpart. (At $a = 0$ the margin growth alone cannot drive CE to zero.) This is
the FGIB analogue of the dead knob --- the same one-corner attractor, now arrived at by
growing the margin rather than by the linearity of the deterministic IB curve.
Three consequences.
\emph{(a) The present implementation is exposed.} The classifier is
\texttt{nn.Linear} trained by Adam with no weight decay, so $c$ is unconstrained. On
separable data the CE gradient drives $\|w_y\|$ upward; it does not diverge in finite
time (the gradient vanishes as $p_y \to 1$, giving logistic growth), but at any finite
scale $c$ the equilibrium $s^*(\beta; c)$ sits on the curve of
Proposition~\ref{prop:sweep} \emph{for that $c$}, and $s^*$ drifts toward $a$ as $c$
grows.
\emph{(b) Testable prediction.} The $\beta$-dependence of the equilibrium should decay
with training length; since the experiments use early stopping with $\texttt{patience} =
10$, part of any observed $\beta$-effect is plausibly an early-stopping effect. This can
be tested by training well past saturation and comparing the converged KL across
$\beta$.
\emph{(c) Architectural fix.} The device that satisfies the fixed-$c$ assumption
\emph{by construction} is a cosine classifier: normalized features, normalized rows,
and a fixed temperature (scale). Weight normalization with a learnable gain does not
guarantee a bounded scale, and weight decay only replaces the free-scale problem with a
joint optimization problem over $(s, c)$ that must be re-derived --- it mitigates the
divergence but does not make Proposition~\ref{prop:sweep} exact. Under a fixed-norm
classifier with fixed temperature the proposition holds as stated, at the price of one
architecture change.
\end{remark}

\textbf{Where the convexity comes from (mechanism).}
Compare the three objectives on their natural solution families, as functions of the
single compression scalar:
\begin{itemize}
\item \emph{dIB on hard clusterings} (Caveats Appendix~B): $L^{\mathrm{dIB}}_\beta
= (1-\beta)H(T)$ --- compression $H(T)$ and prediction $I(Y;T) = H(T)$ are the
\emph{same} scalar, so the objective is monotone in it and the optimum sits at the
endpoint for every $\beta < 1$. Squared-dIB restores interior optima by hand,
$H(T) - \beta H(T)^2$ with unique maximizer $H(T) = 1/(2\beta)$.
\item \emph{Standard IB} (deterministic): the curve is linear ($F'(r) \equiv 1$), every
point shares one supporting slope, and the corner $\langle H(Y), H(Y)\rangle$ is
selected for all $\beta \in (0,1)$. Squared-IB convexifies via $F'(r)/(2r)$ strictly
decreasing.
\item \emph{FGIB}: the compression term is $\operatorname{KL}(q_y \| r_y) =
\tfrac{1}{2\tau^2}\|\mu_y - a v_y\|^2 + {}$(variance terms), equal to
$\tfrac{1}{2\tau^2}(s-a)^2$
at the variance optimum --- \emph{quadratic} in the alignment scalar $s$, hence strictly
convex, because the KL is measured against a \emph{fixed} reference $a v_y$:
$\operatorname{KL}(\mathcal{N}(\mu,\cdot)\|\mathcal{N}(\mu_p,\cdot))
= \tfrac{1}{2\tau^2}\|\mu - \mu_p\|^2 + {}$(variance terms) is quadratic in the
displacement
$\mu - \mu_p$ whenever $\mu_p$ does not move. Entropy-based compression ($H(T)$,
$I(X;T)$) is self-referential --- it depends only on the representation's own
distribution and enters linearly in the clustering probabilities; the fixed anchor
supplies an absolute reference frame and turns compression into a squared distance. The
convexification that squared-IB installs by squaring the compression term is thus built
into FGIB's fixed-anchor KL. Together with the strictly convex prediction term
(Lemma~\ref{lem:lagrangian}), the Lagrangian is strictly convex with a unique interior
minimizer per $\beta$ \emph{at fixed classifier scale} (Proposition~\ref{prop:sweep}),
and the resulting curve is strictly concave (Proposition~\ref{prop:curve}). With the
scale free, the same quadratic gap is driven to zero by margin growth alone
(Remark~\ref{rem:c-scale}) --- the convexification is necessary but not sufficient
without the scale condition.
\end{itemize}

\begin{remark}[The compressed endpoint is non-collapsed but still a label code --- Caveat 2 remains]\label{rem:endpoint}
VIB's fixed marginal $m(z) = \mathcal{N}(0, I)$ places the $\operatorname{KL}=0$ point
at $\mu = 0$ for \emph{all} classes: the fully compressed solution is the trivial
collapse with $\mathrm{CE} = \ln K$ (chance level), and squaring (SVIB) leaves this
corner in place since $\operatorname{KL}^2$ is stationary at $\operatorname{KL}=0$.
FGIB's per-class orthonormal anchors make the $\operatorname{KL}=0$ point a
\emph{class-separated} configuration: $\mu_y = a v_y$ with
$\langle v_y, v_{y'}\rangle = \delta_{yy'}$, so at the corner
$\mathrm{CE}(a) = \ln(1 + (K-1)e^{-ca})$, which is bounded and tends to $0$ as the
anchor scale $a$ grows, and $a = 0$ recovers the VIB corner. But class-separation does
not buy non-triviality in the sense of the Caveats paper: their Caveat~2 identifies as
uninteresting the variables on the manifold $T_\alpha$ that mix a constant map with an
exact copy of $Y$, compressing by ``forgetting the input with some probability'' rather
than by merging perceptually similar classes --- and FGIB's compressed corner
$\mu_y = a v_y$ is a deterministic function of $Y$ alone, i.e.\ a label code of exactly
that kind. The orthonormal frame is in fact \emph{opposed} to the kind of compression
Caveat~2 holds up as interesting: it places every pair of classes at the same distance
and can never express ``merge collie with retriever, keep both far from teapot''. What
the anchors do is choose \emph{one specific} label-code geometry --- class means on a
scaled simplex, a neural-collapse-like configuration --- instead of leaving it to the
optimizer, and refuse the chance-level corner. The compressed endpoint is thus designed,
not certified, and any claim that the method compresses \emph{input} semantics at that
endpoint is exactly the claim Caveat~2 rules out.
\end{remark}

\begin{remark}[Residual-information interpretation (CEB)]\label{rem:residual}
For any backward encoder, $\mathbb{E}[\operatorname{KL}(q(z|x)\|r(z|y))]
= I(X;Z|Y) + \mathbb{E}_y[\operatorname{KL}(q(z|y)\|r(z|y))]
\ge I(X;Z|Y)$ (Fischer 2020, Eqs.~9--11, 20; Lemma~\ref{lem:gap}). FGIB is therefore a
CEB bottleneck with a \emph{fixed} backward encoder: the regularizer upper-bounds the
residual information, and on the aligned family the residual vanishes identically
($q(z|x)$ depends on $y$ alone), so the regularizer equals the anchor-matching gap
$\tfrac{1}{2\tau^2}(s-a)^2$. The quadratic mechanism of Proposition~\ref{prop:sweep} is
inherited from this gap term, which exists only because the anchor is not allowed to
follow the encoder.
\end{remark}

\begin{remark}[Variance decoupling and training--evaluation consistency]
Since the classifier consumes the deterministic $h$ (and the KL does not see the
classifier), the posterior variance decouples: $\sigma^{2*} = \tau^2$ for all $\beta$, and
the trade-off is purely geometric (mean alignment). This contrasts with VIB, where the
classifier consumes the sampled $z$: the noise couples into the prediction term (by
Jensen, $\mathbb{E}_z[\mathrm{CE}(c(z),y)] \ge \mathrm{CE}(c(\mu),y)$), the rate
$R \ge I(X;Z)$ over-counts compression, and training-time sampling versus
evaluation-time mean creates a distribution mismatch (Fischer 2020, Appendix~A.1).
\end{remark}

\begin{remark}[Idealizations]\label{rem:caveats}
The class-symmetric ansatz idealizes the equilibrium geometry; it is exact in the
$\beta$-dominant regime and captures the direction of the KL force in general. For
regression, the continuous RFF anchors satisfy $\|\mu_p(y)\| \approx a$, so the KL again
contains the quadratic displacement term
$\tfrac{1}{2}\|\mu - \mu_p(y)\|^2 + {}$(variance terms) --- the same convexity mechanism
carries over, again subject to the fixed-scale condition of Remark~\ref{rem:c-scale}
(the regression classifier has no margin scale; its norm is likewise unconstrained in
the implementation). Training starts at $\operatorname{KL} \approx a^2/2$
(zero-initialized logvar head, $\sigma = 1$, $\mu \approx 0$), so the anchor scale sets
the initial regularization pressure; larger $a$ delays the approach to $s^*$. The
finite-scale caveat formerly in this remark is stated, with its fix, in
Remark~\ref{rem:c-scale}.
\end{remark}

% ---------------------------------------------------------------
\subsection{The certificate}
% ---------------------------------------------------------------

\begin{lemma}[Gap identity]\label{lem:gap}
For any conditional $r(z|y)$,
\begin{equation}
\mathbb{E}\big[\operatorname{KL}(q(z|x)\,\|\,r(z|y))\big]
= I(X;Z|Y) + \mathbb{E}_y\big[\operatorname{KL}(q(z|y)\,\|\,r(z|y))\big] .
\label{eq:gap}
\end{equation}
\end{lemma}

\begin{proof}
$\mathbb{E}[\operatorname{KL}(q(z|x)\|r(z|y))] = \mathbb{E}[\ln q(z|x) - \ln r(z|y)]$
and $I(X;Z|Y) = \mathbb{E}[\ln q(z|x) - \ln q(z|y)]$ (Fischer 2020, Eqs.~9--10), where
both expectations are over $p(x,y)q(z|x)$; subtracting,
$\mathbb{E}[\ln q(z|y) - \ln r(z|y)] = \mathbb{E}_y[\operatorname{KL}(q(z|y)\|r(z|y))]$
since the argument depends on $z, y$ only. \hfill\qed
\end{proof}

\begin{lemma}[Moment matching]\label{lem:mm}
Within the per-class Gaussian family with \emph{diagonal covariance}, the gap of
Lemma~\ref{lem:gap} is
minimized by the moment-matched Gaussian $b^*_y = \mathcal{N}(m_y, S_y)$ with
$m_y = \mathbb{E}\,q(z|y)$, $S_y = \operatorname{diag}\operatorname{Var} q(z|y)$.
\end{lemma}

\begin{proof}
$\operatorname{KL}(p\|q)$ as a function of $q$ is minimized by matching the first two
moments of $p$ (standard; for Gaussians the minimizer is unique). \hfill\qed
\end{proof}

\begin{proposition}[Variational ordering: CEB's bound is tighter in value]\label{prop:ordering}
For every fixed encoder $q(z|x)$,
$\mathrm{ReX}(\theta^*) \le \mathrm{ReX}_F$ where $\theta^*$ is CEB's optimal backward
encoder, and the excess is the \emph{anchor gap}
\begin{equation}
\Delta := \mathrm{ReX}_F - \mathrm{ReX}(\theta^*)
= \mathbb{E}_y\big[\operatorname{KL}(q(z|y)\,\|\,N(a v_y, \tau^2 I))\big]
- \mathbb{E}_y\big[\operatorname{KL}(q(z|y)\,\|\,N(m_y, S_y))\big] \ge 0 ,
\label{eq:anchor-gap}
\end{equation}
with equality iff the anchors equal the moment-matched marginals. Here $\theta^*$ ranges
over the \emph{ideal} CEB backward encoder family --- per-class Gaussians with free mean
and \emph{free diagonal covariance} (the family of Lemma~\ref{lem:mm}, matching the
diagonal implementation), i.e.\ any conditional $r(z|y)$ of that family the theory can
express; the code's
single affine \texttt{prior\_net} is only an approximation of this family, and the
ordering is a statement about the idealized CEB, not about the implemented
parameterization.
On the aligned family of the sweep analysis ($q(z|y) = \mathcal{N}(s v_y, \sigma^2 I)$,
so $I(X;Z|Y) = 0$ and $b^* = q(z|y)$ exactly),
\begin{equation}
\Delta = \tfrac{1}{2\tau^2}(s-a)^2 + \tfrac{z}{2}\big(\sigma^2/\tau^2 - 1 - \ln(\sigma^2/\tau^2)\big) ,
\label{eq:aligned-gap}
\end{equation}
i.e.\ $\mathrm{ReX}_F = \Delta$ and $\mathrm{ReX}(\theta^*) = 0$: FGIB's bound value
\emph{equals} the sweep-driving compression term of Proposition~\ref{prop:sweep}, while
CEB's bound collapses to the (zero) residual.
\end{proposition}

\begin{proof}
The anchor $N(a v_y, \tau^2 I)$ is a member of CEB's per-class Gaussian family, so the
minimum over the family is no larger than the value at the anchor: the ordering and
Eq.~\eqref{eq:anchor-gap} follow from Lemmas~\ref{lem:gap} and~\ref{lem:mm}.
On the aligned family the marginal is itself Gaussian, so $b^* = q(z|y)$ gives gap $0$,
$I(X;Z|Y) = 0$ (the posterior is class-constant), and
$\operatorname{KL}(\mathcal{N}(s v_y, \sigma^2 I)\|\mathcal{N}(a v_y, \tau^2 I))
= \tfrac{1}{2\tau^2}(s-a)^2 + \tfrac{z}{2}(\sigma^2/\tau^2-1-\ln(\sigma^2/\tau^2))$.
\hfill\qed
\end{proof}

\begin{remark}[Tightness--force trade-off]\label{rem:force}
Proposition~\ref{prop:ordering} is the price of the fixed anchor, and it is paid for a
reason: in CEB the gap is a function of the \emph{backward encoder alone} ---
$b_\theta \to q(z|y)$ eliminates it with the encoder held fixed, so CEB's bound can
become tight while the representation stays uncompressed: its tightness certifies
nothing about compression, and the gap is an unobservable optimization residual that
tracks the forward encoder. FGIB's gap can shrink only if the \emph{posterior} moves
toward the anchors (the anchor itself cannot move), so the force is persistent and, on
the aligned family, the bound value \emph{is} the anchor-matching gap
$\tfrac{1}{2\tau^2}(s-a)^2$, the sweep-driving compression term
(Eq.~\eqref{eq:aligned-gap}). The qualifier is necessary: ``the posterior'' includes
the side head $A$, which is trainable and unused at inference, so part of the movement
can be absorbed without the deployed representation $h$ moving
(Remark~\ref{rem:A-channel}). On the aligned family, where $A$ is fixed by the ansatz,
the equivalence is exact; off it, bound tightness is monotonically equivalent to
alignment of the posterior, not to compression of the deployed $h$.
\end{remark}

\begin{proposition}[The inverse linear transform: the object of the bound]\label{prop:inverse}
Let the side head $T(h) = Ah$ be \emph{invertible} ($z = d$, unconstrained), and let
the trained posterior be isotropic, $\Sigma(x) = \sigma_p^2 I$ with $\sigma_p^2 = \tau^2$
the anchor variance (the KL variance term is minimized at $\sigma^2 = \tau^2$, so this
holds at the variance optimum).
Then $\tilde z := A^{-1} z$ gives
$q'(\tilde z|x) = \mathcal{N}(h(x), \tau^2 (A^{\top}A)^{-1})$ --- the posterior is
centered at the \emph{inference representation} $h$ with side noise
$\tau^2 (A^{\top}A)^{-1}$ --- and the transformed anchors
$r'(\tilde z|y) = \mathcal{N}(a\, A^{-1} v_y, \tau^2 (A^{\top}A)^{-1})$ form an
orthonormal frame at scale $a$ in the induced metric $A^{\top}A$
($\langle A^{-1}v_i, A^{-1}v_j\rangle_{A^{\top}A} = \delta_{ij}$: the fixed frame
$\{v_y\}$ read in canonical coordinates). By KL invariance under invertible linear maps
(Lemma~\ref{lem:kl-invariance}),
\begin{equation}
\mathrm{ReX}_F
= I(X;\, h + \tau\,(A^{\top}A)^{-1/2}\varepsilon\,|Y)
+ \mathbb{E}_y\big[\operatorname{KL}\big(q'(\tilde z|y)\,\|\,r'(\tilde z|y)\big)\big],
\qquad \varepsilon \sim \mathcal{N}(0, I) .
\label{eq:object}
\end{equation}
The bound's object is the inference representation perturbed by \emph{known,
side-path-only} noise (its covariance is a model parameter), and the anchor geometry is
preserved exactly under the transform, in the induced metric. The perturbation,
however, is small only if the side head is well-conditioned: its magnitude is governed
by $A^{\top}A$, which is unconstrained
(Proposition~\ref{prop:conditioning}(iii)).
\end{proposition}

\begin{proof}
$A^{-1}A = I$ and $A^{-1} z = h + A^{-1}\varepsilon$ for $\varepsilon \sim
\mathcal{N}(0, \Sigma)$; $A^{-1}\Sigma A^{-\top} = \tau^2 A^{-1}A^{-\top}
= \tau^2 (A^{\top}A)^{-1}$ for $\Sigma = \tau^2 I$; the prior covariance maps to
$A^{-1} (\tau^2 I) A^{-\top} = \tau^2 (A^{\top}A)^{-1}$; and
$\langle A^{-1}v_i, A^{-1}v_j\rangle_{A^{\top}A}
= v_i^{\top} A^{-\top} A^{\top} A\, A^{-1} v_j = \delta_{ij}$.
Equation~\eqref{eq:object} is Lemma~\ref{lem:gap} applied to the transformed pair.
\hfill\qed
\end{proof}

\begin{proposition}[Orthogonal-factor conditioning: the reference is fixed and well-conditioned]\label{prop:conditioning}
FGIB's anchor frame $V = [v_1, \dots, v_K]$ is the orthonormal factor of a random matrix
(QR: $G = QR$, $V := Q$, $V^{\top}V = I_K$), and the anchor covariance is
$\tau^2 I$ with $\tau^2 > 0$.
Then, for every anchor scale $a \ge 0$ and throughout training:
\begin{enumerate}
\item[(i)] the \emph{reference} of the matching --- the \emph{unscaled} mean frame $V$
and the covariance $\tau^2 I$ --- is fixed (non-trainable buffers) and has condition
number $1$: the certificate is a distance to a compact, non-degenerate,
class-separated (for $a > 0$) target. (At $a = 0$ the mean matrix $a V$ is zero and
has no condition number; the condition-$1$ statement attaches to $V$ and $\tau^2 I$,
not to $a V$.)
\item[(ii)] in side-mean coordinates the certificate's displacement metric is a
scaled identity: the KL's displacement term is $\tfrac{1}{2\tau^2}\|\mu - a v_y\|^2$
with gradient $(\mu - a v_y)/\tau^2$;
\item[(iii)] on the shared-encoder side the metric is $A^{\top}A$, where $A$ is a
learnable unconstrained map, so its \emph{global} condition number is not controlled by
construction --- and the aligned family does not control it either. An earlier version
inferred condition number $1$ on the class-mean subspace from the identity
$h_i^{\top} A^{\top} A\, h_j = s^2\delta_{ij}$; that inference is invalid, since
$\{h_i\}$ is not a Euclidean orthonormal basis. Counterexample: with
$A = \operatorname{diag}(\varepsilon, 1)$, $h_1 = (1/\varepsilon, 0)$,
$h_2 = (0, 1)$ and $s = 1$, one has $A h_i = v_i$ yet
$\kappa(A^{\top}A) = \varepsilon^{-2}$ --- and the inverse-transform perturbation of
Proposition~\ref{prop:inverse} is correspondingly $\varepsilon^{-2}$-large on the first
direction. The counterexample is not adversarial: it is the direction the gap gradient
points in, since the KL is minimized by fitting the anchor target on $h$ in least
squares, which suppresses within-class directions
(Remark~\ref{rem:A-channel}). Any guarantee on this side therefore requires an explicit
constraint on $A$ (orthogonal parameterization, weight norm, or the like) or an
empirical measurement of $\kappa(A^{\top}A)$ over training.
\end{enumerate}
CEB's backward encoder $b_\theta(z|y) = \mathcal{N}(m_\theta(y), S_\theta(y))$ is an
unconstrained linear map of the label: \emph{both} sides of its matching are trainable,
its covariance diagonal can lie anywhere in the clamp $[e^{-10}, e^{10}]$ (condition
number up to $e^{20}$), its means can be rank-deficient (class-merged codes, or the
collapse solution $m_\theta \equiv 0$), and the variance race drives $S_\theta$ to the
clamp, degenerating the certificate into a mean-matching force of weight $1/s = e^{10}$
(Prop.~\ref{prop:variance}, subject to Remark~\ref{rem:p3-status}). What the fixed QR
reference guarantees for FGIB is therefore the \emph{reference} side only: a compact,
non-degenerate, class-separated target whose conditioning cannot degrade. On the
encoder side FGIB's bound is exposed to the same degeneration mechanism CEB's is, with
$A$ playing the role of the second trainable party.
\end{proposition}

\begin{proof}
$V^{\top}V = I_K$: all singular values of $V$ are $1$; the covariance $\tau^2 I$ has
all eigenvalues $\tau^2$; both are frozen buffers, so (i) holds. The diagonal closed
form of
$\operatorname{KL}(\mathcal{N}(\mu, \sigma^2 I)\,\|\,\mathcal{N}(a v_y, \tau^2 I))$
has displacement term $\tfrac{1}{2\tau^2}\|\mu - a v_y\|^2$, giving (ii). For (iii):
the identity
$h_i^{\top} A^{\top} A\, h_j = \mu_i^{\top}\mu_j = s^2 v_i^{\top} v_j = s^2\delta_{ij}$
holds on the aligned family and only says that $\{h_i\}$ is $A^{\top}A$-orthogonal by
construction; it implies nothing about the spectrum of $A^{\top}A$, and the displayed
example verifies the failure. The KL gradient on $A$ is
$(\mu - a v_y)\,h^{\top}/\tau^2$, an alignment force toward the orthonormal frame that
vanishes only at perfect alignment \emph{on non-degenerate samples} (it also vanishes
pointwise at $h = 0$, or by cancellation across a batch) --- but nothing in that force
resists a vanishing within-class component of $A$. For CEB, the displacement term is
$\tfrac12(m - g)^{\top} S_\theta^{-1}(m - g)$: the metric $S_\theta^{-1}$ has condition
number up to $e^{20}$ under the logvar clamp, and the gradient weight $1/s$ at the
collapse equilibrium is Prop.~\ref{prop:variance}. \hfill\qed
\end{proof}

\begin{remark}[Two senses of tightness]\label{rem:senses}
In the pure variational \emph{value} sense (same encoder, same Gaussian family), CEB's
bound is tighter: $\mathrm{ReX}(\theta^*) \le \mathrm{ReX}_F$
(Prop.~\ref{prop:ordering}), the excess being the anchor gap --- the price of the fixed
anchor. The \emph{object} sense does not reverse this ordering: both bounds are valid
upper bounds on their own noisy codes, and neither bounds the residual of a deployed
deterministic continuous representation, which is infinite under the standing
non-atomic assumption (Prop.~\ref{prop:object}). The structural difference is where the certified-versus-
deployed gap lives --- in CEB it is driven by the prediction term inside the inference
path; in FGIB it is a side-path quantity whose size is governed by the conditioning of
$A$ (Prop.~\ref{prop:conditioning}(iii)) --- not which bound is ``tighter'' or
``valid''. The learnability of the side head $A$ bears on that difference: it is a
second trainable party in the matching, unused at inference
(Remark~\ref{rem:A-channel}), so it can absorb part of the anchor gap without moving
the deployed representation. On the aligned family, where $A$ is fixed by the ansatz,
its optimization drives the gap to the equilibrium value $\tfrac12(s^*-a)^2$ at fixed
classifier scale (Prop.~\ref{prop:sweep}); off the ansatz this is an open question
(Section~\ref{sec:limits}).
\end{remark}

\begin{remark}[Unconstrained side head (the implementation)]\label{rem:orth}
Proposition~\ref{prop:inverse} holds for any invertible linear head: the KL invariance
is exact, and the object claim needs only that the perturbation covariance
$\tau^2(A^{\top}A)^{-1}$ be \emph{known} ($A$ is a model parameter) --- but the
object is informative about $h$ only insofar as that covariance is moderate, which the
construction does not guarantee (Prop.~\ref{prop:conditioning}(iii)). The diagonal
variational family is not closed under $A^{-1}$ (Remark~\ref{rem:closure}), so within
the diagonal parameterization of the codebase the equivalence is exact at the
distribution level and approximate as a variational identity; a full-covariance side
head would make it exact within the family. No orthogonality constraint on $A$ is
required --- the orthogonality of the method lives in the fixed anchor frame and in the
equilibrium it induces (the aligned family).
\end{remark}

\begin{remark}[Anchor-frame reading]\label{rem:frame}
With the orthonormal anchor matrix $V$ (columns $v_y$, $K \le z$), left multiplication
by $V^{\top}$ expresses the anchor-space components of $z$ in class coordinates: the
anchors become $a\,e_y$ and the gap of Lemma~\ref{lem:gap} decomposes into $K$
per-class terms plus, when $K < z$, an orthogonal-complement term that the reading
ignores ($V^{\top}$ is a projection, not a coordinate change). The orthogonal frame is
the geometric content of the prior, and it is preserved by the invertible linear
pull-back $A^{-1}$ of Proposition~\ref{prop:inverse} --- not by $A^{\top}$, which is
the Jacobian factor that appears in gradients and equals $A^{-1}$ only when $A$ is
orthogonal. An earlier version conflated the two.
\end{remark}

\begin{remark}[Positioning against VIB and NIB]\label{rem:position}
VIB's rate bounds the \emph{full} mutual information $I(X;Z)$ with a marginal prior;
CEB removes the retained label information ($\mathrm{ReX} = R - I(Y;Z)$ at the optima),
and FGIB inherits the residual-level object while replacing the learnable backward
encoder with a fixed geometric prior (Prop.~\ref{prop:ordering}).
NIB's non-parametric bound $\widehat I \ge I(X;M)$ is $y$-unconditioned and tight only
when the components share a covariance and form well-separated clusters (Kolchinsky et
al.\ 2017, Eq.~10 and the discussion following Eq.~11); when the noise is small relative
to the cluster spread it saturates at $\ln N$ and its gradient vanishes (their
Section~IVA) --- loose exactly in the regime where the parametric FGIB bound is
well-behaved. The two bounds are complementary: NIB certifies the full compression
non-parametrically when the geometry is favorable; FGIB bounds the residual of a
\emph{noisy copy} of the deployed representation parametrically, with a geometry of its
own choosing --- subject to the conditioning caveat of
Proposition~\ref{prop:conditioning}(iii).
\end{remark}

\begin{remark}[Degenerate anchors]\label{rem:a0}
At $a = 0$ the anchors collapse to $N(0, \tau^2 I)$ for all classes and the gap loses its
class
structure: $\Delta = \mathbb{E}_y[\operatorname{KL}(q(z|y)\|N(0,\tau^2 I))]$ --- a
VIB-style
marginal gap; the structural separation of Proposition~\ref{prop:object}(b) still
holds, but the compressed endpoint becomes the trivial collapse
(Remark~\ref{rem:endpoint}).
\end{remark}

% ---------------------------------------------------------------
\subsection{What is proved and what is open}\label{sec:limits}
% ---------------------------------------------------------------

\begin{remark}[Scope of the results]
\textbf{Proved.} (1) The dead-knob proposition for CEB and the MNI reading of it
(Prop.~\ref{prop:dead}, Remark~\ref{rem:mni}); (2) the gap identity, Gaussian moment
matching, and the value ordering of CEB's bound over FGIB's at a fixed encoder
(Lemmas~\ref{lem:gap}, \ref{lem:mm}; Prop.~\ref{prop:ordering}); (3) the blind
certificate on the $T_\alpha$ manifold and the $1{:}\beta$ exchange rate
(Prop.~\ref{prop:flat}); (4) KL invariance under invertible linear maps and the
pointwise equivalence of the side path to a CEB bottleneck at a fixed side head
(Lemmas~\ref{lem:kl-invariance}, \ref{lem:mi-invariance}; Prop.~\ref{prop:fgib-ceb});
(5) the variance decoupling $\sigma^{2*} = \tau^2$ and the non-collapsed compressed endpoint
(Lemma~\ref{lem:lagrangian}, with $\sigma^{2*} = \tau^2$;
Remark~\ref{rem:endpoint}); (6) the one-dimensional sweep
and its strict concavity \emph{at fixed classifier scale}, as a KL--CE surrogate curve
on the aligned family (Props.~\ref{prop:sweep}, \ref{prop:curve};
Remark~\ref{rem:surrogate} --- not a recovered IB curve); (7) the well-conditioning of the fixed
reference side (Prop.~\ref{prop:conditioning}(i)--(ii)).

\textbf{Open or conditional.} (1) The side head $A$ is a trainable, inference-unused
party in the matching (Remark~\ref{rem:A-channel}); combined with the
counterexample of Prop.~\ref{prop:conditioning}(iii), the perturbation in
Prop.~\ref{prop:inverse} can be arbitrarily large, so the bound's object may be far from
the deployed representation exactly where the residual lives. A constraint on $A$ or a
measurement of $\kappa(A^{\top}A)$ over training would close or quantify this channel.
(2) The global sweep requires a bounded classifier scale
(Remark~\ref{rem:c-scale}); the current implementation leaves the scale free and the
proposition conditional. A fixed-temperature cosine classifier would satisfy the
fixed-$c$ assumption by construction; weight decay only mitigates the scale divergence
and replaces the problem with a new joint optimization over $(s,c)$ that must be
re-derived. (3) Neither CEB's nor FGIB's certificate upper-bounds the deployed
deterministic residual (Prop.~\ref{prop:object}); claims of inference-representation
compression are hypotheses, not conclusions. (4) Caveat~2 is not resolved
(Remark~\ref{rem:endpoint}): the compressed endpoint is a chosen label-code geometry,
not certified semantic compression. (5) The two-stage variance dynamics of
Prop.~\ref{prop:variance}(ii)--(iii) await the measurement of
Remark~\ref{rem:p3-status}. (6) The regression RFF anchor analysis inherits the same
fixed-scale and conditioning caveats and is presently a transfer argument, not a proof.
\end{remark}
。
% 记号沿用 Fischer (2020)、Kolchinsky & Tracey (2019, Caveats)、Kolchinsky et al. (2017, NIB)、
% Alemi et al. (2017, VIB)。
%
% 核心结论（FGIB 论文问题陈述与解法）：
% 第一部分（四条失效通道）：
%   P1 退化继承（死旋钮）：CEB ≡ IB(β=1+γ>1)，确定性场景下所有 γ>0 的唯一最优解钉在角点，
%      γ 扫不出任何中间点（SIB Caveat 1 的退化半区，经 β'=1/(1+γ)∈(0,1) 映射）。
%      该角点恰是 CEB 自己的 MNI 目标，rem:mni 正面处理这一反驳；
%   P2 证书盲视 + 1:β 交换率：T_α 流形上 ReX 恒为 0（证书对标签保留量完全盲视），
%      且目标以 1 nat 标签信息 ↔ β nat 认证残差的比率交换；
%   P3 方差竞赛：KL 两侧皆可训练，分类项压 σ→0、钳位处退化为权重 e^{10} 的均值匹配。
%      (ii)(iii) 是受力方向分析而非收敛证明，且可实测（rem:p3-status）；
%   P4 对象错位：两种方法的界都只界住各自的含噪码，都不界住确定性连续推理表示的残差
%      （后者为 +∞）。差别是结构性的（噪声在不在推理路径上），不是紧度排序。
% 第二部分（FGIB 机制，以及它没能关闭的通道）：
%   M1 等价性：旁路 = 固定 A 下作用在 h+已知噪声上的 CEB 瓶颈（逐步成立，非全程固定参考）；
%   M2 扫参：固定分类器尺度 c 时 L_β(s)=CE(s)+(β/2)(s−a)² 严格凸、每 β 唯一内点、曲线严格凹；
%      c 自由时下确界对所有 β 均为 0（rem:c-scale，给出权重归一化分类器的架构修法）；
%   M3 证书：gap 恒等式 + 矩匹配 + 固定编码器下的值排序；固定参考的力持续性成立，
%      但侧头 A 是推理无关的第二个自由方（rem:A-channel），构成未关闭的通道。
%
% 修订记录（第二轮）：理论/实现分层后统一引入锚点方差 τ²（先验 N(a v_y, τ²I)，默认 τ²=1，
% 后验方差最优 σ²*=τ²，σ_p 为公共方差记号；point-anchor 极限统一写为 τ→0）；理论旁路固定为
% 可逆线性头 T(h)=Ah（实现用带 bias 的 nn.Linear 近似，见理论-代码分层 caveat）；
% 明确 sweep 证明的是对齐族上 KL–CE surrogate 曲线（rem:surrogate），不是 IB 互信息曲线；
% 无穷互信息结论统一挂靠分布性 standing assumption；rem:c-scale 的架构修法收紧为
% 固定温度 cosine 分类器（weight decay 只缓解不发散）；a=0 退化点、条件数表述、
% 梯度消失条件、row-space 表述等小项修正。
%
% 修订记录（第一轮）：删除三处错误论断——(1) P3 原称 ReX≈0 时 I(X;Z|Y) 发散，与 gap 恒等式矛盾；
% (2) prop:object(c) 原称 FGIB 的界"更紧/唯一有效"，由 DPI 不成立；(3) prop:conditioning(iii)
% 原从 Gram 恒等式推出条件数为 1，有反例。修正 P2 的换算、Caveat 2 的主张、rem:frame 的
% A^T/V^T、mu_head 的仿射与非推理路径身份。未证结论统一收入 sec:limits。

\newtheorem{lemma}{Lemma}
\newtheorem{proposition}{Proposition}
\newtheorem{remark}{Remark}

% =====================================================================
\section*{Part I. The four failure channels of the conditional bottleneck}
% =====================================================================

\textbf{Notation.}
$X$ is the input and $Y$ the label with $K$ classes; the encoder induces the Markov chain
$Y - X - Z$. The \emph{deterministic scenario} is $Y = f(X)$ ($H(Y|X) = 0$). Following
Fischer (2020), we write the MNI point $I(X;Y) = I(X;Z) = I(Y;Z)$ (Eq.~1), the
chain-rule identity $I(X;Z|Y) = I(X;Z) - I(Y;Z)$ (Eq.~4), the CEB objective
$\mathrm{CEB} \equiv \min_Z I(X;Z|Y) - \gamma I(Y;Z)$ (Eq.~5), the IB objective
$\mathrm{IB} \equiv \min_Z I(X;Z) - \beta_{\mathrm{IB}} I(Y;Z)$ (Eq.~16), the
variational objective
$\mathrm{VCEB} \equiv \min_{e,b,c} \langle \log e(z|x)\rangle - \langle \log b(z|y)\rangle
- \gamma \langle \log c(y|z)\rangle$ (Eq.~15), and the residual certificate
$\mathrm{ReX} \equiv \langle \log e(z|x)\rangle - \langle \log b(z|y)\rangle
\ge I(X;Z|Y)$ (Eq.~20), where $b(z|y)$ is a \emph{trainable} backward encoder (we write
$q(z|x)$ for Fischer's encoder $e(z|x)$; the two notations are interchangeable
throughout). In the loss convention $\mathrm{CE} + \beta\,\mathrm{KL}$ used throughout, the weights are
$\beta = 1/\gamma$ and $\beta_{\mathrm{IB}} = 1 + \gamma = 1 + 1/\beta$. Following
Kolchinsky \& Tracey (2019), we write the IB curve $F(r)$ (their Eq.~1), the IB
Lagrangian $L^{\mathrm{IB}}_\beta = I(Y;T) - \beta I(X;T)$ (their Eq.~2), the
DPI-saturating manifold $T_\alpha := B_\alpha \cdot Y$ (their Eq.~6), the bound
$L^{\mathrm{IB}}_\beta \le (1-\beta)H(Y)$ (their Eq.~7), and their footnote~4 (the
$\beta > 1$ regime of trivial maximizers). The variational objects are
\begin{align*}
R &= \mathbb{E}\big[\operatorname{KL}(q(z|x)\,\|\,m(z))\big] \ge I(X;Z)
&&\text{(VIB rate, marginal prior $m$; Fischer Eq.~19, Alemi et al.\ Eq.~14)} \\
\mathrm{ReX} &= \mathbb{E}\big[\operatorname{KL}(q(z|x)\,\|\,b_\theta(z|y))\big]
\ge I(X;Z|Y)
&&\text{(CEB residual; Fischer Eq.~20)} \\
\widehat I &= -\frac{1}{N}\sum_i \ln\Big[\frac{1}{N}\sum_j
e^{-\|f_i-f_j\|^2/(2\sigma^2)}\Big] \ge I(X;M)
&&\text{(NIB, non-parametric; Eq.~16, the homoscedastic case of Eq.~10)} \\
\mathrm{ReX}_F &= \mathbb{E}\big[\operatorname{KL}(q(z|x)\,\|\,r(z|y))\big]
\ge I(X;Z|Y)
&&\text{(FGIB, fixed anchor)}
\end{align*}
(Fischer 2020, Eqs.~18--20; Alemi et al.\ 2017, Eq.~14; Kolchinsky et al.\ 2017,
Eqs.~10 and 16). Fischer's comparison of the first two: at the respective prior optima
($m = q(z)$, $b = q(z|y)$), $\mathrm{ReX} = R - I(Y;Z)$ --- the conditional prior removes
exactly the retained label information, so CEB's bound is tighter than VIB's by
$I(Y;Z)$. We ask the same question one level down: where does FGIB's fixed-anchor bound
sit relative to CEB's learnable-backward bound?

\begin{proposition}[Channel P1, dead knob: every $\gamma > 0$ selects the same corner]\label{prop:dead}
Assume $Y = f(X)$. For every $\gamma > 0$,
\begin{equation}
\min_Z L_\gamma(Z) = -\gamma H(Y),
\qquad
L_\gamma(Z) := I(X;Z|Y) - \gamma I(Y;Z),
\label{eq:min}
\end{equation}
attained exactly on the corner family
$\{Z : I(Y;Z) = H(Y),\; I(X;Z|Y) = 0\}$ (e.g.\ $Z = Y$, or $Z = (Y, N)$ with
$N \perp X \mid Y$): the corner is unique in the information plane, while at the
representation level it is an entire family. Consequently the map
$\gamma \mapsto \big(I(X;Z^*|Y),\, I(Y;Z^*)\big)$ is constant: no interior point of the
IB curve is ever a Lagrangian minimizer, and CEB inherits the deterministic degeneration
of the Caveats paper (their Issue~1) over its \emph{entire} parameter range, since
$\beta_{\mathrm{IB}} = 1 + \gamma > 1$ always.
\end{proposition}

\begin{proof}
By Fischer's Eq.~4, $L_\gamma(Z) = I(X;Z) - (1+\gamma)I(Y;Z)$. The DPI
($I(Y;Z) \le I(X;Z)$, by the Markov chain $Y - X - Z$) gives
\begin{equation}
L_\gamma(Z) \ge I(Y;Z) - (1+\gamma)\,I(Y;Z)
= -\gamma\,I(Y;Z) \ge -\gamma H(Y),
\label{eq:chain}
\end{equation}
with equality iff both inequalities bind: $I(X;Z) = I(Y;Z)$ and $I(Y;Z) = H(Y)$.
$Z = Y$ attains the bound. Fischer (2020, p.~7) states the underlying equivalence
explicitly: CEB and IB are the same objective for $\gamma = \beta_{\mathrm{IB}} - 1$.
In the Caveats paper's maximizing convention, minimizing $L_\gamma$ is maximizing
$L^{\mathrm{IB}}$ with $\beta' = 1/\beta_{\mathrm{IB}} = 1/(1+\gamma) \in (0,1)$ ---
precisely the range in which their Eq.~7 pins every maximizer at the copy corner (their
Issue~1); their footnote~4 covers the complementary regime $\beta' > 1$ (i.e.\
$\gamma < 0$, outside CEB's range), whose maximizers are the trivial collapse. \hfill\qed
\end{proof}

\begin{remark}[The corner is the MNI point: why that does not rescue the knob]\label{rem:mni}
In the deterministic scenario the corner of Proposition~\ref{prop:dead} \emph{is} the MNI
point: $I(Y;Z) = H(Y) = I(X;Y)$ together with $I(X;Z|Y) = 0$ gives
$I(X;Y) = I(X;Z) = I(Y;Z)$, which is Fischer's Eq.~1. CEB's declared target and the
degenerate attractor therefore coincide, and Kolchinsky \& Tracey concede the point
themselves (their Section~4): ``if $Y = f(X)$ and one's goal is precisely to find the corner
point $\langle H(Y), H(Y)\rangle$ (maximum compression at no prediction loss), then our
results suggest that the IB Lagrangian provides a robust way to do so, being invariant to
the particular choice of $\beta$.'' Read that way, Proposition~\ref{prop:dead} states
hyperparameter robustness rather than failure. Three things keep it from closing the matter.
\begin{enumerate}
\item[(i)] \emph{The corner is one point in the information plane and an entire equivalence
class of representations.} The objective is exactly constant on
$\{Z : I(Y;Z) = H(Y),\ I(X;Z|Y) = 0\}$, which contains $Z = Y$, $Z = (Y,N)$ with
$N \perp X \mid Y$, and every bijective recoding of these. Adversarial robustness, OoD
detection and calibration --- the properties CEB is argued for --- are properties of the
\emph{geometry} of $Z$, on which the objective has no opinion; which member of the class is
reached is settled by the optimizer. That is channels P3 and P4 below.
\item[(ii)] \emph{CEB's own reading of the knob is unavailable at the exact optimum.}
Fischer (2020, p.~7) writes ``$\rho < 0$ may capture less than the MNI. $\rho > 0$ may
capture more'', whereas Proposition~\ref{prop:dead} places every $\gamma > 0$ at the MNI
point exactly. His argument for preferring $\gamma = 1$ is correspondingly not about
minimizers but about the optimizer: weighting the two appearances of $I(Y;Z)$ differently
``forces the optimization procedure to count bits of $I(Y;Z)$ in two different ways'' (his
Eq.~8 and surrounding discussion).
\item[(iii)] \emph{The knob does move trained models, which locates what it acts on.} On
Fashion MNIST --- a deterministically labelled dataset --- Fischer's own Figure~4 reports
the rate lower bound $R_X$ rising from $\approx 2.0$ to $\approx 5.0$ nats as $\rho$ goes
from $-1$ to $5$, with $\rho = 0$ landing on $I(X;Y) = \log 10 \approx 2.3$ nats. Since
every $\gamma > 0$ shares one exact optimum, that axis is not a path along the IB curve: it
is the distance between the optimum of the variational objective over a restricted family
and the point at which training was stopped. A knob acting through the approximation error
is precisely the knob that Propositions~\ref{prop:flat} and~\ref{prop:variance} show cannot
be read off the certificate.
\end{enumerate}
Kolchinsky \& Tracey state that their results ``are based on analytically-provable
properties of the IB curve, i.e., global optima'' and ``do not concern practical issues of
optimization''; Proposition~\ref{prop:flat} supplies the finite-$\beta$ variational
counterpart. The account is owed by any method reporting a $\beta$-sweep on
deterministically labelled data --- FGIB and the experiments of this paper included
(Remark~\ref{rem:c-scale}).
\end{remark}

\begin{proposition}[Channel P2, a blind certificate and a $1{:}\beta$ exchange rate]\label{prop:flat}
On the manifold $T_\alpha = B_\alpha \cdot Y$ of the Caveats paper (their Eq.~6),
$I(X;T_\alpha|Y) = 0$ for every $\alpha$, and with the optimal backward encoder and
classifier,
\begin{equation}
\mathrm{ReX}(T_\alpha) = 0,
\qquad
\mathrm{CE}(T_\alpha) = H(Y) - I(Y;T_\alpha) =: \delta_\alpha,
\qquad
\mathrm{VCEB}(T_\alpha) = \gamma\,\delta_\alpha .
\label{eq:vceb}
\end{equation}
Two consequences.
\emph{(a) The certificate is blind to label retention.} $\mathrm{ReX} = 0$ holds from the
total collapse ($\alpha = 0$, no label information) to the full copy ($\alpha = 1$), so its
tightness is compatible with retaining all, part, or none of $Y$. Fischer's reading of the
certificate as a meter --- ``observing ReX tells us exactly how many bits we are from the
MNI point'' (p.~7) --- therefore holds only along the $I(X;Z|Y)$ axis; on the necessity axis
the meter reads $0$ at chance level.
\emph{(b) Exchange rate.} With the optimal classifier the objective is
$L = \mathrm{CE} + \beta\,\mathrm{ReX} = \delta + \beta\,\mathrm{ReX}$, so it is indifferent
between giving up $\delta$ nats of label information and saving $\delta/\beta$ nats of
certified residual. In the heavy-regularization regime the bottleneck is supposed to act in,
the whole label axis --- up to $H(Y)$ nats --- is purchasable for $H(Y)/\beta$ nats of
certified compression.
\end{proposition}

\begin{proof}
$T_\alpha$ is a function of $Y$ alone, so $T_\alpha \perp X \mid Y$ and
$I(X;T_\alpha|Y) = 0$; the Caveats paper shows $I(Y;T_\alpha)$ sweeps $[0, H(Y)]$.
With the true backward encoder ($b = p(z|y)$) the certificate is tight:
$\mathrm{ReX} = I(X;T_\alpha|Y) = 0$ (Fischer Eqs.~9--11, 20). The optimal classifier
achieves $\langle \log c(y|z)\rangle = -H(Y|T_\alpha)$ (Fischer Eq.~12), i.e.\
$\mathrm{CE} = H(Y|T_\alpha) = \delta_\alpha$, and the VCEB value follows from Eq.~15.
Part~(b) is the definition of the objective evaluated at $\mathrm{CE} = \delta$.
\hfill\qed
\end{proof}

\begin{remark}[What the flatness claim is, and what it is not]\label{rem:convention}
An earlier form of this argument read the near-optimality band off the \emph{value}
$\mathrm{VCEB} = \gamma\delta \to 0$. That reading is a scaling artifact:
$\mathrm{VCEB} = \gamma\,(\mathrm{CE} + \beta\,\mathrm{ReX}) = \gamma L$, so the two
conventions differ by a positive factor and share their minimizers, and on $T_\alpha$ itself
the objective is $L = \delta$ in either convention and strictly prefers $\alpha = 1$. What
is convention-free is the \emph{ratio} of Proposition~\ref{prop:flat}(b): the prediction
term is weaker than the compression term by exactly the factor $\beta$, so in units of the
compression term the entire manifold lies within $H(Y)/\beta$ of optimal. Under a
scale-invariant optimizer (Adam, used throughout the experiments) only the ratio matters;
under plain SGD at fixed learning rate the convention additionally rescales the effective
step size. Combined with Proposition~\ref{prop:flat}(a) --- no certified compression is
bought anywhere on the manifold --- the operative statement is that the label axis is cheap
and unmetered, not that the objective is literally flat along it.
Pushing $\beta$ up also drives $\beta_{\mathrm{IB}} = 1 + 1/\beta \to 1^{+}$: the objective
approaches the Caveats paper's $\beta = 1$ degeneracy (their Eq.~7 discussion), at which
\emph{every} point of the diagonal is simultaneously optimal, from above. Fischer (2020,
p.~7) anticipates the direction --- ``$\rho < 0$ may capture less than the MNI'' --- but by
Proposition~\ref{prop:dead} that cannot happen at the exact optimum;
Proposition~\ref{prop:flat} gives the rate at which the variational objective is willing to
let it happen.
\end{remark}

\begin{remark}[Approximately deterministic labels]\label{rem:approx}
The Caveats paper (Appendix~C, Issue~1) shows that when $Y$ is only $\epsilon$-close to
deterministic, all optimizers fall within $O(-\epsilon\log\epsilon)$ of a single corner:
the dead-knob statement of Proposition~\ref{prop:dead} degrades gracefully, and the
near-optimality band of Proposition~\ref{prop:flat} is perturbed by the same order.
\end{remark}

\begin{proposition}[Channel P3, collusion collapse: the variance race and the overfit channel]\label{prop:variance}
For the diagonal Gaussian family $e(z|x) = \mathcal{N}(m(x), \sigma^2 I)$,
$b(z|y) = \mathcal{N}(g(y), \tau^2 I)$ (the implementation's setup, with the
\texttt{logvar} clamp $[-10, 10]$ of \texttt{kl\_divergence}), and a \emph{linear}
classifier $c(z) = Wz + b$ (the implementation's classifier is \texttt{nn.Linear}):
\begin{enumerate}
\item[(i)] for fixed means, $\operatorname{KL}$ is minimized over $\sigma^2$ at
$\sigma^2 = \tau^2$, with value $\tfrac{1}{2\tau^2}\|m - g\|^2$ --- a pure mean-matching
term;
\item[(ii)] the classifier term favors $\sigma^2 \to 0$:
$\mathrm{CE}(c(z), y) = -\log \mathrm{softmax}_y(Wz+b)$ is convex in $z$ (log-sum-exp
of an affine function of $z$, minus a linear term --- the convexity claim requires $c$
affine), so Jensen gives
$\mathbb{E}_z[\mathrm{CE}(c(z), y)] \ge \mathrm{CE}(c(m), y)$. The KL's derivative in
$\sigma^2$ vanishes at $\sigma^2 = \tau^2$, but along the matched line it \emph{resists}
shrinking $\tau^2$ while the means are unmatched
($\partial \operatorname{KL}/\partial \tau^2|_{\sigma^2=\tau^2}
= -\|m-g\|^2/(2\tau^4) < 0$). The joint dynamics therefore runs in two stages: the mean
matching $m \to g$ completes first (with weight $1/\tau^2$, which grows as $\tau$
shrinks), and once $\|m-g\| \to 0$ the variance race is unopposed --- the classifier
drives $\sigma^2 \to 0$ with $\tau^2$ following the moment-matched value --- until the
clamp at $s = e^{-10}$, where the encoder collapses toward the deterministic code
$z \approx g(y)$ and
$\mathrm{ReX} = \tfrac{1}{2s}\mathbb{E}\|m(x) - g(y)\|^2$: a mean-matching force of
weight $1/s = e^{10} \approx 2.2\times10^{4}$ between two \emph{trainable} maps;
\item[(iii)] the certified object and the deployed object separate. By the gap lemma
(Lemma~\ref{lem:gap}) $\mathrm{ReX} = I(X;Z|Y) + \Delta \ge I(X;Z|Y)$, so a small
$\mathrm{ReX}$ \emph{entails} a small residual for the stochastic training code --- the
certificate cannot under-report its own object. What it does not control is the relation
between the training code and the deployed representation. Exact mean matching
($m(x) = g(y)$) makes both objects class-constant, so both residuals vanish; the
scenario of interest is the \emph{near}-match regime, where
$\mathbb{E}\,\|m(X) - g(Y)\|^2 = o(s)$ keeps $\mathrm{ReX}$ small while the deployed
continuous $m$ retains within-class variation, and any non-class-constant continuous
$m$ has $I(X;m(X)|Y) = +\infty$ under the standing non-atomic assumption (the
pathology of deterministic networks, Amjad \&
Geiger 2018; Kolchinsky et al.\ 2017). The variance race therefore moves the
\emph{certified} object toward a noise-dominated class-constant code while the
\emph{deployed} object can keep diverging residual --- and the encoder is fitted to the
backward encoder's per-class targets, memorizing the label code $g(y)$ rather than
being told what residual to erase, while the certificate reports $\mathrm{ReX} \approx
0$ on a code that inference does not use.
\end{enumerate}
\end{proposition}

\begin{proof}
The closed form
$\operatorname{KL} = \tfrac{d}{2}\big[\sigma^2/\tau^2 + \|m-g\|^2/(d\tau^2)
- 1 + \ln(\tau^2/\sigma^2)\big]$ has derivative $\tfrac{d}{2}(1/\tau^2 - 1/\sigma^2)$
in $\sigma^2$, so the variance optimum is $\sigma^2 = \tau^2$, giving (i); at
$\sigma^2 = \tau^2 = s$ the value is $\|m-g\|^2/(2s)$. For (ii): $\mathrm{CE}(c(z), y)$ has Hessian
$\operatorname{diag}(p) - pp^{\top} \succeq 0$ (the categorical Fisher information), so
it is convex in $z$ and Jensen gives the stated inequality; the same Hessian identity makes
$\sigma^2 \mapsto \mathbb{E}_z[\mathrm{CE}(c(z),y)]$ non-decreasing (Gaussians of larger
variance dominate in convex order), strictly increasing in $\sigma^2$ wherever the
classifier is unsaturated, so that term pushes $\sigma^2$ down. The KL's partial derivatives at the
matched point are $\partial_{\sigma^2}\operatorname{KL} = 0$ and
$\partial_{\tau^2}\operatorname{KL} = -\|m-g\|^2/(2\tau^4) < 0$, so the KL pulls
$\tau^2$ up until the means match, while the mean-matching gradient on the encoder is
$(m - g)/\tau^2$; once $\|m-g\| \to 0$ the race to the clamp is unopposed. For (iii): the
gap identity (Lemma~\ref{lem:gap}) forces the certified residual and the certificate to
vanish together; the divergence attaches to the deployed mean $m(x)$ whenever it is not
class-constant (Amjad \& Geiger; Kolchinsky et al.). \hfill\qed
\end{proof}

\begin{remark}[Status of the two-stage account]\label{rem:p3-status}
Parts~(ii) and~(iii) are a stationarity-and-direction-of-force argument, not a convergence
proof: the ``two stages'' presume a time-scale separation (mean matching completes before
the variance race) that is not established here. The account is directly falsifiable in
this codebase: log the mean posterior log-variance and the fraction of units at the clamp
$s = e^{-10}$ during CEB training. \emph{We ran the test.} On MNIST, across all six
$\beta$ settings and 5 seeds (Section~\ref{sec:ablation} of the main paper), the trained
posterior and prior log-variances of CEB remain near their initialization ($\sigma^2 =
\tau^2 = 1$, mean $\log\sigma^2$ within $0.74$ of zero, clamp fraction $0.0$ at every
epoch of every run). At the early-stopping time scale of the experimental protocol, the
variance race therefore does not materialize, and
Proposition~\ref{prop:variance}(iii) is empirically void for those runs: Part~(iii)
stands as a sharp hypothesis about what longer training or different schedules would
produce, the force analysis of (i)--(ii) as the only proved content, and the
measurement itself as the finding.
\end{remark}

\begin{remark}[Train--evaluation mismatch]\label{rem:mismatch}
The classifier is trained on the sampled code $z$ and evaluated on the mean $m(x)$
(Fischer 2020, Appendix~A.1). Proposition~\ref{prop:variance}(ii) is the force that
closes this gap by destroying the noise altogether ($\sigma \to 0$), at which point the
bottleneck's stochasticity --- the object the bound is defined on --- disappears.
\end{remark}

\begin{proposition}[Channel P4, certificate-object mismatch: neither bound is an upper bound on the deployed representation]\label{prop:object}
(a)~\emph{CEB}. The classifier consumes the sampled code $Z$ at training and the mean
$\mu(X)$ at evaluation (Fischer 2020, Appendix~A.1). At the variance-optimal posterior,
$\sigma(x) \equiv \sqrt{S_\theta(y)}$ (the backward encoder's per-class standard
deviation) is class-constant, so $Z = \mu(X) + \sqrt{S_\theta(Y)}\,\varepsilon$ is a
noisy function of $\mu(X)$ given $Y$; by the DPI,
\begin{equation}
I(X; Z | Y) \;\le\; I(X;\, \mu(X) | Y) ,
\label{eq:dpi}
\end{equation}
strictly so whenever the noise destroys within-class detail. Hence CEB's certificate
$\mathrm{ReX} \ge I(X;Z|Y)$ is a bound on the \emph{training code}, not on the residual
of the representation deployed at inference. The gap is not a small correction: whenever
$\mu$ is a continuous map that is not class-constant, $I(X;\mu(X)|Y) = +\infty$ under
the standing non-atomic assumption (Amjad \& Geiger 2020; Kolchinsky et al.\ 2017), so
no finite quantity --- CEB's certificate included --- bounds the deployed residual from
above. The CE term presses the noise
toward $0$ (by Jensen, $\mathbb{E}_z[\mathrm{CE}(c(z),y)] \ge \mathrm{CE}(c(\mu),y)$),
but no finite $\sigma$ closes the gap.
(b)~\emph{FGIB}. The inference path is $h$ at both training and evaluation; the side
noise of Proposition~\ref{prop:inverse} is not part of any inference path. By the DPI
applied to the same perturbed-variable construction,
\begin{equation}
I(X;\, h + \tau (A^{\top}A)^{-1/2}\varepsilon \,|\, Y)
\;\le\; I(X; H | Y) ,
\label{eq:dpi-f}
\end{equation}
so the direction of the bound is the same as CEB's: $\mathrm{ReX}_F$ upper-bounds the
noisy side-path object, and \emph{no} finite value upper-bounds the deployed residual.
The point-anchor limit $\tau \to 0$ does not repair this: it is a change of model whose
bound value $\|\mu - a v_y\|^2/(2\tau^2)$ typically diverges
(Remark~\ref{rem:limit}), and even the limiting object needs $A$ well-conditioned
(Proposition~\ref{prop:conditioning}(iii)).
(c)~\emph{Conclusion}. In the pure variational value sense (same encoder, same Gaussian
family), CEB's bound is tighter (Proposition~\ref{prop:ordering}). In the object sense
the two are not ordered: both are valid upper bounds on their respective noisy codes, and
neither bounds the deterministic deployed representation. What remains structural --- not
a tightness claim --- is \emph{where the gap lives}: CEB's perturbation is inside the
inference path and is driven by the prediction term (so removing it changes the model),
while FGIB's perturbation lives on a side path that inference never reads and is
controlled by side-path parameters alone. The latter is a cleaner separation of
regularizer from model, not a tighter bound.
\end{proposition}

\begin{remark}[The bound is a training instrument, not an inference one]\label{rem:training-role}
No method computes its bound at inference: the bound's roles are \emph{training-time}
--- a regularizer (a differentiable surrogate for the intractable information term) and
a meter (Fischer 2020, p.~7: ``observing ReX tells us exactly how many bits we are from
the MNI point''). In the $Z \perp X \mid Y$ limit CEB's mean is also class-constant
($m(x) = \mathbb{E}[Z|x] = \mathbb{E}[Z|y]$), so the mismatch of
Proposition~\ref{prop:object}(a) is not a limit failure: it is a finite-$\beta$
\emph{under-reporting} --- the regularizer controls, and the meter reports, the training
code, which is strictly more compressed than the deployed representation $\mu(X)$, by a
positive, trained deficit that cannot be measured from the certificate. FGIB's
side-path noise is a known side-path parameter and is removable in principle
($\tau \to 0$ changes the side path, never the deployed model), but for finite
$\tau$ its regularizer and meter act on a \emph{perturbed copy} of the deployed
object, with the perturbation size governed by the conditioning of $A$
(Proposition~\ref{prop:conditioning}(iii)).
The one genuine inference-time use of a certificate value is as a detection score
(Fischer 2020, RG3, uses the marginal rate $R$): CEB's label-conditioned
$\mathrm{ReX}$ is not even computable at inference (the label is unknown --- Fischer
trains a separate marginal for this purpose), while FGIB's geometry gives a label-free
score, the distance to the nearest anchor
$\min_y \operatorname{KL}(q(z|x)\|\mathcal{N}(a v_y, \tau^2 I))$, with no extra network
(the minimum is an internal enumeration over the $K$ fixed anchor buffers --- the single
QR orthonormal frame --- and consults no ground-truth label, no more than softmax
consults labels when it enumerates the $K$ logits). Whether this score is useful for OoD
detection is a hypothesis; nothing here proves it.
\end{remark}

\textbf{The problem addressed by FGIB.}
The four channels are addressed by the FGIB design as follows; two are closed by
construction, one holds conditionally, and one is only partially addressed. The precise
scope of each claim is documented in Part~II and the open items are collected in
Section~\ref{sec:limits}.
\begin{itemize}
\item \emph{Dead knob and cheap label axis} (Props.~\ref{prop:dead}, \ref{prop:flat}):
FGIB's compression term is the anchor gap $\tfrac{1}{2\tau^2}(s-a)^2$, quadratic in the
alignment scalar, and \emph{at fixed classifier scale} the Lagrangian
$\mathrm{CE}(s) + \tfrac{\beta}{2\tau^2}(s-a)^2$ is strictly convex with a unique interior
minimizer per $\beta$ --- the $\beta$-sweep is one-to-one onto the whole trade-off curve
(Props.~\ref{prop:sweep}, \ref{prop:curve}). With the scale free the infimum is $0$ for
every $\beta$; Remark~\ref{rem:c-scale} states the condition and the architectural fix.
\item \emph{Blind certificate} (Prop.~\ref{prop:flat}): with the fixed anchor
$r(z|y) = \mathcal{N}(a v_y, \tau^2 I)$ the gap
$\mathbb{E}_y[\operatorname{KL}(q(z|y)\,\|\,r(z|y))]$ can shrink only if the
\emph{posterior} moves toward the anchors --- the reference itself cannot chase the
encoder, so the force is persistent (Prop.~\ref{prop:ordering} and
Remark~\ref{rem:force}). The side head $A$, however, is a second trainable party in that
matching, and it is not part of the inference path
(Remark~\ref{rem:A-channel}).
\item \emph{Variance collapse} (Prop.~\ref{prop:variance}): the classifier consumes the
deterministic $h$ at both training and evaluation, so $\sigma^{2*} = \tau^2$ for every $\beta$
(the variance decouples from the prediction term; Lemma~\ref{lem:lagrangian}), the
side path keeps a finite covariance and the bound stays well-defined
(Remark~\ref{rem:limit}), and the trivial collapse $\mu \equiv 0$ pays
the fixed positive price $\operatorname{KL} = \tfrac12 a^2 > 0$ for every class --- the
degenerate corner is excluded by geometry rather than by the optimizer's good behavior.
This channel is closed cleanly, by construction.
\item \emph{Certificate-object mismatch} (Prop.~\ref{prop:object}): the classifier
consumes the deterministic $h$ at both training and evaluation --- there is no
training-code/evaluation-mean split --- but the bound is an upper bound on a
\emph{noisy side-path copy} of $h$, not on $I(X;H|Y)$, which no finite certificate
bounds. This channel is \emph{not} closed; what FGIB buys is that the perturbation is
side-path-only, and its size is governed by the conditioning of $A$
(Prop.~\ref{prop:conditioning}(iii)).
\end{itemize}

% =====================================================================
\section*{Part II. The FGIB theory}
% =====================================================================

\textbf{FGIB notation.}
The deterministic main path computes $h = f_{\mathrm{enc}}(X)$ and the classifier
consumes $h$ directly; the side path parameterizes $q(z|x) = \mathcal{N}(\mu(x),
\Sigma(x))$ with \emph{linear} mean head $\mu(x) = A\,h(x)$ for
$A \in \mathbb{R}^{z \times z}$ (the theory assumes $z = d$ and $A$ invertible; the
inverse transforms of Propositions~\ref{prop:fgib-ceb} and~\ref{prop:inverse} are
defined under this assumption), and a separate logvar head for $\Sigma(x)$. The
side path is \emph{not part of inference}: at evaluation the model returns the classifier
output and never evaluates the mean head (so $A$ is a trainable module that exists only
inside the regularizer --- structurally the same position CEB's backward encoder
occupies on the other side of the KL). The label-conditional prior is the
\emph{fixed} anchor
$r(z|y) = \mathcal{N}(a\,v_y, \tau^2 I)$ with \emph{anchor variance} $\tau^2 > 0$
(default $1.0$): the $v_y$ are orthonormal directions obtained by
QR-decomposing a random matrix ($K \le z$; the orthonormal factor of the random matrix),
scaled by the anchor scale $a$ (default $4.0$), stored as
non-trainable buffers. At the variance optimum of the KL the posterior variance matches
the anchor variance, $\sigma^{2*} = \tau^2$; we write $\sigma_p^2$ for this common
variance where the distinction does not matter. The orthogonality of the method lies in
this fixed frame and in the equilibrium it induces --- the KL drives the
class-conditional means $\mu_y = A\,\mathbb{E}[h|y]$ toward the orthonormal anchors
$a v_y$, so the transformed class representations become mutually orthogonal (the
aligned family below) --- while $A$ itself remains an unconstrained linear map. For
regression, fixed random-Fourier anchors
$\mu_p(y) = a\sqrt{2/z}\,\cos(2\pi\omega y + b)$ with frozen
$\omega \sim \mathcal{N}(0, 2^2)$, $b \sim U[0, 2\pi)$ replace the per-class table
($\|\mu_p(y)\| \approx a$). The loss is
$L_{\mathrm{FGIB}} = \mathbb{E}[\mathrm{CE}(c(h), y)]
+ \beta\,\mathbb{E}[\operatorname{KL}(q(z|x)\,\|\,r(z|y))]$, with gradients reaching the
main path only through the shared encoder; the zero-initialized logvar head fixes
$\sigma^2 = 1$ at initialization, which equals $\tau^2$ in the default configuration
$\tau^2 = 1$ (for general $\tau$ the theory's variance optimum $\sigma^2 = \tau^2$ is
reached by the logvar head, whose initialization is a scheduling choice), so training
starts at $\operatorname{KL} \approx \tfrac{1}{2\tau^2} a^2$.

\emph{Implementation caveat (theory--code layering).} The theory works with a fixed
orthonormal factor, an invertible linear side head $T(h) = Ah$, and per-class Gaussian
priors with anchor variance $\tau^2$. The code implements an \emph{affine} head
(\texttt{nn.Linear} with bias) --- which the theory absorbs as an idealization, the bias
being immaterial when $A$ is learned jointly --- and does not enforce orthogonality,
invertibility, or conditioning constraints on $A$. The guarantees below (orthonormal
frames, invariance, conditioning) are properties of the idealized model; the
implementation approximates them, and the guarantees do not transfer automatically
(Remarks~\ref{rem:A-channel}, \ref{rem:rank}).
\textbf{Distributional standing assumption.} Several results state that a deterministic
continuous representation carries infinite conditional information. Those statements
assume the standard continuous setting: $X \mid Y$ is non-atomic, and the induced
representation $H(X) \mid Y$ is a non-atomic continuous variable. For the finite
empirical distribution of the training data no such conclusion is claimed; there the
quantities are finite and only their growth with sample size is governed by these
results.

% ---------------------------------------------------------------
\subsection{The equivalence: the side path is a CEB bottleneck on the main path}
% ---------------------------------------------------------------

\textbf{Motivation.}
Stochastic bottlenecks such as VIB and CEB are \emph{stochastic at training time but
deterministic at evaluation time}: the classifier is trained on sampled codes
$z \sim q(z|x)$ yet evaluated on the mean $\mu_q(x)$ (Fischer 2020, Appendix A.1), so
its training and evaluation distributions differ. In the same models the cross-entropy
gradient and the KL gradient act on the same heads: the classification term can exploit
the variance (by Jensen, $\mathbb{E}_z[\mathrm{CE}(c(z), y)] \ge \mathrm{CE}(c(\mu_q),
y)$, so loosening the noise relaxes the upper bound), while the KL term pushes the
posterior back toward the prior. Finally, the regularizer constrains the noisy code $z$,
whereas inference consumes the noise-free mean --- the object being regularized is not
the object being used. FGIB resolves all three mismatches by design: the classifier
consumes the deterministic representation $h$ at both training and evaluation time, the
two objectives act on disjoint parameter sets (the classification gradient does not
enter the side heads, and the KL gradient reaches the main path only through the shared
encoder), and, by Proposition~\ref{prop:fgib-ceb} below, the side-path regularizer is,
\emph{at each fixed value of the side head}, a CEB bottleneck read in the coordinates in
which the posterior is centered at the representation that inference actually uses
(Remark~\ref{rem:A-channel} states why this is weaker than a fixed constraint on $h$).
Confining the stochasticity to a side path also keeps the variational objective
well-defined: for continuous deterministic representations the residual information
diverges (Amjad \& Geiger; Kolchinsky et al.), so a purely deterministic bottleneck
admits no finite variational bound, while the finite side-path covariance of FGIB does
(Remark~\ref{rem:limit}). Whether the robustness benefits of the stochastic bottleneck
survive the deterministic inference path is the empirical question studied in the
experiments.

\begin{lemma}[Linear invariance of the Gaussian KL]\label{lem:kl-invariance}
Let $q = \mathcal{N}(\mu_q, \Sigma_q)$ and $p = \mathcal{N}(\mu_p, \Sigma_p)$ on
$\mathbb{R}^{z}$, and let $A \in \mathbb{R}^{z \times z}$ be invertible. Define the
common pushforwards $q' = A_\# q = \mathcal{N}(A\mu_q, A\Sigma_q A^\top)$ and
$p' = A_\# p$. Then $\operatorname{KL}(q'\,\|\,p') = \operatorname{KL}(q\,\|\,p)$.
\end{lemma}

\begin{proof}
With $d = z$, the closed form reads
$\operatorname{KL}(q\,\|\,p)
= \tfrac{1}{2}\big[\operatorname{tr}(\Sigma_p^{-1}\Sigma_q)
+ (\mu_q - \mu_p)^\top \Sigma_p^{-1}(\mu_q - \mu_p)
- d + \ln|\Sigma_p| - \ln|\Sigma_q|\big]$.
Under the common map $A$ the three non-constant terms are unchanged:
\begin{align*}
\operatorname{tr}\big((A\Sigma_p A^\top)^{-1}(A\Sigma_q A^\top)\big)
&= \operatorname{tr}(A^{-\top}\Sigma_p^{-1}A^{-1} A\,\Sigma_q A^\top)
 = \operatorname{tr}(A^{-\top}\Sigma_p^{-1}\Sigma_q A^\top) \\
&= \operatorname{tr}(\Sigma_p^{-1}\Sigma_q A^\top A^{-\top})
 = \operatorname{tr}(\Sigma_p^{-1}\Sigma_q), \\[2pt]
(A\Delta\mu)^\top (A\Sigma_p A^\top)^{-1} (A\Delta\mu)
&= \Delta\mu^\top A^\top A^{-\top}\Sigma_p^{-1}A^{-1}A\,\Delta\mu
 = \Delta\mu^\top \Sigma_p^{-1}\Delta\mu, \\[2pt]
\ln|A\Sigma_p A^\top| - \ln|A\Sigma_q A^\top|
&= \big(\ln|\Sigma_p| + 2\ln|A|\big) - \big(\ln|\Sigma_q| + 2\ln|A|\big)
 = \ln|\Sigma_p| - \ln|\Sigma_q|,
\end{align*}
where $\Delta\mu = \mu_q - \mu_p$ and the trace step uses the cyclic property. \hfill\qed
\end{proof}

\begin{lemma}[Invariance of the CEB objective and the MNI point]\label{lem:mi-invariance}
For any bijective measurable map $\phi$,
$I(X;\phi(Z)|Y) = I(X;Z|Y)$ and $I(Y;\phi(Z)) = I(Y;Z)$.
Consequently the CEB objective (Eq.~5) is invariant under bijections of $Z$, and
$Z$ attains the MNI point (Eq.~1) if and only if $\phi(Z)$ does.
\end{lemma}

\begin{proof}
Mutual information is invariant under bijective transformations of either argument
(data processing inequality applied in both directions). \hfill\qed
\end{proof}

\begin{proposition}[The FGIB side path is a CEB bottleneck on the main path, at each fixed side head]\label{prop:fgib-ceb}
Let $z = d$ and let the learned linear head $T(h) = A h$ be invertible (unconstrained,
so a.s.\ so at initialization). The \emph{fixed}
prior $r(z|y) = \mathcal{N}(\mu_p(y), \Sigma_p(y))$ --- for FGIB, the orthonormal anchor
table $\mu_p(y) = a\,v_y$, $\Sigma_p = \tau^2 I$ --- is, by construction, the image
under $T$ of the canonical prior
$r'(\tilde{z}|y) = \mathcal{N}(A^{-1}\mu_p(y),\, A^{-1}\Sigma_p(y) A^{-\top})$
(since $A A^{-1} = I$), and the posterior $q(z|x) = \mathcal{N}(A\,h(x), \Sigma_q(x))$
is the image of
$q'(\tilde{z}|x) = \mathcal{N}(h(x), A^{-1}\Sigma_q(x) A^{-\top})$.
\emph{At a fixed value of $A$}, the prior is never transformed by training:
$A^{-1}$ only reads the fixed anchor geometry in the coordinates in which the posterior
is centered at $h$; during training $A$ moves, so the pulled-back prior moves with
it (Remark~\ref{rem:A-channel}).
By KL invariance under the common invertible linear map $T^{-1} = A^{-1}$
(Lemma~\ref{lem:kl-invariance}), pointwise for every $(x, y)$,
\begin{equation}
\operatorname{KL}\big(q(z|x)\,\|\,r(z|y)\big)
= \operatorname{KL}\big(q'(\tilde{z}|x)\,\|\,r'(\tilde{z}|y)\big),
\label{eq:fgib-ceb}
\end{equation}
That is, at each fixed side head the side-path regularizer is pointwise identical to a
CEB bottleneck with a backward encoder \emph{frozen at the current head}, applied to
a noisy version of the main-path representation,
$\tilde{z} = h + \tilde{\epsilon}$ with
$\tilde{\epsilon} \sim \mathcal{N}(0, A^{-1}\Sigma_q A^{-\top})$,
using the canonical prior $r'$ as that backward encoder. For FGIB the canonical backward
encoder is the frame $\mathcal{N}(a\,A^{-1} v_y, \tau^2 A^{-1}A^{-\top})$, whose
mean anchors $A^{-1}(a v_y)$ are orthonormal at scale $a$ in the induced metric
$A^{\top}A$
($\langle A^{-1}v_i, A^{-1}v_j\rangle_{A^{\top}A} = \delta_{ij}$): the orthogonality
enters through the fixed frame $\{v_y\}$ and the equilibrium it induces, not through a
constraint on $A$.
Hence the regularizer's gradient with respect to $h$ (through the shared encoder)
coincides, up to the fixed linear map $A^{\top}$, with the gradient that a CEB model
applies to its representation mean.
\end{proposition}

\begin{proof}
The map $z \mapsto \tilde{z} = A^{-1}z$ is a bijection, and
$q' = (A^{-1})_\# q$, $r' = (A^{-1})_\# r$:
$A^{-1}A\,h = h$ and $A^{-1}\Sigma_q A^{-\top}$ give the canonical posterior, and
$A^{-1}\mu_p$, $A^{-1}\Sigma_p A^{-\top}$ give the canonical prior, with
$A\,(A^{-1}\mu_p) = \mu_p$ and $A\,(A^{-1}\Sigma_p A^{-\top})\,A^{\top}
= \Sigma_p$, so $r$ is exactly the image of $r'$ under $T$.
Equation~\eqref{eq:fgib-ceb} is Lemma~\ref{lem:kl-invariance} applied to $A^{-1}$.
Sampling $\tilde{z} \sim q'$ is equivalent to
$\tilde{z} = h + \tilde{\epsilon}$ with the stated noise.
Combined with the variational bound of Fischer (2020, Eqs.~9--11), the regularizer
upper-bounds the residual information of the noisy main-path representation:
$I(X;\tilde{Z}|Y) \le \mathbb{E}[\operatorname{KL}(q'(\tilde{z}|x)\,\|\,r'(\tilde{z}|y))]$.
\hfill\qed
\end{proof}

\begin{remark}[Trainable-implementation caveat: the side head is a second trainable party in the matching]\label{rem:A-channel}
The identity of Proposition~\ref{prop:fgib-ceb} holds pointwise at each value of the
side head, so the asymmetry advertised against CEB --- one trainable side versus two ---
is weaker than it looks. In the idealized theory $A$ may be treated as fixed
(orthonormal or any fixed invertible map); in the \emph{trainable implementation} $A$ is
free and \emph{unused at inference}:
like CEB's backward encoder $b(z|y)$, it lives inside the KL and can lower it without
moving the deployed representation $h$. Since $h$ is linear-separable by the jointly
trained classifier, the linear fit of the anchor target $a v_y$ on $h$ fits well, and
much of that fit can be absorbed by $A$ alone. The variance term
$\tfrac{z}{2}(\sigma^2/\tau^2 - 1 - \ln(\sigma^2/\tau^2))$ does not depend on $A$ at
all, so nothing in
the regularizer resists $A$ drifting toward a projector onto the $K$-dimensional
between-class subspace with a vanishing component on within-class directions. Full
collapse is blocked --- with invertible $A$, $\operatorname{KL} = 0$ forces $h$ itself
to be class-constant --- but the approach to the rank-$K$ regime is not blocked, and
that regime is exactly where the conditioning failure and the perturbation blow-up of
Propositions~\ref{prop:conditioning} and~\ref{prop:inverse} live. This is the one
channel of Part~I that FGIB has not closed; Section~\ref{sec:limits} collects it.
\end{remark}

\begin{remark}[Family closure]\label{rem:closure}
Lemma~\ref{lem:kl-invariance} holds for the underlying Gaussian distributions. Within a
restricted variational family the equality of Proposition~\ref{prop:fgib-ceb} is exact
only if the family is closed under the inverse map $A^{-1}$: the full-covariance family
is closed (and is the parameterization used by Fischer 2020); the diagonal family is not
--- for diagonal $q, r$ the transformed pair is diagonal only if $A^{-1}$ preserves
diagonality (e.g., permutations and sign flips) and approximate otherwise. The diagonal
implementation of the codebase is therefore covered at the distribution level, with the
family approximation confined to the noise covariance $A^{-1}\Sigma_q A^{-\top}$
(Remark~\ref{rem:orth}).
\end{remark}

\begin{remark}[Rank deficiency]\label{rem:rank}
If $z < d$ or $A$ is rank-deficient, only the projection of $h$ onto the row space of
$A$ (the image of $A^{\top}$ in $h$-space) receives gradient from the regularizer; the
null-space component is unconstrained.
In the default configuration $z = d$ and a randomly initialized $A$ is almost surely
invertible, but invertibility carries no conditioning margin --- near-singularity is
enough to produce the effects of Proposition~\ref{prop:conditioning}(iii), and it is the
direction the gap gradient points in (Remark~\ref{rem:A-channel}). If $A \in
\mathbb{R}^{z \times d}$ has full column rank ($z \ge d$), a left inverse $B$ with
$BA = I$ exists and the data-processing inequality for KL gives, per $(x,y)$,
$\operatorname{KL}(q\|r) \ge \operatorname{KL}(B_\# q\,\|\,B_\# r)
\ge I(X;\,BZ\,|\,Y)$ --- a valid lower-level bound, but one whose equality
(invariance) is lost, and whose pulled-back prior again depends on $B$. The exact
invariance argument of Proposition~\ref{prop:fgib-ceb} genuinely requires the invertible
linear map $T(h) = Ah$ with $z = d$.
\end{remark}

\begin{remark}[Noise-free limit]\label{rem:limit}
The equivalence of Proposition~\ref{prop:fgib-ceb} should be understood at the level of
the variational objective, not as a mutual-information limit: as the side-path noise
vanishes, $I(X;Z|Y)$ diverges for continuous deterministic $H$ under the standing
non-atomic assumption (the pathology of
deterministic networks discussed by Amjad \& Geiger and Kolchinsky et al.), and the
variational KL bound becomes loose. Keeping the side-path covariance finite --- without
injecting noise into the inference path --- is precisely what keeps the regularizer
well-defined; this is the design rationale of FGIB.
\end{remark}

% ---------------------------------------------------------------
\subsection{The sweep}
% ---------------------------------------------------------------

\textbf{The degeneracy to be addressed (restated in our terms).}
When $Y = f(X)$, the IB curve saturates the data-processing bound and is piecewise
linear, $F(r) = \min\{r, H(Y)\}$ (Caveats paper, Section~3). Two consequences (their
Caveat~1): (i)~by the DPI, $L^{\mathrm{IB}}_\beta(T) = I(Y;T) - \beta I(X;T)
\le (1-\beta)I(Y;T) \le (1-\beta)H(Y)$, and $T_{\mathrm{copy}} := f(X) = Y$ attains the
bound for \emph{every} $\beta \in [0,1]$: all $\beta \in (0,1)$ are pinned at the single
corner point $\langle H(Y), H(Y)\rangle$ of the curve; (ii)~conversely, at $\beta = 1$
every point of the increasing part simultaneously maximizes $L^{\mathrm{IB}}_1$. The
mechanism is exposed by writing $\max_T L^{\mathrm{IB}}_\beta = \max_r [F(r) - \beta r]$
(the Legendre transform of $-F$): on the linear increasing part every point has the
\emph{same} derivative $F'(r) = 1$, hence the same supporting slope --- a Lagrangian of
the form ``prediction $-\beta\cdot{}$compression'' sweeps a Pareto curve only if each
point of the curve has a \emph{unique} supporting slope, i.e.\ the curve is strictly
concave. The dIB variant fails identically: on the hard-clustering family $T = g(Y)$ one
has $H(T|Y) = 0$, hence $I(Y;T) = H(T)$, so $L^{\mathrm{dIB}}_\beta = (1-\beta)H(T)$
--- a \emph{monotone} function of the single compression scalar $H(T)$, maximized by the
finest partition for every $\beta < 1$. The squared fixes repair this by making the
objective strictly concave in the compression scalar: $F'(r)/(2r) = \beta$ is strictly
decreasing, so each $\beta$ selects a unique $r$ (Eq.~8 argument), and squared-dIB has
the unique maximizer $H(T) = 1/(2\beta)$ (their Eq.~A20). By Proposition~\ref{prop:dead}
this degeneracy is inherited verbatim by CEB ($\beta_{\mathrm{IB}} = 1+\gamma > 1$),
which is the first channel of Part~I.

\textbf{Reduced model of FGIB.}
The KL depends on the side mean $\mu$ alone, and the classifier term is invariant under
the coordinate change $\tilde w_j = A^{-\top} w_j$ (since $w_j \cdot h =
\tilde w_j \cdot \mu$): with $A$ invertible ($z = d$ in the default configuration, so a
random $A$ is a.s.\ invertible), every equilibrium in $(h, w)$ coordinates maps
bijectively to one in $(\mu, \tilde w)$ coordinates. We therefore analyze the reduced
model $\mu(x) = h(x)$ without loss of generality.
Consider the \emph{class-symmetric aligned family}: the encoder maps every input of
class $y$ to $h_y = s\,v_y$ ($s \ge 0$), the classifier is held at the symmetric
configuration $w_j = c\,v_j$, $b_j = 0$ with fixed scale $c > 0$, and the posterior is
$q_y = \mathcal{N}(s v_y, \sigma^2 I)$. The family is the natural equilibrium manifold of
the two forces: the KL gradient
$\nabla_h \operatorname{KL} = (\mu - a v_y)/\tau^2$ acts exactly
along $v_y$, and the classification force on the class center acts along $v_y$ at the
symmetric configuration; the ansatz is a \emph{working approximation}: it captures the
direction of both forces, and is motivated as exact in the $\beta$-dominant regime, but
that exactness is not proved here --- the cross terms are of order $O(c/\beta)$ and are
idealized away (see Remark~\ref{rem:caveats}).

\begin{lemma}[FGIB Lagrangian on the aligned family]\label{lem:lagrangian}
On the aligned family,
\begin{equation}
\operatorname{KL}(s, \sigma^2)
= \tfrac{z}{2}\big(\sigma^2/\tau^2 - 1 - \ln(\sigma^2/\tau^2)\big)
+ \tfrac{1}{2\tau^2}(s-a)^2,
\qquad
\mathrm{CE}(s) = \ln\!\big(1 + (K-1)e^{-cs}\big),
\label{eq:kl-ce}
\end{equation}
and the optimizer over $\sigma^2$ is $\sigma^2 = \tau^2$ for \emph{every} $\beta$, since
the CE term does not depend on $\sigma^2$ (the classifier consumes the deterministic
$h$, not $z$).
\end{lemma}

\begin{proof}
With $h_y = s v_y$ and $w_j = c v_j$, the logits are $c\,\langle v_j, s v_y\rangle
= cs\,\delta_{jy}$, so $p_y = e^{cs}/(e^{cs} + K - 1)$ and
$\mathrm{CE} = -\ln p_y = \ln(1 + (K-1)e^{-cs})$.
For the KL, the diagonal closed form with $\Sigma_q = \sigma^2 I$,
$\Sigma_p = \tau^2 I$ reads
$\tfrac{1}{2}\sum_d\big[\sigma^2/\tau^2 + (\mu_d - a v_{yd})^2/\tau^2 - 1
- \ln(\sigma^2/\tau^2)\big]$, giving
Eq.~\eqref{eq:kl-ce} since $\|v_y\| = 1$. The variance term is non-negative with a
unique minimum at $\sigma^2 = \tau^2$, and $\mathrm{CE}$ is independent of $\sigma^2$.
\hfill\qed
\end{proof}

\begin{proposition}[Unique interior optimum per $\beta$ at fixed classifier scale --- the sweep]\label{prop:sweep}
\emph{At a fixed classifier scale $c$}, for every $\beta > 0$, the FGIB Lagrangian on
the aligned family,
\begin{equation}
L_\beta(s) := \ln\!\big(1 + (K-1)e^{-cs}\big) + \tfrac{\beta}{2\tau^2}(s-a)^2,
\label{eq:fgib-lag}
\end{equation}
is strictly convex in $s$ and has a unique minimizer $s^*(\beta) > a$, given implicitly
by
\begin{equation}
\frac{\beta}{\tau^2}\,(s-a) = c\,\varphi(cs),
\qquad
\varphi(u) := \frac{(K-1)e^{-u}}{1 + (K-1)e^{-u}} \in (0,1) .
\label{eq:stationarity}
\end{equation}
The map $\beta \mapsto s^*(\beta)$ is continuous and strictly decreasing, with
$s^*(\beta) \to a$ as $\beta \to \infty$ and $s^*(\beta) \to \infty$ as $\beta \to 0$.
Consequently the map $\beta \mapsto \big(\operatorname{KL}(s^*), \mathrm{CE}(s^*)\big)$
is one-to-one onto the entire curve
$\{\big(\tfrac{1}{2\tau^2}(s-a)^2,\; \ln(1+(K-1)e^{-cs})\big) : s \in (a, \infty)\}$:
no two
$\beta$ share the same optimum, and every interior point of the trade-off curve is
attained --- the exact failure mode of $L^{\mathrm{IB}}_\beta$ (all $\beta$ pinned at
$T_{\mathrm{copy}}$) is absent. The fixed-$c$ condition is load-bearing and is analyzed
in Remark~\ref{rem:c-scale}.
\end{proposition}

\begin{proof}
Compute $\mathrm{CE}'(s) = -c\varphi(cs)$ and
$\mathrm{CE}''(s) = c^2 (K-1) e^{-cs}/(1 + (K-1)e^{-cs})^2 > 0$, so CE is strictly
convex and strictly decreasing in $s$; the compression term $\tfrac{1}{2\tau^2}(s-a)^2$
is strictly convex. Hence $L_\beta''(s) = \mathrm{CE}''(s) + \beta/\tau^2 > 0$:
$L_\beta$ is
strictly convex with at most one minimizer, characterized by
$0 = L_\beta'(s) = -c\varphi(cs) + \beta(s-a)/\tau^2$, i.e.\
Eq.~\eqref{eq:stationarity}.
The left side $\beta(s-a)/\tau^2$ is strictly increasing from $-\beta a/\tau^2$
(at $s=0$) through $0$
(at $s=a$) to $+\infty$; the right side $c\varphi(cs)$ is strictly decreasing in $s$
with $c\varphi(cs) \in (0, c)$. The unique crossing therefore satisfies $s^* > a$.
Differentiating Eq.~\eqref{eq:stationarity} with respect to $\beta$ gives
$(s-a)/\tau^2 + (\beta/\tau^2) s' = c^2\varphi'(cs)\,s'$ with $\varphi' < 0$, so
$s' = -(s-a)/(\beta - c^2\tau^2\varphi'(cs)) < 0$: $s^*$ is strictly decreasing.
The limits follow by monotonicity: as $\beta \to \infty$ the left side can only match
the bounded right side when $s \to a$; as $\beta \to 0$ the crossing requires
$\varphi(cs)\to 0$, i.e.\ $s \to \infty$. \hfill\qed
\end{proof}

\begin{proposition}[The trade-off curve is strictly concave in the compression--prediction plane]\label{prop:curve}
On the curve of Proposition~\ref{prop:sweep} (fixed classifier scale $c$),
\begin{equation}
\frac{d\,\mathrm{CE}}{d\,\operatorname{KL}}
= \frac{\mathrm{CE}'(s)}{(s-a)/\tau^2}
= -\frac{\tau^2\,c\varphi(cs)}{s-a},
\label{eq:slope}
\end{equation}
which is strictly increasing from $-\infty$ (as $s \to a^+$) to $0^-$ (as $s \to
\infty$). Thus the curve $\operatorname{KL} \mapsto -\mathrm{CE}$ (compression versus
achievable prediction, on the KL--CE surrogate plane of
Remark~\ref{rem:surrogate}) is strictly concave, and each of its points has a \emph{unique}
supporting slope $-\beta$. By the criterion of the Caveats paper (their Fig.~2 right
panel and the $F'(r)/(2r)$ argument of Eq.~8), this is precisely the property that makes
the Lagrangian $-\mathrm{CE} - \beta\,\operatorname{KL}$ sweep the full \emph{surrogate}
curve on the aligned family at fixed scale: the convexity of the compression term
supplies the strict-concavity condition \emph{without squaring it}. This is a statement
about the KL--CE surrogate of Proposition~\ref{prop:sweep}, not about the Caveats
information plane (Remark~\ref{rem:surrogate}).
\end{proposition}

\begin{proof}
$\varphi$ is positive and strictly decreasing and $s-a$ is positive and strictly
increasing, so $\tau^2 c\varphi(cs)/(s-a)$ is strictly decreasing in $s$; the signs give
the stated limits. The stationarity condition Eq.~\eqref{eq:stationarity} is exactly
$-d\mathrm{CE}/d\operatorname{KL} = \beta$, i.e.\ $\beta$ is the supporting slope at the
optimum. \hfill\qed
\end{proof}

\begin{remark}[What the sweep sweeps: a KL--CE surrogate curve, not the IB curve]\label{rem:surrogate}
Propositions~\ref{prop:sweep} and~\ref{prop:curve} prove strict convexity of the
FGIB Lagrangian and one-to-one parameterization by $\beta$ \emph{on the class-symmetric
aligned family, at fixed classifier scale}, in the plane
$\big(\operatorname{KL}_{\mathrm{anchor}},\, \mathrm{CE}\big)$. That plane is not the
Caveats paper's IB plane $\big(I(X;T),\, I(Y;T)\big)$, and the claim does not extend to
it. On the aligned family the posterior is class-constant by construction, so
$I(X;Z|Y) = 0$ for every $s$, while the deterministic class prototypes $h_y = s v_y$
already retain the label completely for any $s > 0$ ($I(Y;H) = H(Y)$); what varies
along the sweep is the anchor mismatch and the classifier margin, not a proved
compression of mutual information. Whether the deployed representation's residual
$I(X;H|Y)$ moves monotonically with $\beta$ is a separate, empirical question; the
theory guarantees a strictly concave surrogate, not a recovered IB curve. This
qualification is carried into Section~\ref{sec:limits}.
\end{remark}

\begin{remark}[The classifier scale is load-bearing, and the current implementation leaves it free]\label{rem:c-scale}
If $c$ is free, the pair $(s, c) = (a, c)$ with $c \to \infty$ sends
$\operatorname{KL} = \tfrac{1}{2\tau^2}(s-a)^2 = 0$ and
$\mathrm{CE} = \ln(1 + (K-1)e^{-ca}) \to 0$ simultaneously (for $a > 0$): the infimum
of $L_\beta$ over the aligned family is $0$ for \emph{every} $\beta$, no minimizer is
attained, and the one-to-one sweep of Proposition~\ref{prop:sweep} has no global
counterpart. (At $a = 0$ the margin growth alone cannot drive CE to zero.) This is
the FGIB analogue of the dead knob --- the same one-corner attractor, now arrived at by
growing the margin rather than by the linearity of the deterministic IB curve.
Three consequences.
\emph{(a) The present implementation is exposed.} The classifier is
\texttt{nn.Linear} trained by Adam with no weight decay, so $c$ is unconstrained. On
separable data the CE gradient drives $\|w_y\|$ upward; it does not diverge in finite
time (the gradient vanishes as $p_y \to 1$, giving logistic growth), but at any finite
scale $c$ the equilibrium $s^*(\beta; c)$ sits on the curve of
Proposition~\ref{prop:sweep} \emph{for that $c$}, and $s^*$ drifts toward $a$ as $c$
grows.
\emph{(b) Testable prediction.} The $\beta$-dependence of the equilibrium should decay
with training length; since the experiments use early stopping with $\texttt{patience} =
10$, part of any observed $\beta$-effect is plausibly an early-stopping effect. This can
be tested by training well past saturation and comparing the converged KL across
$\beta$.
\emph{(c) Architectural fix.} The device that satisfies the fixed-$c$ assumption
\emph{by construction} is a cosine classifier: normalized features, normalized rows,
and a fixed temperature (scale). Weight normalization with a learnable gain does not
guarantee a bounded scale, and weight decay only replaces the free-scale problem with a
joint optimization problem over $(s, c)$ that must be re-derived --- it mitigates the
divergence but does not make Proposition~\ref{prop:sweep} exact. Under a fixed-norm
classifier with fixed temperature the proposition holds as stated, at the price of one
architecture change.
\end{remark}

\textbf{Where the convexity comes from (mechanism).}
Compare the three objectives on their natural solution families, as functions of the
single compression scalar:
\begin{itemize}
\item \emph{dIB on hard clusterings} (Caveats Appendix~B): $L^{\mathrm{dIB}}_\beta
= (1-\beta)H(T)$ --- compression $H(T)$ and prediction $I(Y;T) = H(T)$ are the
\emph{same} scalar, so the objective is monotone in it and the optimum sits at the
endpoint for every $\beta < 1$. Squared-dIB restores interior optima by hand,
$H(T) - \beta H(T)^2$ with unique maximizer $H(T) = 1/(2\beta)$.
\item \emph{Standard IB} (deterministic): the curve is linear ($F'(r) \equiv 1$), every
point shares one supporting slope, and the corner $\langle H(Y), H(Y)\rangle$ is
selected for all $\beta \in (0,1)$. Squared-IB convexifies via $F'(r)/(2r)$ strictly
decreasing.
\item \emph{FGIB}: the compression term is $\operatorname{KL}(q_y \| r_y) =
\tfrac{1}{2\tau^2}\|\mu_y - a v_y\|^2 + {}$(variance terms), equal to
$\tfrac{1}{2\tau^2}(s-a)^2$
at the variance optimum --- \emph{quadratic} in the alignment scalar $s$, hence strictly
convex, because the KL is measured against a \emph{fixed} reference $a v_y$:
$\operatorname{KL}(\mathcal{N}(\mu,\cdot)\|\mathcal{N}(\mu_p,\cdot))
= \tfrac{1}{2\tau^2}\|\mu - \mu_p\|^2 + {}$(variance terms) is quadratic in the
displacement
$\mu - \mu_p$ whenever $\mu_p$ does not move. Entropy-based compression ($H(T)$,
$I(X;T)$) is self-referential --- it depends only on the representation's own
distribution and enters linearly in the clustering probabilities; the fixed anchor
supplies an absolute reference frame and turns compression into a squared distance. The
convexification that squared-IB installs by squaring the compression term is thus built
into FGIB's fixed-anchor KL. Together with the strictly convex prediction term
(Lemma~\ref{lem:lagrangian}), the Lagrangian is strictly convex with a unique interior
minimizer per $\beta$ \emph{at fixed classifier scale} (Proposition~\ref{prop:sweep}),
and the resulting curve is strictly concave (Proposition~\ref{prop:curve}). With the
scale free, the same quadratic gap is driven to zero by margin growth alone
(Remark~\ref{rem:c-scale}) --- the convexification is necessary but not sufficient
without the scale condition.
\end{itemize}

\begin{remark}[The compressed endpoint is non-collapsed but still a label code --- Caveat 2 remains]\label{rem:endpoint}
VIB's fixed marginal $m(z) = \mathcal{N}(0, I)$ places the $\operatorname{KL}=0$ point
at $\mu = 0$ for \emph{all} classes: the fully compressed solution is the trivial
collapse with $\mathrm{CE} = \ln K$ (chance level), and squaring (SVIB) leaves this
corner in place since $\operatorname{KL}^2$ is stationary at $\operatorname{KL}=0$.
FGIB's per-class orthonormal anchors make the $\operatorname{KL}=0$ point a
\emph{class-separated} configuration: $\mu_y = a v_y$ with
$\langle v_y, v_{y'}\rangle = \delta_{yy'}$, so at the corner
$\mathrm{CE}(a) = \ln(1 + (K-1)e^{-ca})$, which is bounded and tends to $0$ as the
anchor scale $a$ grows, and $a = 0$ recovers the VIB corner. But class-separation does
not buy non-triviality in the sense of the Caveats paper: their Caveat~2 identifies as
uninteresting the variables on the manifold $T_\alpha$ that mix a constant map with an
exact copy of $Y$, compressing by ``forgetting the input with some probability'' rather
than by merging perceptually similar classes --- and FGIB's compressed corner
$\mu_y = a v_y$ is a deterministic function of $Y$ alone, i.e.\ a label code of exactly
that kind. The orthonormal frame is in fact \emph{opposed} to the kind of compression
Caveat~2 holds up as interesting: it places every pair of classes at the same distance
and can never express ``merge collie with retriever, keep both far from teapot''. What
the anchors do is choose \emph{one specific} label-code geometry --- class means on a
scaled simplex, a neural-collapse-like configuration --- instead of leaving it to the
optimizer, and refuse the chance-level corner. The compressed endpoint is thus designed,
not certified, and any claim that the method compresses \emph{input} semantics at that
endpoint is exactly the claim Caveat~2 rules out.
\end{remark}

\begin{remark}[Residual-information interpretation (CEB)]\label{rem:residual}
For any backward encoder, $\mathbb{E}[\operatorname{KL}(q(z|x)\|r(z|y))]
= I(X;Z|Y) + \mathbb{E}_y[\operatorname{KL}(q(z|y)\|r(z|y))]
\ge I(X;Z|Y)$ (Fischer 2020, Eqs.~9--11, 20; Lemma~\ref{lem:gap}). FGIB is therefore a
CEB bottleneck with a \emph{fixed} backward encoder: the regularizer upper-bounds the
residual information, and on the aligned family the residual vanishes identically
($q(z|x)$ depends on $y$ alone), so the regularizer equals the anchor-matching gap
$\tfrac{1}{2\tau^2}(s-a)^2$. The quadratic mechanism of Proposition~\ref{prop:sweep} is
inherited from this gap term, which exists only because the anchor is not allowed to
follow the encoder.
\end{remark}

\begin{remark}[Variance decoupling and training--evaluation consistency]
Since the classifier consumes the deterministic $h$ (and the KL does not see the
classifier), the posterior variance decouples: $\sigma^{2*} = \tau^2$ for all $\beta$, and
the trade-off is purely geometric (mean alignment). This contrasts with VIB, where the
classifier consumes the sampled $z$: the noise couples into the prediction term (by
Jensen, $\mathbb{E}_z[\mathrm{CE}(c(z),y)] \ge \mathrm{CE}(c(\mu),y)$), the rate
$R \ge I(X;Z)$ over-counts compression, and training-time sampling versus
evaluation-time mean creates a distribution mismatch (Fischer 2020, Appendix~A.1).
\end{remark}

\begin{remark}[Idealizations]\label{rem:caveats}
The class-symmetric ansatz idealizes the equilibrium geometry; it is exact in the
$\beta$-dominant regime and captures the direction of the KL force in general. For
regression, the continuous RFF anchors satisfy $\|\mu_p(y)\| \approx a$, so the KL again
contains the quadratic displacement term
$\tfrac{1}{2}\|\mu - \mu_p(y)\|^2 + {}$(variance terms) --- the same convexity mechanism
carries over, again subject to the fixed-scale condition of Remark~\ref{rem:c-scale}
(the regression classifier has no margin scale; its norm is likewise unconstrained in
the implementation). Training starts at $\operatorname{KL} \approx a^2/2$
(zero-initialized logvar head, $\sigma = 1$, $\mu \approx 0$), so the anchor scale sets
the initial regularization pressure; larger $a$ delays the approach to $s^*$. The
finite-scale caveat formerly in this remark is stated, with its fix, in
Remark~\ref{rem:c-scale}.
\end{remark}

% ---------------------------------------------------------------
\subsection{The certificate}
% ---------------------------------------------------------------

\begin{lemma}[Gap identity]\label{lem:gap}
For any conditional $r(z|y)$,
\begin{equation}
\mathbb{E}\big[\operatorname{KL}(q(z|x)\,\|\,r(z|y))\big]
= I(X;Z|Y) + \mathbb{E}_y\big[\operatorname{KL}(q(z|y)\,\|\,r(z|y))\big] .
\label{eq:gap}
\end{equation}
\end{lemma}

\begin{proof}
$\mathbb{E}[\operatorname{KL}(q(z|x)\|r(z|y))] = \mathbb{E}[\ln q(z|x) - \ln r(z|y)]$
and $I(X;Z|Y) = \mathbb{E}[\ln q(z|x) - \ln q(z|y)]$ (Fischer 2020, Eqs.~9--10), where
both expectations are over $p(x,y)q(z|x)$; subtracting,
$\mathbb{E}[\ln q(z|y) - \ln r(z|y)] = \mathbb{E}_y[\operatorname{KL}(q(z|y)\|r(z|y))]$
since the argument depends on $z, y$ only. \hfill\qed
\end{proof}

\begin{lemma}[Moment matching]\label{lem:mm}
Within the per-class Gaussian family with \emph{diagonal covariance}, the gap of
Lemma~\ref{lem:gap} is
minimized by the moment-matched Gaussian $b^*_y = \mathcal{N}(m_y, S_y)$ with
$m_y = \mathbb{E}\,q(z|y)$, $S_y = \operatorname{diag}\operatorname{Var} q(z|y)$.
\end{lemma}

\begin{proof}
$\operatorname{KL}(p\|q)$ as a function of $q$ is minimized by matching the first two
moments of $p$ (standard; for Gaussians the minimizer is unique). \hfill\qed
\end{proof}

\begin{proposition}[Variational ordering: CEB's bound is tighter in value]\label{prop:ordering}
For every fixed encoder $q(z|x)$,
$\mathrm{ReX}(\theta^*) \le \mathrm{ReX}_F$ where $\theta^*$ is CEB's optimal backward
encoder, and the excess is the \emph{anchor gap}
\begin{equation}
\Delta := \mathrm{ReX}_F - \mathrm{ReX}(\theta^*)
= \mathbb{E}_y\big[\operatorname{KL}(q(z|y)\,\|\,N(a v_y, \tau^2 I))\big]
- \mathbb{E}_y\big[\operatorname{KL}(q(z|y)\,\|\,N(m_y, S_y))\big] \ge 0 ,
\label{eq:anchor-gap}
\end{equation}
with equality iff the anchors equal the moment-matched marginals. Here $\theta^*$ ranges
over the \emph{ideal} CEB backward encoder family --- per-class Gaussians with free mean
and \emph{free diagonal covariance} (the family of Lemma~\ref{lem:mm}, matching the
diagonal implementation), i.e.\ any conditional $r(z|y)$ of that family the theory can
express; the code's
single affine \texttt{prior\_net} is only an approximation of this family, and the
ordering is a statement about the idealized CEB, not about the implemented
parameterization.
On the aligned family of the sweep analysis ($q(z|y) = \mathcal{N}(s v_y, \sigma^2 I)$,
so $I(X;Z|Y) = 0$ and $b^* = q(z|y)$ exactly),
\begin{equation}
\Delta = \tfrac{1}{2\tau^2}(s-a)^2 + \tfrac{z}{2}\big(\sigma^2/\tau^2 - 1 - \ln(\sigma^2/\tau^2)\big) ,
\label{eq:aligned-gap}
\end{equation}
i.e.\ $\mathrm{ReX}_F = \Delta$ and $\mathrm{ReX}(\theta^*) = 0$: FGIB's bound value
\emph{equals} the sweep-driving compression term of Proposition~\ref{prop:sweep}, while
CEB's bound collapses to the (zero) residual.
\end{proposition}

\begin{proof}
The anchor $N(a v_y, \tau^2 I)$ is a member of CEB's per-class Gaussian family, so the
minimum over the family is no larger than the value at the anchor: the ordering and
Eq.~\eqref{eq:anchor-gap} follow from Lemmas~\ref{lem:gap} and~\ref{lem:mm}.
On the aligned family the marginal is itself Gaussian, so $b^* = q(z|y)$ gives gap $0$,
$I(X;Z|Y) = 0$ (the posterior is class-constant), and
$\operatorname{KL}(\mathcal{N}(s v_y, \sigma^2 I)\|\mathcal{N}(a v_y, \tau^2 I))
= \tfrac{1}{2\tau^2}(s-a)^2 + \tfrac{z}{2}(\sigma^2/\tau^2-1-\ln(\sigma^2/\tau^2))$.
\hfill\qed
\end{proof}

\begin{remark}[Tightness--force trade-off]\label{rem:force}
Proposition~\ref{prop:ordering} is the price of the fixed anchor, and it is paid for a
reason: in CEB the gap is a function of the \emph{backward encoder alone} ---
$b_\theta \to q(z|y)$ eliminates it with the encoder held fixed, so CEB's bound can
become tight while the representation stays uncompressed: its tightness certifies
nothing about compression, and the gap is an unobservable optimization residual that
tracks the forward encoder. FGIB's gap can shrink only if the \emph{posterior} moves
toward the anchors (the anchor itself cannot move), so the force is persistent and, on
the aligned family, the bound value \emph{is} the anchor-matching gap
$\tfrac{1}{2\tau^2}(s-a)^2$, the sweep-driving compression term
(Eq.~\eqref{eq:aligned-gap}). The qualifier is necessary: ``the posterior'' includes
the side head $A$, which is trainable and unused at inference, so part of the movement
can be absorbed without the deployed representation $h$ moving
(Remark~\ref{rem:A-channel}). On the aligned family, where $A$ is fixed by the ansatz,
the equivalence is exact; off it, bound tightness is monotonically equivalent to
alignment of the posterior, not to compression of the deployed $h$.
\end{remark}

\begin{proposition}[The inverse linear transform: the object of the bound]\label{prop:inverse}
Let the side head $T(h) = Ah$ be \emph{invertible} ($z = d$, unconstrained), and let
the trained posterior be isotropic, $\Sigma(x) = \sigma_p^2 I$ with $\sigma_p^2 = \tau^2$
the anchor variance (the KL variance term is minimized at $\sigma^2 = \tau^2$, so this
holds at the variance optimum).
Then $\tilde z := A^{-1} z$ gives
$q'(\tilde z|x) = \mathcal{N}(h(x), \tau^2 (A^{\top}A)^{-1})$ --- the posterior is
centered at the \emph{inference representation} $h$ with side noise
$\tau^2 (A^{\top}A)^{-1}$ --- and the transformed anchors
$r'(\tilde z|y) = \mathcal{N}(a\, A^{-1} v_y, \tau^2 (A^{\top}A)^{-1})$ form an
orthonormal frame at scale $a$ in the induced metric $A^{\top}A$
($\langle A^{-1}v_i, A^{-1}v_j\rangle_{A^{\top}A} = \delta_{ij}$: the fixed frame
$\{v_y\}$ read in canonical coordinates). By KL invariance under invertible linear maps
(Lemma~\ref{lem:kl-invariance}),
\begin{equation}
\mathrm{ReX}_F
= I(X;\, h + \tau\,(A^{\top}A)^{-1/2}\varepsilon\,|Y)
+ \mathbb{E}_y\big[\operatorname{KL}\big(q'(\tilde z|y)\,\|\,r'(\tilde z|y)\big)\big],
\qquad \varepsilon \sim \mathcal{N}(0, I) .
\label{eq:object}
\end{equation}
The bound's object is the inference representation perturbed by \emph{known,
side-path-only} noise (its covariance is a model parameter), and the anchor geometry is
preserved exactly under the transform, in the induced metric. The perturbation,
however, is small only if the side head is well-conditioned: its magnitude is governed
by $A^{\top}A$, which is unconstrained
(Proposition~\ref{prop:conditioning}(iii)).
\end{proposition}

\begin{proof}
$A^{-1}A = I$ and $A^{-1} z = h + A^{-1}\varepsilon$ for $\varepsilon \sim
\mathcal{N}(0, \Sigma)$; $A^{-1}\Sigma A^{-\top} = \tau^2 A^{-1}A^{-\top}
= \tau^2 (A^{\top}A)^{-1}$ for $\Sigma = \tau^2 I$; the prior covariance maps to
$A^{-1} (\tau^2 I) A^{-\top} = \tau^2 (A^{\top}A)^{-1}$; and
$\langle A^{-1}v_i, A^{-1}v_j\rangle_{A^{\top}A}
= v_i^{\top} A^{-\top} A^{\top} A\, A^{-1} v_j = \delta_{ij}$.
Equation~\eqref{eq:object} is Lemma~\ref{lem:gap} applied to the transformed pair.
\hfill\qed
\end{proof}

\begin{proposition}[Orthogonal-factor conditioning: the reference is fixed and well-conditioned]\label{prop:conditioning}
FGIB's anchor frame $V = [v_1, \dots, v_K]$ is the orthonormal factor of a random matrix
(QR: $G = QR$, $V := Q$, $V^{\top}V = I_K$), and the anchor covariance is
$\tau^2 I$ with $\tau^2 > 0$.
Then, for every anchor scale $a \ge 0$ and throughout training:
\begin{enumerate}
\item[(i)] the \emph{reference} of the matching --- the \emph{unscaled} mean frame $V$
and the covariance $\tau^2 I$ --- is fixed (non-trainable buffers) and has condition
number $1$: the certificate is a distance to a compact, non-degenerate,
class-separated (for $a > 0$) target. (At $a = 0$ the mean matrix $a V$ is zero and
has no condition number; the condition-$1$ statement attaches to $V$ and $\tau^2 I$,
not to $a V$.)
\item[(ii)] in side-mean coordinates the certificate's displacement metric is a
scaled identity: the KL's displacement term is $\tfrac{1}{2\tau^2}\|\mu - a v_y\|^2$
with gradient $(\mu - a v_y)/\tau^2$;
\item[(iii)] on the shared-encoder side the metric is $A^{\top}A$, where $A$ is a
learnable unconstrained map, so its \emph{global} condition number is not controlled by
construction --- and the aligned family does not control it either. An earlier version
inferred condition number $1$ on the class-mean subspace from the identity
$h_i^{\top} A^{\top} A\, h_j = s^2\delta_{ij}$; that inference is invalid, since
$\{h_i\}$ is not a Euclidean orthonormal basis. Counterexample: with
$A = \operatorname{diag}(\varepsilon, 1)$, $h_1 = (1/\varepsilon, 0)$,
$h_2 = (0, 1)$ and $s = 1$, one has $A h_i = v_i$ yet
$\kappa(A^{\top}A) = \varepsilon^{-2}$ --- and the inverse-transform perturbation of
Proposition~\ref{prop:inverse} is correspondingly $\varepsilon^{-2}$-large on the first
direction. The counterexample is not adversarial: it is the direction the gap gradient
points in, since the KL is minimized by fitting the anchor target on $h$ in least
squares, which suppresses within-class directions
(Remark~\ref{rem:A-channel}). Any guarantee on this side therefore requires an explicit
constraint on $A$ (orthogonal parameterization, weight norm, or the like) or an
empirical measurement of $\kappa(A^{\top}A)$ over training.
\end{enumerate}
CEB's backward encoder $b_\theta(z|y) = \mathcal{N}(m_\theta(y), S_\theta(y))$ is an
unconstrained linear map of the label: \emph{both} sides of its matching are trainable,
its covariance diagonal can lie anywhere in the clamp $[e^{-10}, e^{10}]$ (condition
number up to $e^{20}$), its means can be rank-deficient (class-merged codes, or the
collapse solution $m_\theta \equiv 0$), and the variance race drives $S_\theta$ to the
clamp, degenerating the certificate into a mean-matching force of weight $1/s = e^{10}$
(Prop.~\ref{prop:variance}, subject to Remark~\ref{rem:p3-status}). What the fixed QR
reference guarantees for FGIB is therefore the \emph{reference} side only: a compact,
non-degenerate, class-separated target whose conditioning cannot degrade. On the
encoder side FGIB's bound is exposed to the same degeneration mechanism CEB's is, with
$A$ playing the role of the second trainable party.
\end{proposition}

\begin{proof}
$V^{\top}V = I_K$: all singular values of $V$ are $1$; the covariance $\tau^2 I$ has
all eigenvalues $\tau^2$; both are frozen buffers, so (i) holds. The diagonal closed
form of
$\operatorname{KL}(\mathcal{N}(\mu, \sigma^2 I)\,\|\,\mathcal{N}(a v_y, \tau^2 I))$
has displacement term $\tfrac{1}{2\tau^2}\|\mu - a v_y\|^2$, giving (ii). For (iii):
the identity
$h_i^{\top} A^{\top} A\, h_j = \mu_i^{\top}\mu_j = s^2 v_i^{\top} v_j = s^2\delta_{ij}$
holds on the aligned family and only says that $\{h_i\}$ is $A^{\top}A$-orthogonal by
construction; it implies nothing about the spectrum of $A^{\top}A$, and the displayed
example verifies the failure. The KL gradient on $A$ is
$(\mu - a v_y)\,h^{\top}/\tau^2$, an alignment force toward the orthonormal frame that
vanishes only at perfect alignment \emph{on non-degenerate samples} (it also vanishes
pointwise at $h = 0$, or by cancellation across a batch) --- but nothing in that force
resists a vanishing within-class component of $A$. For CEB, the displacement term is
$\tfrac12(m - g)^{\top} S_\theta^{-1}(m - g)$: the metric $S_\theta^{-1}$ has condition
number up to $e^{20}$ under the logvar clamp, and the gradient weight $1/s$ at the
collapse equilibrium is Prop.~\ref{prop:variance}. \hfill\qed
\end{proof}

\begin{remark}[Two senses of tightness]\label{rem:senses}
In the pure variational \emph{value} sense (same encoder, same Gaussian family), CEB's
bound is tighter: $\mathrm{ReX}(\theta^*) \le \mathrm{ReX}_F$
(Prop.~\ref{prop:ordering}), the excess being the anchor gap --- the price of the fixed
anchor. The \emph{object} sense does not reverse this ordering: both bounds are valid
upper bounds on their own noisy codes, and neither bounds the residual of a deployed
deterministic continuous representation, which is infinite under the standing
non-atomic assumption (Prop.~\ref{prop:object}). The structural difference is where the certified-versus-
deployed gap lives --- in CEB it is driven by the prediction term inside the inference
path; in FGIB it is a side-path quantity whose size is governed by the conditioning of
$A$ (Prop.~\ref{prop:conditioning}(iii)) --- not which bound is ``tighter'' or
``valid''. The learnability of the side head $A$ bears on that difference: it is a
second trainable party in the matching, unused at inference
(Remark~\ref{rem:A-channel}), so it can absorb part of the anchor gap without moving
the deployed representation. On the aligned family, where $A$ is fixed by the ansatz,
its optimization drives the gap to the equilibrium value $\tfrac12(s^*-a)^2$ at fixed
classifier scale (Prop.~\ref{prop:sweep}); off the ansatz this is an open question
(Section~\ref{sec:limits}).
\end{remark}

\begin{remark}[Unconstrained side head (the implementation)]\label{rem:orth}
Proposition~\ref{prop:inverse} holds for any invertible linear head: the KL invariance
is exact, and the object claim needs only that the perturbation covariance
$\tau^2(A^{\top}A)^{-1}$ be \emph{known} ($A$ is a model parameter) --- but the
object is informative about $h$ only insofar as that covariance is moderate, which the
construction does not guarantee (Prop.~\ref{prop:conditioning}(iii)). The diagonal
variational family is not closed under $A^{-1}$ (Remark~\ref{rem:closure}), so within
the diagonal parameterization of the codebase the equivalence is exact at the
distribution level and approximate as a variational identity; a full-covariance side
head would make it exact within the family. No orthogonality constraint on $A$ is
required --- the orthogonality of the method lives in the fixed anchor frame and in the
equilibrium it induces (the aligned family).
\end{remark}

\begin{remark}[Anchor-frame reading]\label{rem:frame}
With the orthonormal anchor matrix $V$ (columns $v_y$, $K \le z$), left multiplication
by $V^{\top}$ expresses the anchor-space components of $z$ in class coordinates: the
anchors become $a\,e_y$ and the gap of Lemma~\ref{lem:gap} decomposes into $K$
per-class terms plus, when $K < z$, an orthogonal-complement term that the reading
ignores ($V^{\top}$ is a projection, not a coordinate change). The orthogonal frame is
the geometric content of the prior, and it is preserved by the invertible linear
pull-back $A^{-1}$ of Proposition~\ref{prop:inverse} --- not by $A^{\top}$, which is
the Jacobian factor that appears in gradients and equals $A^{-1}$ only when $A$ is
orthogonal. An earlier version conflated the two.
\end{remark}

\begin{remark}[Positioning against VIB and NIB]\label{rem:position}
VIB's rate bounds the \emph{full} mutual information $I(X;Z)$ with a marginal prior;
CEB removes the retained label information ($\mathrm{ReX} = R - I(Y;Z)$ at the optima),
and FGIB inherits the residual-level object while replacing the learnable backward
encoder with a fixed geometric prior (Prop.~\ref{prop:ordering}).
NIB's non-parametric bound $\widehat I \ge I(X;M)$ is $y$-unconditioned and tight only
when the components share a covariance and form well-separated clusters (Kolchinsky et
al.\ 2017, Eq.~10 and the discussion following Eq.~11); when the noise is small relative
to the cluster spread it saturates at $\ln N$ and its gradient vanishes (their
Section~IVA) --- loose exactly in the regime where the parametric FGIB bound is
well-behaved. The two bounds are complementary: NIB certifies the full compression
non-parametrically when the geometry is favorable; FGIB bounds the residual of a
\emph{noisy copy} of the deployed representation parametrically, with a geometry of its
own choosing --- subject to the conditioning caveat of
Proposition~\ref{prop:conditioning}(iii).
\end{remark}

\begin{remark}[Degenerate anchors]\label{rem:a0}
At $a = 0$ the anchors collapse to $N(0, \tau^2 I)$ for all classes and the gap loses its
class
structure: $\Delta = \mathbb{E}_y[\operatorname{KL}(q(z|y)\|N(0,\tau^2 I))]$ --- a
VIB-style
marginal gap; the structural separation of Proposition~\ref{prop:object}(b) still
holds, but the compressed endpoint becomes the trivial collapse
(Remark~\ref{rem:endpoint}).
\end{remark}

% ---------------------------------------------------------------
\subsection{What is proved and what is open}\label{sec:limits}
% ---------------------------------------------------------------

\begin{remark}[Scope of the results]
\textbf{Proved.} (1) The dead-knob proposition for CEB and the MNI reading of it
(Prop.~\ref{prop:dead}, Remark~\ref{rem:mni}); (2) the gap identity, Gaussian moment
matching, and the value ordering of CEB's bound over FGIB's at a fixed encoder
(Lemmas~\ref{lem:gap}, \ref{lem:mm}; Prop.~\ref{prop:ordering}); (3) the blind
certificate on the $T_\alpha$ manifold and the $1{:}\beta$ exchange rate
(Prop.~\ref{prop:flat}); (4) KL invariance under invertible linear maps and the
pointwise equivalence of the side path to a CEB bottleneck at a fixed side head
(Lemmas~\ref{lem:kl-invariance}, \ref{lem:mi-invariance}; Prop.~\ref{prop:fgib-ceb});
(5) the variance decoupling $\sigma^{2*} = \tau^2$ and the non-collapsed compressed endpoint
(Lemma~\ref{lem:lagrangian}, with $\sigma^{2*} = \tau^2$;
Remark~\ref{rem:endpoint}); (6) the one-dimensional sweep
and its strict concavity \emph{at fixed classifier scale}, as a KL--CE surrogate curve
on the aligned family (Props.~\ref{prop:sweep}, \ref{prop:curve};
Remark~\ref{rem:surrogate} --- not a recovered IB curve); (7) the well-conditioning of the fixed
reference side (Prop.~\ref{prop:conditioning}(i)--(ii)).

\textbf{Open or conditional.} (1) The side head $A$ is a trainable, inference-unused
party in the matching (Remark~\ref{rem:A-channel}); combined with the
counterexample of Prop.~\ref{prop:conditioning}(iii), the perturbation in
Prop.~\ref{prop:inverse} can be arbitrarily large, so the bound's object may be far from
the deployed representation exactly where the residual lives. A constraint on $A$ or a
measurement of $\kappa(A^{\top}A)$ over training would close or quantify this channel.
(2) The global sweep requires a bounded classifier scale
(Remark~\ref{rem:c-scale}); the current implementation leaves the scale free and the
proposition conditional. A fixed-temperature cosine classifier would satisfy the
fixed-$c$ assumption by construction; weight decay only mitigates the scale divergence
and replaces the problem with a new joint optimization over $(s,c)$ that must be
re-derived. (3) Neither CEB's nor FGIB's certificate upper-bounds the deployed
deterministic residual (Prop.~\ref{prop:object}); claims of inference-representation
compression are hypotheses, not conclusions. (4) Caveat~2 is not resolved
(Remark~\ref{rem:endpoint}): the compressed endpoint is a chosen label-code geometry,
not certified semantic compression. (5) The two-stage variance dynamics of
Prop.~\ref{prop:variance}(ii)--(iii) await the measurement of
Remark~\ref{rem:p3-status}. (6) The regression RFF anchor analysis inherits the same
fixed-scale and conditioning caveats and is presently a transfer argument, not a proof.
\end{remark}
